% Generated mechanically from the frozen crystalline.tex target.
% Do not edit this generated file; edit the bound locale source instead.

% 好，从这里开始。
%

\setcounter{chapter}{59}
\chapter{晶体上同调}



\phantomsection
\label{crystalline-section-phantom}




\section{引言}
\label{crystalline-section-introduction}

\noindent
本章以 Johan de Jong 于 2012 年在哥伦比亚大学所作的一系列讲座为基础。
本章旨在快速介绍晶体上同调。参考文献为专著 \cite{Berthelot}。

\medskip\noindent
关于除幂环的较初等纯代数讨论已经移至一个预备章节，因为它在讨论交换代数中的
Tate 分解时也很有用。参见“除幂代数”第
\ref{dpa-section-introduction} 节。













\section{除幂包络}
\label{crystalline-section-divided-power-envelope}

\noindent
以下引理中的构造称为除幂包络。它将在后文发挥重要作用。

\begin{lemma}
\label{crystalline-lemma-divided-power-envelope}
设 $(A, I, \gamma)$ 为除幂环，$A \to B$ 为环同态。设
$J \subset B$ 为满足 $IB \subset J$ 的理想。则存在除幂环同态
$$
(A, I, \gamma) \longrightarrow (D, \bar J, \bar \gamma)
$$
使得
$$
\Hom_{(A, I, \gamma)}((D, \bar J, \bar \gamma), (C, K, \delta)) =
\Hom_{(A, I)}((B, J), (C, K))
$$
并且该等式关于 $(A, I, \gamma)$ 上的除幂代数 $(C, K, \delta)$
具有函子性。这里，左端表示 $(A, I, \gamma)$ 上除幂环的态射，
右端表示 $(A, I)$ 上的（环，理想）对的态射。
\end{lemma}

\begin{proof}
以 $\mathcal{C}$ 表示由除幂环 $(C, K, \delta)$ 组成的范畴。考虑如下定义的函子
$F : \mathcal{C} \longrightarrow \textit{Sets}$：
$$
F(C, K, \delta) =
\left\{
(\varphi, \psi)
\middle|
\begin{matrix}
\varphi : (A, I, \gamma) \to (C, K, \delta)
\text{ 为除幂环同态} \\
\psi : (B, J) \to (C, K)\text{ 为满足 }
\psi(J) \subset K\text{ 的 }A\text{-代数同态}
\end{matrix}
\right\}
$$
我们将证明“除幂代数”引理 \ref{dpa-lemma-a-version-of-brown}
适用于此函子，从而得到本引理。设
$(\varphi, \psi) \in F(C, K, \delta)$。令 $C' \subset C$ 为由
$\varphi(A)$、$\psi(B)$ 以及所有 $f \in J$ 所对应的
$\delta_n(\psi(f))$ 生成的子环。令 $K' \subset K \cap C'$ 为
$C'$ 中由 $\varphi(I)$ 以及 $f \in J$ 所对应的
$\delta_n(\psi(f))$ 生成的理想。于是
$(C', K', \delta|_{K'})$ 是除幂环，并且 $C'$ 的基数由基数
$\kappa = |A| \otimes |B|^{\aleph_0}$ 控制。此外，$\varphi$ 分解为
$A \to C' \to C$，而 $\psi$ 分解为 $B \to C' \to C$。
这就证明“除幂代数”引理 \ref{dpa-lemma-a-version-of-brown}
的假设 (1) 成立。假设 (2) 是显然的，因为除幂环范畴中的极限与遗忘函子
$(C, K, \delta) \mapsto (C, K)$ 可交换；参见“除幂代数”引理
\ref{dpa-lemma-limits} 及其证明。
\end{proof}

\begin{definition}
\label{crystalline-definition-divided-power-envelope}
设 $(A, I, \gamma)$ 为除幂环，$A \to B$ 为环同态，并设
$J \subset B$ 为满足 $IB \subset J$ 的理想。引理
\ref{crystalline-lemma-divided-power-envelope} 中构造的除幂代数
$(D, \bar J, \bar\gamma)$ 称为{\it $B$ 中 $J$ 相对于
$(A, I, \gamma)$ 的除幂包络}，记作 $D_B(J)$ 或
$D_{B, \gamma}(J)$。
\end{definition}

\noindent
设 $(A, I, \gamma) \to (C, K, \delta)$ 为除幂环同态。
$D_{B, \gamma}(J) = (D, \bar J, \bar \gamma)$ 的泛性质为
$$
\begin{matrix}
\text{把 }J\text{ 映入 }K\text{ 的环同态 }B \to C
\end{matrix}
\longleftrightarrow
\begin{matrix}
\text{除幂同态} \\
(D, \bar J, \bar \gamma) \to (C, K, \delta)
\end{matrix}
$$
该对应由预复合与 $\text{id}_D$ 对应的映射 $B \to D$ 给出。
下面列出由泛性质直接推出的 $(D, \bar J, \bar \gamma)$ 的若干性质。
存在 $A$-代数同态
\begin{equation}
\label{crystalline-equation-divided-power-envelope}
B \longrightarrow D \longrightarrow B/J
\end{equation}
第一个箭头把 $J$ 映入 $\bar J$，而 $\bar J$ 是第二个箭头的核。
当 $n > 0$ 且 $x$ 属于 $J \to D$ 的像时，元素
$\bar\gamma_n(x)$ 作为 $D$ 中的理想生成 $\bar J$，并作为
$B$-代数生成 $D$。

\begin{lemma}
\label{crystalline-lemma-divided-power-envelop-quotient}
设 $(A, I, \gamma)$ 为除幂环。设
$\varphi : B' \to B$ 为核等于 $K$ 的 $A$-代数满射，
$IB \subset J \subset B$ 为理想，并令 $J' \subset B'$ 为 $J$ 的逆像。
记 $D_{B', \gamma}(J') = (D', \bar J', \bar\gamma)$。则
$D_{B, \gamma}(J) = (D'/K', \bar J'/K', \bar\gamma)$，其中
$K'$ 是由 $n \geq 1$、$k \in K$ 时的元素 $\bar\gamma_n(k)$
生成的理想。
\end{lemma}

\begin{proof}
记 $D_{B, \gamma}(J) = (D, \bar J, \bar \gamma)$。$D'$ 的泛性质给出
除幂代数同态 $D' \to D$。由于 $B' \to B$ 和 $J' \to J$ 均为满射，
可知 $D' \to D$ 也是满射（见上文的说明）。显然，当 $n \geq 1$、
$k \in K$ 时，$\bar\gamma_n(k)$ 属于其核；因而得到同态
$D'/K' \to D$。反过来，$\bar J'/K' \subset D'/K'$ 上存在除幂结构；
参见“除幂代数”引理 \ref{dpa-lemma-kernel}。于是 $D$ 的泛性质给出逆映射
$D \to D'/K'$，从而得证。
\end{proof}

\noindent
在定义 \ref{crystalline-definition-divided-power-envelope} 的情形中，可以选取满射
$P \to B$，其中 $P$ 是 $A$ 上的多项式代数，并令 $J' \subset P$
为 $J$ 的逆像。前一引理用 $D_{P, \gamma}(J')$ 描述了
$D_{B, \gamma}(J)$。注意，由“除幂代数”引理
\ref{dpa-lemma-gamma-extends}，$\gamma$ 可延拓为 $IP$ 上的除幂结构
$\gamma'$。因此 $D_{P, \gamma}(J') = D_{P, \gamma'}(J')$ 是下一引理
所述除幂包络特殊情形的一个例子。

\begin{lemma}
\label{crystalline-lemma-describe-divided-power-envelope}
设 $(B, I, \gamma)$ 为除幂代数，$I \subset J \subset B$ 为理想。
设 $(D, \bar J, \bar \gamma)$ 为 $J$ 相对于 $\gamma$ 的除幂包络。
选取元素 $f_t \in J$（$t \in T$），使得 $J = I + (f_t)$。则存在满射
$$
\Psi :
(B\langle x_t \rangle, IB\langle x_t \rangle + B\langle x_t \rangle_+, \delta)
\longrightarrow
(D, \bar J, \bar \gamma)
$$
它是除幂环同态，并把 $x_t$ 映为 $f_t$ 在 $D$ 中的像。$\Psi$ 的核由
元素 $x_t-f_t$ 以及所有如下元素生成：
$$
\delta_n\left(\sum r_t x_t - r_0\right)
$$
其中 $r_t \in B$、$r_0 \in I$，并且在 $B$ 中
$\sum r_t f_t = r_0$。
\end{lemma}

\begin{proof}
在本引理的陈述中，把 $B\langle x_t \rangle$ 看成以如下理想为指定理想的除幂环：
$J' = IB\langle x_t \rangle + B\langle x_t \rangle_{+}$；参见“除幂代数”注
\ref{dpa-remark-divided-power-polynomial-algebra}。
$\Psi$ 的存在性来自除幂多项式环的泛性质。其像是 $D$ 的除幂子环，
故由 $D$ 的泛性质等于 $D$；这就说明 $\Psi$ 是满射。显然，
$x_t-f_t$ 属于其核。令
$$
\mathcal{R} = \{(r_0, r_t) \in I \oplus \bigoplus\nolimits_{t \in T} B
\mid \sum r_t f_t = r_0 \text{ 于 }B\text{ 中}\}
$$
若 $(r_0,r_t)\in\mathcal R$，则显然
$\sum r_t x_t-r_0$ 属于该核。由于 $\Psi$ 是除幂环同态且
$\sum r_tx_t-r_0\in J'$，故
$\delta_n(\sum r_t x_t-r_0)$ 也属于该核。令
$K\subset B\langle x_t\rangle$ 为由 $x_t-f_t$ 以及
$(r_0,r_t)\in\mathcal R$ 时的元素
$\delta_n(\sum r_t x_t-r_0)$ 生成的理想。
为了证明 $K=\Ker(\Psi)$，只需证明 $\delta$ 可延拓至
$B\langle x_t\rangle/K$。事实上，若能如此，$D$ 的泛性质便给出
$\Psi$ 的逆映射 $D\to B\langle x_t\rangle/K$。因此，需要证明
$K\cap J'$ 在 $\delta_n$ 下保持不变；参见“除幂代数”引理
\ref{dpa-lemma-kernel}。令 $K'\subset B\langle x_t\rangle$
为由下列元素生成的理想：
\begin{enumerate}
\item 当 $m>0$ 且 $(r_0,r_t)\in\mathcal R$ 时的
$\delta_m(\sum r_t x_t-r_0)$；
\item 当 $m>0$ 且 $t',t\in I$ 时的 $x_{t'}^{[m]}(x_t-f_t)$。
\end{enumerate}
我们断言 $K'=K\cap J'$。根据“除幂代数”引理
\ref{dpa-lemma-kernel} (2)(c) 的判据，以及对上列元素的
$\delta_n$ 的计算（留给读者），该断言说明当 $n>0$ 时
$K\cap J'$ 在 $\delta_n$ 下保持不变。为证明该断言，先注意到
$K'\subset K\cap J'$。反过来，若 $h\in K\cap J'$，则模 $K'$ 可写成
$$
h = \sum r_t (x_t - f_t)
$$
其中 $r_t\in B$。由于 $h\in K\cap J'\subset J'$，可知
$r_0=\sum r_t f_t\in I$，从而 $(r_0,r_t)\in\mathcal R$，并有
$$
h = \sum r_t x_t - r_0
$$
属于 $K'$，正合所需。
\end{proof}

\begin{lemma}
\label{crystalline-lemma-divided-power-envelope-add-variables}
设 $(A,I,\gamma)$ 为除幂环，$B$ 为 $A$-代数，
$IB\subset J\subset B$ 为理想，并设 $x_i$ 为一族变量。则
$$
D_{B[x_i], \gamma}(JB[x_i] + (x_i)) = D_{B, \gamma}(J) \langle x_i \rangle
$$
\end{lemma}

\begin{proof}
一种证明方法是由引理 \ref{crystalline-lemma-describe-divided-power-envelope} 推出，
因为 $B[x_i]$ 中变量 $x_i$ 之间的任何关系都是平凡的。另一方面，
本引理也可由除幂多项式代数的泛性质与除幂包络的泛性质推出。
\end{proof}

\noindent
若 $B\to B'$ 在 $V(IB')\subset\Spec(B')$ 的所有素点处均平坦，
则下一引理的条件 (1) 与 (2) 成立；而这两个条件也与该平坦性条件
密切相关，参见“代数”引理 \ref{algebra-lemma-what-does-it-mean}。
特别地，该引理说明取除幂包络与局部化可交换。

\begin{lemma}
\label{crystalline-lemma-flat-base-change-divided-power-envelope}
设 $(A,I,\gamma)$ 为除幂环，$B\to B'$ 为 $A$-代数同态。假设
\begin{enumerate}
\item $B/IB \to B'/IB'$ 是平坦的；
\item $\text{Tor}_1^B(B', B/IB) = 0$.
\end{enumerate}
则对任意理想 $IB\subset J\subset B$，典范映射
$$
D_B(J) \otimes_B B' \longrightarrow D_{B'}(JB')
$$
是同构。
\end{lemma}

\begin{proof}
令 $D=D_B(J)$，并以 $\bar J\subset D$ 表示其除幂理想，
相应的除幂结构记作 $\bar\gamma$。$D$ 的泛性质给出 $B$-代数同态
$D\to D_{B'}(JB')$，从而给出引理中的映射。只需证明
$\bar\gamma$ 可延拓至 $D\otimes_B B'$；因为一旦如此，
$D_{B'}(JB')$ 的泛性质便给出引理中映射的逆映射
$D_{B'}(JB')\to D\otimes_B B'$。

\medskip\noindent
选取满射 $P\to B'$，其中 $P$ 是 $B$ 上的多项式代数。特别地，
$B\to P$ 平坦，故由“代数”引理 \ref{algebra-lemma-flat-base-change}，
$D\to D\otimes_B P$ 也平坦。于是由“除幂代数”引理
\ref{dpa-lemma-gamma-extends}，$\bar\gamma$ 可延拓至
$D\otimes_B P$；仍以 $\bar\gamma$ 表示该延拓。令
$\mathfrak a=\Ker(P\to B')$，于是有短正合列
$$
0 \to \mathfrak a \to P \to B' \to 0
$$
因此，$\text{Tor}_1^B(B',B/IB)=0$ 蕴含
$\mathfrak a\cap IP=I\mathfrak a$。现有如下交换图：
$$
\xymatrix{
B/J \otimes_B \mathfrak a \ar[r]_\beta &
B/J \otimes_B P \ar[r] &
B/J \otimes_B B' \\
D \otimes_B \mathfrak a \ar[r]^\alpha \ar[u] &
D \otimes_B P \ar[r] \ar[u] &
D \otimes_B B' \ar[u] \\
\bar J \otimes_B \mathfrak a \ar[r] \ar[u] &
\bar J \otimes_B P \ar[r] \ar[u] &
\bar J \otimes_B B' \ar[u]
}
$$
即使在上方和右方补上 $0$，该图仍然正合。我们须证明理想
$\bar J \otimes_B P$ 上的除幂保持理想
$\Im(\alpha) \cap \bar J \otimes_B P$；参见“除幂代数”引理
\ref{dpa-lemma-kernel}。考虑正合列
$$
0 \to \mathfrak a/I\mathfrak a \to P/IP \to B'/IB' \to 0
$$
（这里使用了上面得到的 $\mathfrak a\cap IP=I\mathfrak a$）。
由于 $B'/IB'$ 在 $B/IB$ 上平坦，对该序列施加
$B/J\otimes_{B/IB}-$ 后仍为正合；参见“代数”引理
\ref{algebra-lemma-flat-tor-zero}。因此
$$
\Ker(B/J \otimes_{B/IB} \mathfrak a/I\mathfrak a \to
B/J \otimes_{B/IB} P/IP) =
\Ker(\mathfrak a/J\mathfrak a \to P/JP)
$$
为零。因此 $\beta$ 是单射，从而
$\Im(\alpha) \cap \bar J \otimes_B P$ 是
$\bar J\otimes\mathfrak a$ 的像。若 $f\in\bar J$ 且
$a\in\mathfrak a$，则
$\bar\gamma_n(f \otimes a) = \bar\gamma_n(f) \otimes a^n$
故结论显然成立。
\end{proof}

\noindent
下一引理是 \cite[命题 2.1.7]{dJ-crystalline} 的一个特殊情形，
而后者又推广了 \cite[命题 2.8.2]{Berthelot}。

\begin{lemma}
\label{crystalline-lemma-flat-extension-divided-power-envelope}
设 $(B, I, \gamma) \to (B', I', \gamma')$ 为除幂环同态，
$I \subset J \subset B$ 与 $I' \subset J' \subset B'$ 为理想。假设
\begin{enumerate}
\item $B/I \to B'/I'$ 是平坦的；
\item $J' = JB' + I'$.
\end{enumerate}
则典范映射
$$
D_{B, \gamma}(J) \otimes_B B' \longrightarrow D_{B', \gamma'}(J')
$$
是同构。
\end{lemma}

\begin{proof}
令 $D = D_{B, \gamma}(J)$。选取生成 $J/I$ 的元素 $f_t \in J$。
如引理 \ref{crystalline-lemma-describe-divided-power-envelope} 的证明中那样，令
$\mathcal{R} = \{(r_0, r_t) \in I \oplus \bigoplus\nolimits_{t \in T} B
\mid \sum r_t f_t = r_0 \text{ 于 }B\text{ 中}\}$。该引理说明
$$
D = B\langle x_t \rangle/ K
$$
其中 $K$ 由元素 $x_t - f_t$ 以及 $(r_0, r_t) \in \mathcal{R}$ 时的
$\delta_n(\sum r_t x_t - r_0)$ 生成。因此
\begin{equation}
\label{crystalline-equation-base-change}
D \otimes_B B' = B'\langle x_t \rangle/K'
\end{equation}
其中 $K'$ 由上述 $K$ 的生成元在 $B'\langle x_t \rangle$ 中的像生成。
令 $f'_t \in B'$ 为 $f_t$ 的像。由假设 (1)，元素 $f'_t \in J'$
生成 $J'/I'$，并且 $x_t - f'_t \in K'$。令
$$
\mathcal{R}' =
\{(r'_0, r'_t) \in I' \oplus \bigoplus\nolimits_{t \in T} B'
\mid \sum r'_t f'_t = r'_0 \text{ 于 }B'\text{ 中}\}
$$
为完成证明，须对 $(r'_0, r'_t) \in \mathcal{R}'$ 证明
$\delta'_n(\sum r'_t x_t - r'_0) \in K'$。因为一旦如此，
$D \otimes_B B'$ 的表示式 (\ref{crystalline-equation-base-change}) 就与引理
\ref{crystalline-lemma-describe-divided-power-envelope} 由生成元 $f'_t$ 得到的
$D_{B', \gamma'}(J')$ 的表示完全相同。设
$(r'_0, r'_t) \in \mathcal{R}'$。则
$\sum r'_t f'_t = 0$ 在 $B'/I'$ 中成立。由假设 (1)，
$B/I \to B'/I'$ 平坦，故可应用平坦性的方程判据（“代数”引理
\ref{algebra-lemma-flat-eq}），得到某个 $m > 0$，以及
$r_{jt} \in B$ 和 $c_j \in B'$（$j = 1, \ldots, m$），使得
$$
r_{j0} = \sum\nolimits_t r_{jt} f_t  \in I \text{，其中 } j = 1, \ldots, m
$$
以及
$$
i'_t  = r'_t - \sum\nolimits_j c_j r_{jt} \in I' \text{，对所有 }t
$$
注意，这还蕴含
$r'_0 = \sum_t i'_t f_t + \sum_j c_j r_{j0}$。于是
\begin{align*}
\delta'_n(\sum\nolimits_t r'_t x_t  - r'_0)
& =
\delta'_n(
\sum\nolimits_t i'_t x_t +
\sum\nolimits_{t, j} c_j r_{jt} x_t -
\sum\nolimits_t i'_t f_t -
\sum\nolimits_j c_j r_{j0}) \\
& =
\delta'_n(
\sum\nolimits_t i'_t(x_t - f_t) +
\sum\nolimits_j c_j (\sum\nolimits_t r_{jt} x_t  - r_{j0}))
\end{align*}
由于 $\delta_n(a + b) = \sum_{m = 0, \ldots, n} \delta_m(a) \delta_{n - m}(b)$，
并且当 $m > 0$ 时，$\delta_m(\sum i'_t(x_t - f_t))$ 属于由
$x_t - f_t \in K'$ 生成的理想，只需证明
$\delta_n(\sum c_j (\sum r_{jt} x_t  - r_{j0}))$ 属于 $K'$。为此使用
$$
\delta_n(\sum\nolimits_j c_j (\sum\nolimits_t r_{jt} x_t  - r_{j0}))
=
\sum c_1^{n_1} \ldots c_m^{n_m}
\delta_{n_1}(\sum r_{1t} x_t  - r_{10}) \ldots
\delta_{n_m}(\sum r_{mt} x_t  - r_{m0})
$$
其中求和遍及所有满足 $n_1 + \ldots + n_m = n$ 的指标。这正是所需结论。
\end{proof}





\section{若干显式除幂加厚}
\label{crystalline-section-explicit-thickenings}

\noindent
本节的构造将帮助我们定义晶体位点上模晶体的联络。

\begin{lemma}
\label{crystalline-lemma-divided-power-first-order-thickening}
设 $(A, I, \gamma)$ 为除幂环，$M$ 为 $A$-模。令
$B = A \oplus M$ 为 $A$-代数，其中 $M$ 是平方为零的理想；并令
$J = I \oplus M$。对 $x \in I$ 与 $z \in M$，令
$$
\delta_n(x + z) = \gamma_n(x) + \gamma_{n - 1}(x)z
$$
则 $\delta$ 是除幂结构，并且 $A \to B$ 是从
$(A, I, \gamma)$ 到 $(B, J, \delta)$ 的除幂环同态。
\end{lemma}

\begin{proof}
须检验“除幂代数”定义 \ref{dpa-definition-divided-powers}
中的条件 (1)--(5)。这里直接证明；下一引理的证明给出了一种避免计算的方法。
条件 (1) 与 (3) 显然。条件 (2) 由下式推出：
\begin{align*}
\delta_n(x + z)\delta_m(x + z)
& =
(\gamma_n(x) + \gamma_{n - 1}(x)z)(\gamma_m(x) + \gamma_{m - 1}(x)z) \\
& = \gamma_n(x)\gamma_m(x) + \gamma_n(x)\gamma_{m - 1}(x)z +
\gamma_{n - 1}(x)\gamma_m(x)z \\
& =
\frac{(n + m)!}{n!m!} \gamma_{n + m}(x) +
\left(\frac{(n + m - 1)!}{n!(m - 1)!} +
\frac{(n + m - 1)!}{(n - 1)!m!}\right)
\gamma_{n + m - 1}(x) z \\
& =
\frac{(n + m)!}{n!m!}\delta_{n + m}(x + z)
\end{align*}
条件 (5) 由下式推出：
\begin{align*}
\delta_n(\delta_m(x + z))
& =
\delta_n(\gamma_m(x) + \gamma_{m - 1}(x)z) \\
& =
\gamma_n(\gamma_m(x)) + \gamma_{n - 1}(\gamma_m(x))\gamma_{m - 1}(x)z \\
& =
\frac{(nm)!}{n! (m!)^n} \gamma_{nm}(x) +
\frac{((n - 1)m)!}{(n - 1)! (m!)^{n - 1}}
\gamma_{(n - 1)m}(x) \gamma_{m - 1}(x) z \\
& = \frac{(nm)!}{n! (m!)^n}(\gamma_{nm}(x) + \gamma_{nm - 1}(x) z)
\end{align*}
这里使用了初等数论。为证明 (4)，须验证
$$
\delta_n(x + x' + z + z')
=
\gamma_n(x + x') + \gamma_{n - 1}(x + x')(z + z')
$$
等于
$$
\sum\nolimits_{i = 0}^n
(\gamma_i(x) + \gamma_{i - 1}(x)z)
(\gamma_{n - i}(x') + \gamma_{n - i - 1}(x')z')
$$
分别比较 $1$、$z$ 与 $z'$ 的系数，再使用 $\gamma$ 的条件 (4)，
便容易得到该等式。
\end{proof}

\begin{lemma}
\label{crystalline-lemma-divided-power-second-order-thickening}
设 $(A, I, \gamma)$ 为除幂环，$M,N$ 为 $A$-模，并设
$q : M \times M \to N$ 为 $A$-双线性映射。令
$B = A \oplus M \oplus N$ 为具有如下乘法的 $A$-代数：
$$
(x, z, w)\cdot (x', z', w') = (xx', xz' + x'z, xw' + x'w + q(z, z') + q(z', z))
$$
并令 $J = I \oplus M \oplus N$。对 $(x, z, w) \in J$，令
$$
\delta_n(x, z, w) = (\gamma_n(x), \gamma_{n - 1}(x)z,
\gamma_{n - 1}(x)w + \gamma_{n - 2}(x)q(z, z))
$$
则 $\delta$ 是除幂结构，并且 $A \to B$ 是从
$(A, I, \gamma)$ 到 $(B, J, \delta)$ 的除幂环同态。
\end{lemma}

\begin{proof}
先证明“除幂代数”定义 \ref{dpa-definition-divided-powers}
中的性质 (4)。在 $J$ 中选取 $(x,z,w)$ 与 $(x',z',w')$，并选取映射
$$
A_0 = \mathbf{Z}\langle s, s'\rangle \longrightarrow A,\quad
s \longmapsto x,
s' \longmapsto x'
$$
除幂多项式环的泛性质保证可以如此选取。令
$M_0 = A_0 \oplus A_0$ 以及
$N_0 = A_0 \oplus A_0 \oplus M_0 \otimes_{A_0} M_0$，并令
$q_0 : M_0 \times M_0 \to N_0$ 为自然映射。把 $M_0 \to M$
定义为将 $M_0$ 的两个基向量分别映至 $z$ 与 $z'$ 的 $A_0$-线性映射。
把 $N_0 \to N$ 定义为如下 $A_0$-线性映射：将 $N_0$ 的前两个
基向量分别映至 $w$ 与 $w'$，并在最后一个直和项上使用
$M_0 \otimes_{A_0} M_0 \to M \otimes_A M \xrightarrow{q} N$
于是，只需对情形 $(A_0,M_0,N_0,q_0)$ 证明恒等式 (4)。
其他恒等式也同理。这就把问题归约到 $A$ 为无
$\mathbf{Z}$-挠环且各模均无 $A$-挠的情形。在此情形中，只须证明
$$
n! \delta_n(x, z, w) = (x, z, w)^n
$$
在环 $A$ 中成立；参见“除幂代数”引理 \ref{dpa-lemma-silly}。
为此，注意到
$$
(x, z, w)^2 = (x^2, 2xz, 2xw + 2q(z, z))
$$
并且由归纳法
$$
(x, z, w)^n = (x^n, nx^{n - 1}z, nx^{n - 1}w + n(n - 1)x^{n - 2}q(z, z))
$$
另一方面，
$$
n! \delta_n(x, z, w) = (n!\gamma_n(x), n!\gamma_{n - 1}(x)z,
n!\gamma_{n - 1}(x)w + n!\gamma_{n - 2}(x) q(z, z))
$$
两式相符，证明完成。
\end{proof}




\section{兼容性}
\label{crystalline-section-compatibility}

\noindent
本节并非必读内容；它说明这里的讨论如何与 \cite{Berthelot} 中的讨论衔接。
考虑以下技术性概念。

\begin{definition}
\label{crystalline-definition-compatible}
设 $(A, I, \gamma)$ 与 $(B, J, \delta)$ 为除幂环，
$A \to B$ 为环同态。若 $J + IB$ 上存在除幂结构 $\bar\gamma$，使得
$$
(A, I, \gamma) \to (B, J + IB, \bar \gamma)\quad\text{且}\quad
(B, J, \delta) \to (B, J + IB, \bar \gamma)
$$
均为除幂环同态，则称{\it $\delta$ 与 $\gamma$ 兼容}。
\end{definition}

\noindent
设 $p$ 为素数，$(A,I,\gamma)$ 为除幂环，$A\to C$ 为环同态，
并且 $p$ 在 $C$ 中幂零。假设 $\gamma$ 可延拓至 $IC$（参见“除幂代数”
引理 \ref{dpa-lemma-gamma-extends}）。在此情形中，\cite{Berthelot}
所定义的 $\Spec(C)$ 在 $\Spec(A)$ 上的（大仿射）晶体位点，
是由以下系统组成的范畴的反范畴：
$$
(B, J, \delta, A \to B, C \to B/J)
$$
其中
\begin{enumerate}
\item $(B,J,\delta)$ 是除幂环，且 $p$ 在 $B$ 中幂零；
\item $\delta$ 与 $\gamma$ 兼容；
\item 下图
$$
\xymatrix{
B \ar[r] & B/J \\
A \ar[u] \ar[r] & C \ar[u]
}
$$
交换。
\end{enumerate}
在 \cite{Berthelot} 中，条件“$\gamma$ 可延拓至 $C$ 且 $\delta$ 与
$\gamma$ 兼容”用于保证 $\Spec(C)$ 的晶体上同调与
$\Spec(C/IC)$ 的晶体上同调相同。我们将只考虑满足 $IC=0$ 的 $C$，
从而避开这一问题\footnote{当然，这需要付出代价。}。在此情形中，
对上述系统 $(B,J,\delta,A\to B,C\to B/J)$，上图的交换性蕴含
$IB\subset J$；而兼容性等价于
$(A,I,\gamma)\to(B,J,\delta)$ 是除幂环同态。




\section{仿射晶体位点}
\label{crystalline-section-affine-site}

\noindent
本节讨论晶体位点的代数版本。以下是我们讨论这些内容时采用的基本情形。

\begin{situation}
\label{crystalline-situation-affine}
这里，$p$ 是素数，$(A,I,\gamma)$ 是除幂环，$A$ 是
$\mathbf{Z}_{(p)}$-代数；并且 $A\to C$ 是满足 $IC=0$ 且
$p$ 在 $C$ 中幂零的环同态。
\end{situation}

\noindent
通常，素数 $p$ 包含在除幂理想 $I$ 中。

\begin{definition}
\label{crystalline-definition-affine-thickening}
在情形 \ref{crystalline-situation-affine} 中。
\begin{enumerate}
\item $C$ 在 $(A,I,\gamma)$ 上的一个{\it 除幂加厚}，是一个除幂代数同态
$(A,I,\gamma)\to(B,J,\delta)$（其中 $p$ 在 $B$ 中幂零），
以及一个使下图交换的环同态 $C\to B/J$：
$$
\xymatrix{
B \ar[r] & B/J \\
& C \ar[u] \\
A \ar[uu] \ar[r] & A/I \ar[u]
}
$$
\item 一个{\it 除幂加厚同态}
$$
(B, J, \delta, C \to B/J) \longrightarrow (B', J', \delta', C \to B'/J')
$$
是除幂 $A$-代数同态 $\varphi:B\to B'$，使得复合映射
$C\to B/J\to B'/J'$ 等于给定映射 $C\to B'/J'$。
\item 以 $\text{CRIS}(C/A,I,\gamma)$，或简记为 $\text{CRIS}(C/A)$，
表示 $C$ 在 $(A,I,\gamma)$ 上的除幂加厚所组成的范畴。
\item 以 $\text{Cris}(C/A,I,\gamma)$，或简记为 $\text{Cris}(C/A)$，
表示由满足 $C\to B/J$ 为同构的对象 $(B,J,\delta,C\to B/J)$
组成的全子范畴。通常将这种对象记作 $(B\to C,\delta)$，其中约定
$J=\Ker(B\to C)$。
\end{enumerate}
\end{definition}

\noindent
注意，对上述除幂加厚 $(B,J,\delta)$，理想 $J$ 是局部幂零的；
参见“除幂代数”引理 \ref{dpa-lemma-nil}。有典范函子
\begin{equation}
\label{crystalline-equation-forget-affine}
\text{CRIS}(C/A) \longrightarrow C\text{-代数},\quad
(B, J, \delta) \longmapsto B/J
\end{equation}
一般而言，该范畴没有等化子或纤维积，也没有始对象（即空余极限）。

\begin{lemma}
\label{crystalline-lemma-affine-thickenings-colimits}
在情形 \ref{crystalline-situation-affine} 中。
\begin{enumerate}
\item $\text{CRIS}(C/A)$ 有有限积（但不一定有无限积）；
\item $\text{CRIS}(C/A)$ 有所有有限非空余极限，且
(\ref{crystalline-equation-forget-affine}) 与这些余极限可交换；
\item $\text{Cris}(C/A)$ 有所有有限非空余极限，且
$\text{Cris}(C/A) \to \text{CRIS}(C/A)$ 与这些余极限可交换。
\end{enumerate}
\end{lemma}

\begin{proof}
空积，即 $C$ 在 $(A,I,\gamma)$ 上的除幂加厚范畴中的终对象，
是如下零环：将其视为 $A$-代数，赋予零理想以及零理想上唯一的除幂，
再赋予从 $C$ 到零环的唯一同态。若
$(B_t,J_t,\delta_t)_{t\in T}$ 是 $\text{CRIS}(C/A)$ 中的一族对象，
则可按“除幂代数”引理 \ref{dpa-lemma-limits} 构造其积
$(\prod_t B_t,\prod_t J_t,\prod_t\delta_t)$。映射
$C\to\prod B_t/\prod J_t=\prod B_t/J_t$ 是显然的。然而，只有当
$T$ 有限时，才能保证 $p$ 在 $\prod_tB_t$ 中幂零。

\medskip\noindent
给定 $\text{CRIS}(C/A)$ 中的两个对象 $(B,J,\gamma)$ 与
$(B',J',\gamma')$，可以在除幂环范畴中构造余笛卡尔图
$$
\xymatrix{
(B, J, \delta) \ar[r] & (B'', J'', \delta'') \\
(A, I, \gamma) \ar[r] \ar[u] & (B', J', \delta') \ar[u]
}
$$
于是有
$$
B''/J'' = B/J \otimes_{A/I} B'/J' \longleftarrow C \otimes_{A/I} C
$$
参见“除幂代数”注 \ref{dpa-remark-pushout}。以
$J''\subset K\subset B''$ 表示使下图为推出的理想：
$$
\xymatrix{
B''/J'' \ar[r] & B''/K \\
C \otimes_{A/I} C \ar[r] \ar[u] & C \ar[u]
}
$$
也就是说，$B''/K\cong B/J\otimes_C B'/J'$。令
$D_{B''}(K)=(D,\bar K,\bar\delta)$ 为 $B''$ 中 $K$ 相对于
$(B'',J'',\delta'')$ 的除幂包络。容易验证，
$(D,\bar K,\bar\delta)$ 是 $(B,J,\delta)$ 与
$(B',J',\delta')$ 在 $\text{CRIS}(C/A)$ 中的余积。

\medskip\noindent
下面讨论余等化子。设
$\alpha,\beta:(B,J,\delta)\to(B',J',\delta')$ 为
$\text{CRIS}(C/A)$ 中的态射。考虑
$B''=B'/(\alpha(b)-\beta(b))$，并令 $J''\subset B''$ 为 $J'$ 的像。
令 $D_{B''}(J'')=(D,\bar J,\bar\delta)$ 为 $B''$ 中 $J''$ 相对于
$(B',J',\delta')$ 的除幂包络。容易验证，
$(D,\bar J,\bar\delta)$ 是 $(B,J,\delta)$ 与
$(B',J',\delta')$ 在 $\text{CRIS}(C/A)$ 中的余等化子。

\medskip\noindent
由“范畴”引理 \ref{categories-lemma-almost-finite-colimits-exist}，
$\text{CRIS}(C/A)$ 有所有有限非空余极限。上述构造表明
(\ref{crystalline-equation-forget-affine}) 与它们可交换。由于
$\text{Cris}(C/A)$ 是 (\ref{crystalline-equation-forget-affine}) 在 $C$ 上的纤维范畴，
这形式地蕴含 (3)。
\end{proof}

\begin{remark}
\label{crystalline-remark-completed-affine-site}
在情形 \ref{crystalline-situation-affine} 中，以 $\text{Cris}^\wedge(C/A)$
表示如下范畴：其对象是满足下列条件的对 $(B\to C,\delta)$：
\begin{enumerate}
\item $B$ 是 $p$-进完备的 $A$-代数；
\item $B \to C$ 是 $A$-代数满射；
\item $\delta$ 是 $\Ker(B \to C)$ 上的除幂结构；
\item $A \to B$ 是除幂环同态。
\end{enumerate}
态射按定义 \ref{crystalline-definition-affine-thickening} 定义。于是
$\text{Cris}(C/A)\subset\text{Cris}^\wedge(C/A)$ 是由满足
$p$ 在 $B$ 中幂零的对象组成的全子范畴。反过来，
$\text{Cris}^\wedge(C/A)$ 的任意对象 $(B\to C,\delta)$ 等于极限
$$
(B \to C, \delta) = \lim_e (B/p^eB \to C, \delta)
$$
其中当 $e\gg0$ 时，对象 $(B/p^eB\to C,\delta)$ 属于
$\text{Cris}(C/A)$；参见“除幂代数”引理
\ref{dpa-lemma-extend-to-completion}。特别地，
$\text{Cris}^\wedge(C/A)$ 是 $\text{Cris}(C/A)$ 的 pro-对象范畴的
全子范畴；参见“范畴”注 \ref{categories-remark-pro-category}。
\end{remark}

\begin{lemma}
\label{crystalline-lemma-list-properties}
在情形 \ref{crystalline-situation-affine} 中，设 $P\to C$ 为核等于 $J$ 的
$A$-代数满射，并记 $D_{P,\gamma}(J)=(D,\bar J,\bar\gamma)$。
令 $(D^\wedge,J^\wedge,\bar\gamma^\wedge)$ 为 $D$ 的 $p$-进完备化；
参见“除幂代数”引理 \ref{dpa-lemma-extend-to-completion}。
对每个 $e\geq1$，令 $P_e=P/p^eP$，令 $J_e\subset P_e$ 为 $J$ 的像，
并记 $D_{P_e,\gamma}(J_e)=(D_e,\bar J_e,\bar\gamma)$。
则对所有充分大的 $e$，有：
\begin{enumerate}
\item $p^eD\subset\bar J$ 与 $p^eD^\wedge\subset\bar J^\wedge$
均在除幂下保持不变；
\item 作为除幂环，$D^\wedge/p^eD^\wedge=D/p^eD=D_e$；
\item $(D_e,\bar J_e,\bar\gamma)$ 是 $\text{Cris}(C/A)$ 的对象；
\item $(D^\wedge,\bar J^\wedge,\bar\gamma^\wedge)$ 等于
$\lim_e(D_e,\bar J_e,\bar\gamma)$；
\item $(D^\wedge,\bar J^\wedge,\bar\gamma^\wedge)$ 是
$\text{Cris}^\wedge(C/A)$ 的对象。
\end{enumerate}
\end{lemma}

\begin{proof}
(1) 来自“除幂代数”引理 \ref{dpa-lemma-extend-to-completion}。
$p$-进完备化的一般性质给出
$D/p^eD=D^\wedge/p^eD^\wedge$。由于 $D/p^eD$ 是除幂环，且
$P\to D/p^eD$ 经过 $P_e$ 分解，$D_e$ 的泛性质给出映射
$D_e\to D/p^eD$。反过来，$D$ 的泛性质给出映射 $D\to D_e$，
它经过 $D/p^eD$ 分解。略去验证这两个映射互逆的步骤，从而得到 (2)。
当 $e$ 充分大时，$p^eC=0$，故 (3) 成立。(4) 来自“除幂代数”引理
\ref{dpa-lemma-extend-to-completion}，而 (5) 由定义显然成立。
\end{proof}

\begin{lemma}
\label{crystalline-lemma-set-generators}
在情形 \ref{crystalline-situation-affine} 中，设 $P$ 为 $A$ 上的多项式代数，
$P\to C$ 为核等于 $J$ 的 $A$-代数满射。取引理
\ref{crystalline-lemma-list-properties} 中的 $(D_e,\bar J_e,\bar\gamma)$。
对 $\text{CRIS}(C/A)$ 的每个对象 $(B,J_B,\delta)$，存在某个 $e$
以及 $\text{CRIS}(C/A)$ 中的态射 $D_e\to B$。
\end{lemma}

\begin{proof}
可以找到提升映射 $C\to B/J_B$ 的 $A$-代数同态 $P\to B$。
由 $\text{CRIS}(C/A)$ 的定义，对某个 $e$ 有 $p^eB=0$，故
$P\to B$ 分解为 $P\to P_e\to B$。由除幂包络的泛性质，
$P_e\to B$ 经过 $D_e$ 分解。
\end{proof}

\begin{lemma}
\label{crystalline-lemma-generator-completion}
在情形 \ref{crystalline-situation-affine} 中，设 $P$ 为 $A$ 上的多项式代数，
$P\to C$ 为核等于 $J$ 的 $A$-代数满射。令
$(D,\bar J,\bar\gamma)$ 为 $D_{P,\gamma}(J)$ 的 $p$-进完备化。
对 $\text{Cris}^\wedge(C/A)$ 的每个对象 $(B\to C,\delta)$，
存在 $\text{Cris}^\wedge(C/A)$ 中的态射 $D\to B$。
\end{lemma}

\begin{proof}
可以找到与到 $C$ 的映射兼容的 $A$-代数同态 $P\to B$。
由 $\text{Cris}^\wedge(C/A)$ 的定义，$P\to B$ 分解为
$P\to D_{P,\gamma}(J)\to B$。由于 $B$ 是 $p$-进完备的，
该映射可进一步经过 $D$ 分解。
\end{proof}



\section{微分模}
\label{crystalline-section-differentials}

\noindent
本节建立除幂环的微分模理论。

\begin{definition}
\label{crystalline-definition-derivation}
设 $A$ 为环，$(B,J,\delta)$ 为除幂环，$A\to B$ 为环同态，
$M$ 为 $B$-模。取值于 $M$ 的一个{\it 除幂 $A$-导子}，是满足如下条件的
映射 $\theta:B\to M$：它可加，在 $A$ 的元素上为零，满足莱布尼茨法则
$\theta(bb')=b\theta(b')+b'\theta(b)$，并且
$$
\theta(\delta_n(x)) = \delta_{n - 1}(x)\theta(x)
$$
对所有 $n\geq1$ 与 $x\in J$ 成立。
\end{definition}

\noindent
在该定义的情形中，与通常的导子一样，存在{\it 泛除幂 $A$-导子}
$$
\text{d}_{B/A, \delta} : B \to \Omega_{B/A, \delta}
$$
使得任意除幂 $A$-导子 $\theta:B\to M$ 都可唯一地写成
$\theta=\xi\circ d_{B/A,\delta}$，其中 $\xi:\Omega_{B/A,\delta}\to M$
为 $B$-线性映射。若
$(A,I,\gamma)\to(B,J,\delta)$ 是除幂环同态，则可忘掉 $A$ 上的除幂，
并考虑 $B$ 在 $A$ 上的除幂导子。下面给出（除幂）微分泛模的若干基本性质。

\begin{lemma}
\label{crystalline-lemma-omega}
设 $A$ 为环，$(B,J,\delta)$ 为除幂环，$A\to B$ 为环同态。
\begin{enumerate}
\item 考虑带有除幂理想 $(JB[x],\delta')$ 的 $B[x]$，其中
$\delta'$ 是 $\delta$ 到 $B[x]$ 的延拓。则
$$
\Omega_{B[x]/A, \delta'} =
\Omega_{B/A, \delta} \otimes_B B[x] \oplus B[x]\text{d}x.
$$
\item 考虑带有除幂理想
$(JB\langle x \rangle+B\langle x \rangle_{+},\delta')$ 的
$B\langle x \rangle$。则
$$
\Omega_{B\langle x\rangle/A, \delta'} =
\Omega_{B/A, \delta} \otimes_B B\langle x \rangle \oplus
B\langle x\rangle \text{d}x.
$$
\item 设 $K\subset J$ 为对所有 $n>0$ 均在 $\delta_n$ 下保持不变的理想。
令 $B'=B/K$，并以 $\delta'$ 表示 $J/K$ 上诱导的除幂。则
$\Omega_{B'/A,\delta'}$ 是 $\Omega_{B/A,\delta}\otimes_BB'$ 对如下
$B'$-子模的商：该子模由 $k\in K$ 时的 $\text{d}k$ 生成。
\end{enumerate}
\end{lemma}

\begin{proof}
这些结论直接来自 $\Omega_{B/A,\delta}$ 的构造：它是由元素
$\text{d}b$ 生成的自由 $B$-模对下列关系的商：
\begin{enumerate}
\item $\text{d}(b + b') = \text{d}b + \text{d}b'$, $b, b' \in B$,
\item $\text{d}a = 0$, $a \in A$,
\item $\text{d}(bb') = b \text{d}b' + b' \text{d}b$, $b, b' \in B$,
\item $\text{d}\delta_n(f) = \delta_{n - 1}(f)\text{d}f$, $f \in J$, $n > 1$.
\end{enumerate}
注意，最后一个关系解释了为何除幂多项式代数与通常的多项式代数给出
“相同”的答案：在前一种情形中，$x$ 是除幂理想的元素，因此
$\text{d}x^{[n]}=x^{[n-1]}\text{d}x$。
\end{proof}

\noindent
设 $(A,I,\gamma)$ 为除幂环。在这一背景下，$I$ 的幂应由如下除幂给出：
$$
I^{[n]} = \text{由下列元素生成的理想：}
\gamma_{e_1}(x_1) \ldots \gamma_{e_t}(x_t)
\text{ 其中 }\sum e_j \geq n\text{ 且 }x_j \in I.
$$
当然，$I^n\subset I^{[n]}$。注意 $I^{[1]}=I$；有时也令 $I^{[0]}=A$。

\begin{lemma}
\label{crystalline-lemma-diagonal-and-differentials}
设 $(A,I,\gamma)\to(B,J,\delta)$ 为除幂环同态。令
$(B(1),J(1),\delta(1))$ 为 $(B,J,\delta)$ 与自身在
$(A,I,\gamma)$ 上的余积，即使得下图
$$
\xymatrix{
(B, J, \delta) \ar[r] & (B(1), J(1), \delta(1)) \\
(A, I, \gamma) \ar[r] \ar[u] & (B, J, \delta) \ar[u]
}
$$
为余笛卡尔图。记 $K=\Ker(B(1)\to B)$。则
$K\cap J(1)\subset J(1)$ 在除幂结构下保持不变，并且典范地有
$$
\Omega_{B/A, \delta} = K/ \left(K^2 + (K \cap J(1))^{[2]}\right)
$$
\end{lemma}

\begin{proof}
$K\cap J(1)\subset J(1)$ 在除幂结构下保持不变，是因为
$B(1)\to B$ 是除幂环同态。

\medskip\noindent
回忆 $K/K^2$ 具有典范的 $B$-模结构。以
$s_0,s_1:B\to B(1)$ 表示两个余投影，并考虑由
$b\mapsto s_1(b)-s_0(b)$ 给出的映射
$\text{d}:B\to K/(K^2+(K\cap J(1))^{[2]})$。显然，$\text{d}$ 可加，
在 $A$ 上为零，并满足莱布尼茨法则。我们断言 $\text{d}$ 是除幂
$A$-导子。取 $x\in J$，令 $y=s_1(x)$、$z=s_0(x)$。下文以
$\delta$ 代替 $\delta(1)$ 表示 $J(1)$ 上的除幂结构。须证明对
$n\geq1$，模 $K^2+(K\cap J(1))^{[2]}$ 有
$\delta_n(y)-\delta_n(z)=\delta_{n-1}(y)(y-z)$。当 $n=1$ 时等式成立。
设 $n>1$。注意，当 $i>1$ 时，$\delta_i(y-z)$ 属于
$(K\cap J(1))^{[2]}$。模 $K^2+(K\cap J(1))^{[2]}$ 计算，得到
$$
\delta_n(z) = \delta_n(z - y + y) =
\sum\nolimits_{i = 0}^n \delta_i(z - y)\delta_{n - i}(y) =
\delta_{n - 1}(y) \delta_1(z - y) + \delta_n(y)
$$
这就证明了所需等式。

\medskip\noindent
设 $M$ 为 $B$-模，$\theta:B\to M$ 为除幂 $A$-导子。令
$D=B\oplus M$，其中 $M$ 是平方为零的理想。对 $n>1$，令
$\delta_n(x+m)=\delta_n(x)+\delta_{n-1}(x)m$，从而在
$J\oplus M\subset D$ 上定义除幂结构；参见引理
\ref{crystalline-lemma-divided-power-first-order-thickening}。有两个除幂代数同态
$B\to D$：第一个是包含映射，第二个由 $b\mapsto b+\theta(b)$ 给出。
因此，得到 $(A,I,\gamma)$ 上除幂代数的典范同态 $B(1)\to D$。
它诱导映射 $K\to M$；该映射在 $K^2$ 上为零（因为 $M$ 是平方为零的理想），
并且在 $(K\cap J(1))^{[2]}$ 上也为零，因为 $M^{[2]}=0$。
由构造，复合映射 $B\to K/K^2+(K\cap J(1))^{[2]}\to M$ 等于
$\theta$。因此 $\text{d}$ 是泛除幂 $A$-导子，从而得证。
\end{proof}

\begin{remark}
\label{crystalline-remark-filtration-differentials}
设 $A\to B$ 为环同态，$(J,\delta)$ 为 $B$ 上的除幂结构。
泛模 $\Omega_{B/A,\delta}$ 还带有一点额外结构：即由
$\text{d}_{B/A,\delta}(J)$ 生成的 $B$-子模
$N\subset\Omega_{B/A,\delta}$。在引理
\ref{crystalline-lemma-diagonal-and-differentials} 给出的同构下，它对应于
$K\cap J(1)$ 在 $\Omega_{B/A,\delta}$ 中的像。考虑 $A$-代数
$D=B\oplus\Omega^1_{B/A,\delta}$，其理想为 $\bar J=J\oplus N$，
除幂 $\bar\delta$ 如该引理证明中所定义。则
$(D,\bar J,\bar\delta)$ 是除幂环，并且由 $b\mapsto b$ 与
$b\mapsto b+\text{d}_{B/A,\delta}(b)$ 给出的两个映射 $B\to D$
都是 $A$ 上的除幂环同态。此外，$N$ 是使这一结论成立的
$\Omega_{B/A,\delta}$ 的最小子模。
\end{remark}

\begin{lemma}
\label{crystalline-lemma-diagonal-and-differentials-affine-site}
在情形 \ref{crystalline-situation-affine} 中，设 $(B,J,\delta)$ 为
$\text{CRIS}(C/A)$ 的对象。令 $(B(1),J(1),\delta(1))$ 为
$(B,J,\delta)$ 与自身在 $\text{CRIS}(C/A)$ 中的余积，并记
$K=\Ker(B(1)\to B)$。则 $K\cap J(1)\subset J(1)$ 在除幂结构下
保持不变，并且典范地有
$$
\Omega_{B/A, \delta} = K/ \left(K^2 + (K \cap J(1))^{[2]}\right)
$$
\end{lemma}

\begin{proof}
逐字照搬引理 \ref{crystalline-lemma-diagonal-and-differentials} 的证明即可。
唯一需要检验的是，除幂环 $D=B\oplus M$ 是
$\text{CRIS}(C/A)$ 的对象，且两个映射 $B\to D$ 是
$\text{CRIS}(C/A)$ 中的态射。由于 $D/(J\oplus M)=B/J$，
可以利用 $C\to B/J$ 将 $D$ 视为 $\text{CRIS}(C/A)$ 的对象；
关于态射的断言由构造显然成立。
\end{proof}

\begin{lemma}
\label{crystalline-lemma-module-differentials-divided-power-envelope}
设 $(A,I,\gamma)$ 为除幂环，$A\to B$ 为环同态，
$IB\subset J\subset B$ 为理想。令
$D_{B,\gamma}(J)=(D,\bar J,\bar\gamma)$ 为除幂包络。则
$$
\Omega_{D/A, \bar\gamma} = \Omega_{B/A} \otimes_B D
$$
\end{lemma}

\begin{proof}[第一种证明]
设 $M$ 为 $D$-模。我们断言，$A$-导子 $\vartheta:B\to M$
与除幂 $A$-导子 $\theta:D\to M$ 是同一回事。由米田引理，
该断言蕴含结论。

\medskip\noindent
考虑 $D$ 的平方零加厚 $D\oplus M$。若令 $M$ 上的高阶除幂运算为零，
则 $\bar J\oplus M$ 上有除幂结构 $\delta$。换言之，对任意
$x\in\bar J$ 与 $m\in M$，令
$\delta_n(x+m)=\bar\gamma_n(x)+\bar\gamma_{n-1}(x)m$；参见引理
\ref{crystalline-lemma-divided-power-first-order-thickening}。考虑 $A$-代数映射
$B\to D\oplus M$，其第一分量由 $B\to D$ 给出，第二分量为
$\vartheta$。由泛性质，得到对应的除幂代数同态
$D\to D\oplus M$；其第二分量正是与 $\vartheta$ 对应的除幂
$A$-导子 $\theta$。
\end{proof}

\begin{proof}[第二种证明]
先设 $B$ 在 $A$ 上平坦。此时 $\gamma$ 可延拓为 $IB$ 上的除幂结构
$\gamma'$；参见“除幂代数”引理 \ref{dpa-lemma-gamma-extends}。
因此 $D=D_{B,\gamma'}(J)$ 等于除幂环 $(D',J',\delta)$ 对元素
$x_t-f_t$ 与 $\delta_n(\sum r_t x_t-r_0)$ 所生成理想的商，其中
$D'=B\langle x_t\rangle$ 且
$J'=IB\langle x_t\rangle+B\langle x_t\rangle_+$；记号与说明参见引理
\ref{crystalline-lemma-describe-divided-power-envelope}。以
$\text{d}:D'\to\Omega_{D'/A,\delta}$ 表示泛导子。注意
$$
\Omega_{D'/A, \delta} =
\Omega_{B/A} \otimes_B D' \oplus \bigoplus D' \text{d}x_t,
$$
参见引理 \ref{crystalline-lemma-omega}。由此可知，$\Omega_{D/A,\bar\gamma}$
是 $\Omega_{D'/A,\delta}\otimes_{D'}D$ 对如下子模的商：该子模由
$\text{d}$ 作用于上列 $D'\to D$ 的核的生成元所得元素生成；
参见引理 \ref{crystalline-lemma-omega}。由于
$\text{d}(x_t-f_t)=-\text{d}f_t+\text{d}x_t$，在该商中有
$\text{d}x_t=\text{d}f_t$。特别地，映射
$\Omega_{B/A}\otimes_BD\to\Omega_{D/A,\gamma}$ 是满射；其核由
$\text{d}$ 作用于元素 $\delta_n(\sum r_t x_t-r_0)$ 所得的像给出。
然而，给定 $B$ 中的关系 $\sum r_tf_t-r_0=0$，其中
$r_t\in B$ 且 $r_0\in IB$，可知
\begin{align*}
\text{d}\delta_n(\sum r_t x_t - r_0)
& =
\delta_{n - 1}(\sum r_t x_t - r_0)\text{d}(\sum r_t x_t - r_0)
\\
& =
\delta_{n - 1}(\sum r_t x_t - r_0)
\left(
\sum r_t\text{d}(x_t - f_t) + \sum (x_t - f_t)\text{d}r_t
\right)
\end{align*}
因为 $\sum r_tf_t-r_0=0$ 在 $B$ 中成立。因此，它在
$\Omega_{B/A}\otimes_A D$ 中已经为零，从而当 $B$ 在 $A$ 上平坦时得证。

\medskip\noindent
在一般情形中，把 $B$ 写成多项式环的商 $P\to B$，并令
$J'\subset P$ 为 $J$ 的逆像。采用引理
\ref{crystalline-lemma-divided-power-envelop-quotient} 的记号，有 $D=D'/K'$。
由证明第一段处理的情形，
$\Omega_{D'/A,\bar\gamma'}=\Omega_{P/A}\otimes_PD'$。于是
$\Omega_{D/A,\bar\gamma}$ 是 $\Omega_{P/A}\otimes_PD$ 对如下子模的商：
该子模由 $k\in\Ker(P\to B)$ 时的 $\text{d}\bar\gamma_n'(k)$ 生成；
参见引理 \ref{crystalline-lemma-omega} 以及引理
\ref{crystalline-lemma-divided-power-envelop-quotient} 对 $K'$ 的描述。由于
$\text{d}\bar\gamma_n'(k)=\bar\gamma'_{n-1}(k)\text{d}k$，仍只需再对
$k\in\Ker(P\to B)$ 时的 $\text{d}k$ 所生成的子模取商。又因为
$\Omega_{B/A}$ 正是 $\Omega_{P/A}\otimes_AB$ 对这些元素所生成子模的商
（“代数”引理 \ref{algebra-lemma-differential-seq}），故得证。
\end{proof}

\begin{remark}
\label{crystalline-remark-divided-powers-de-rham-complex}
设 $A\to B$ 为环同态，$(J,\delta)$ 为 $B$ 上的除幂结构。令
$\Omega_{B/A,\delta}^i=\wedge_B^i\Omega_{B/A,\delta}$，其中
$\Omega_{B/A,\delta}$ 是泛除幂 $A$-导子
$\text{d}=\text{d}_{B/A}:B\to\Omega_{B/A,\delta}$ 的值域。注意，
$\Omega_{B/A,\delta}$ 是 $\Omega_{B/A}$ 对如下 $B$-子模的商：
该子模由 $x\in J$ 时的元素
$\text{d}\delta_n(x)-\delta_{n-1}(x)\text{d}x$ 生成。我们断言“代数”
引理 \ref{algebra-lemma-de-rham-complex} 适用。为此，只需验证
$\Omega_B$ 的元素
$\text{d}\delta_n(x)-\delta_{n-1}(x)\text{d}x$ 在
$\Omega^2_{B/A,\delta}$ 中的像为零。事实上，
$$
\text{d}(\delta_{n - 1}(x)) \wedge \text{d}x
= \delta_{n - 2}(x) \text{d}x \wedge \text{d}x = 0
$$
在 $\Omega^2_{B/A,\delta}$ 中成立，正合所需。因此得到{\it 除幂德拉姆复形}
$$
\Omega^0_{B/A, \delta} \to \Omega^1_{B/A, \delta} \to
\Omega^2_{B/A, \delta} \to \ldots
$$
它将在后文发挥重要作用。
\end{remark}

\begin{remark}
\label{crystalline-remark-connection}
设 $A\to B$ 为环同态，$\Omega_{B/A}\to\Omega$ 为满足“代数”引理
\ref{algebra-lemma-de-rham-complex} 之假设的商。设 $M$ 为 $B$-模。
一个{\it 联络}是可加映射
$$
\nabla : M \longrightarrow M \otimes_B \Omega
$$
使得对 $b\in B$ 与 $m\in M$ 有
$\nabla(bm)=b\nabla(m)+m\otimes\text{d}b$。在此情形中，可按规则
$$
\nabla : M \otimes_B \Omega^i \longrightarrow M \otimes_B \Omega^{i + 1}
$$
定义这些映射：
$\nabla(m\otimes\omega)=\nabla(m)\wedge\omega+m\otimes\text{d}\omega$。
这是良定义的，因为若 $b\in B$，则
\begin{align*}
\nabla(bm \otimes \omega) - \nabla(m \otimes b\omega)
& =
\nabla(bm) \wedge \omega + bm \otimes \text{d}\omega
- \nabla(m) \wedge b\omega - m \otimes \text{d}(b\omega) \\
& =
b\nabla(m) \wedge \omega + m \otimes \text{d}b \wedge \omega
+ bm \otimes \text{d}\omega \\
& \ \ \ \ \ \ - b\nabla(m) \wedge \omega - bm \otimes \text{d}(\omega)
- m \otimes \text{d}b \wedge \omega = 0
\end{align*}
依照惯例，当且仅当复合映射
$$
M \xrightarrow{\nabla} M \otimes_B \Omega^1
\xrightarrow{\nabla} M \otimes_B \Omega^2
$$
为零时，称该联络{\it 可积}。在这种情形下，得到复形
$$
M \xrightarrow{\nabla} M \otimes_B \Omega^1
\xrightarrow{\nabla} M \otimes_B \Omega^2
\xrightarrow{\nabla} M \otimes_B \Omega^3
\xrightarrow{\nabla} M \otimes_B \Omega^4 \to \ldots
$$
称为该联络的德拉姆复形。
\end{remark}

\begin{remark}
\label{crystalline-remark-base-change-connection}
考虑如下环的交换图：
$$
\xymatrix{
B \ar[r]_\varphi & B' \\
A \ar[u] \ar[r] & A' \ar[u]
}
$$
设 $\Omega_{B/A}\to\Omega$ 与 $\Omega_{B'/A'}\to\Omega'$ 为满足
“代数”引理 \ref{algebra-lemma-de-rham-complex} 之假设的商。
假设存在映射 $\varphi:\Omega\to\Omega'$，使下图交换：
$$
\xymatrix{
\Omega_{B/A} \ar[r] \ar[d] &
\Omega_{B'/A'} \ar[d] \\
\Omega \ar[r]^{\varphi} &
\Omega'
}
$$
其中上方水平箭头是由 $\varphi:B\to B'$ 诱导的典范映射
$\Omega_{B/A}\to\Omega_{B'/A'}$。在此情形中，给定任意一对
$(M,\nabla)$，其中 $M$ 为 $B$-模且
$\nabla:M\to M\otimes_B\Omega$ 为联络，可得到其{\it 基变换}
$(M\otimes_BB',\nabla')$，其中
$$
\nabla' :
M \otimes_B B'
\longrightarrow
(M \otimes_B B') \otimes_{B'} \Omega' = M \otimes_B \Omega'
$$
按下式定义：
$$
\nabla'(m \otimes b') =
\sum m_i \otimes b'\text{d}\varphi(b_i) + m \otimes \text{d}b' 
$$
这里 $\nabla(m)=\sum m_i\otimes\text{d}b_i$。若 $\nabla$ 可积，
则 $\nabla'$ 也可积；此时有德拉姆复形之间的典范映射（注
\ref{crystalline-remark-connection}）
\begin{equation}
\label{crystalline-equation-base-change-map-complexes}
M \otimes_B \Omega^\bullet
\longrightarrow
(M \otimes_B B') \otimes_{B'} (\Omega')^\bullet =
M \otimes_B (\Omega')^\bullet
\end{equation}
它把 $m\otimes\eta$ 映为 $m\otimes\varphi(\eta)$。
\end{remark}

\begin{lemma}
\label{crystalline-lemma-differentials-completion}
设 $A\to B$ 为环同态，$(J,\delta)$ 为 $B$ 上的除幂结构，$p$ 为素数。
假设 $A$ 是 $\mathbf{Z}_{(p)}$-代数，且 $p$ 在 $B/J$ 中幂零。则
$$
\lim_e \Omega_{B_e/A, \bar\delta} =
\lim_e \Omega_{B/A, \delta}/p^e\Omega_{B/A, \delta} =
\lim_e \Omega_{B^\wedge/A, \delta^\wedge}/p^e \Omega_{B^\wedge/A, \delta^\wedge}
$$
记号与说明参见证明。
\end{lemma}

\begin{proof}
由“除幂代数”引理 \ref{dpa-lemma-extend-to-completion}，对所有充分大的
$e$，$\delta$ 可延拓至 $B_e=B/p^eB$，故第一个极限有意义。
该引理还在完备化 $B^\wedge=\lim_eB_e$ 上给出除幂结构
$\delta^\wedge$，故最后一个极限也有意义。由引理 \ref{crystalline-lemma-omega}
以及恒有 $\text{d}p^e=0$，满射
$\Omega_{B/A,\delta}\to\Omega_{B_e/A,\bar\delta}$ 的核为
$p^e\Omega_{B/A,\delta}$。映射
$\Omega_{B^\wedge/A,\delta^\wedge}\to\Omega_{B_e/A,\bar\delta}$
的核也同理。因此结论显然。
\end{proof}



\section{除幂概形}
\label{crystalline-section-divided-power-schemes}

\noindent
下面说明如何将前述概念整体化。

\begin{definition}
\label{crystalline-definition-divided-power-structure}
设 $\mathcal C$ 为位点，$\mathcal O$ 为 $\mathcal C$ 上的环层，
$\mathcal I\subset\mathcal O$ 为理想层。$\mathcal I$ 上的一个
{\it 除幂结构 $\gamma$}，是映射序列
$\gamma_n:\mathcal I\to\mathcal I$（$n\geq1$），使得对
$\mathcal C$ 的任意对象 $U$，三元组
$$
(\mathcal{O}(U), \mathcal{I}(U), \gamma)
$$
为除幂环。
\end{definition}

\noindent
当然，这尤其适用于拓扑空间上的环层。不过，稍作一般化是有益的，
因为晶体位点的结构层正是定义在一个……位点上的！本章有时把上述定义中的
三元组 $(\mathcal C,\mathcal I,\gamma)$ 称为{\it 除幂拓扑斯}。
给定另一个 $(\mathcal C',\mathcal I',\gamma')$，以及环化拓扑斯的态射
$(f, f^\sharp) : (\Sh(\mathcal{C}), \mathcal{O}) \to
(\Sh(\mathcal{C}'), \mathcal{O}')$
若 $f^\sharp(f^{-1}\mathcal I')\subset\mathcal I$，且对所有 $n\geq1$
下图均交换，便称 $(f,f^\sharp)$ 诱导一个{\it 除幂拓扑斯态射}：
$$
\xymatrix{
f^{-1}\mathcal{I}' \ar[d]_{f^{-1}\gamma'_n} \ar[r]_{f^\sharp} &
\mathcal{I} \ar[d]^{\gamma_n} \\
f^{-1}\mathcal{I}' \ar[r]^{f^\sharp} & \mathcal{I}
}
$$
若 $f$ 来自由函子 $u:\mathcal C'\to\mathcal C$ 诱导的位点态射，
这恰好意味着对所有 $U'\in\Ob(\mathcal C')$，
$$
(\mathcal{O}'(U'), \mathcal{I}'(U'), \gamma')
\longrightarrow
(\mathcal{O}(u(U')), \mathcal{I}(u(U')), \gamma)
$$
均为除幂环同态。

\medskip\noindent
在概形的情形中，还要求除幂理想为{\bf 拟凝聚}的。除此以外，
定义与拓扑斯情形完全相同。具体如下。

\begin{definition}
\label{crystalline-definition-divided-power-scheme}
一个{\it 除幂概形}是三元组 $(S,\mathcal I,\gamma)$，其中 $S$ 是概形，
$\mathcal I$ 是拟凝聚理想层，而 $\gamma$ 是 $\mathcal I$ 上的除幂结构。
一个{\it 除幂概形态射}
$(S, \mathcal{I}, \gamma) \to (S', \mathcal{I}', \gamma')$，
是满足 $f^{-1}\mathcal I'\mathcal O_S\subset\mathcal I$ 的概形态射
$f:S\to S'$，并且
$$
(\mathcal{O}_{S'}(U'), \mathcal{I}'(U'), \gamma')
\longrightarrow
(\mathcal{O}_S(f^{-1}U'), \mathcal{I}(f^{-1}U'), \gamma)
$$
对每个开集 $U'\subset S'$ 都是除幂环同态。
\end{definition}

\noindent
回忆拟凝聚理想层与闭浸入之间存在一一对应，见
“态射”第 \ref{morphisms-section-closed-immersions} 节。
因此，给定除幂概形 $(T,\mathcal J,\gamma)$，由 $\mathcal J$ 定义出
一个典范闭浸入 $U\to T$。反过来，给定闭浸入 $U\to T$，以及与
$U\to T$ 相伴的理想层 $\mathcal J$ 上的除幂结构 $\gamma$，便得到
除幂概形 $(T,\mathcal J,\gamma)$。在许多情形中，我们只想考虑
$U\to T$ 为加厚时的这类三元组 $(U,T,\gamma)$；见
“态射进阶”定义 \ref{more-morphisms-definition-thickening}。

\begin{definition}
\label{crystalline-definition-divided-power-thickening}
若 $U\to T$ 是加厚，则称上述三元组 $(U,T,\gamma)$ 为一个
{\it 除幂加厚}。
\end{definition}

\noindent
当三个除幂概形中有一个是除幂加厚时，其纤维积存在。正式陈述如下。

\begin{lemma}
\label{crystalline-lemma-fibre-product}
设 $f:(T,\mathcal J,\delta)\to(S,\mathcal I,\gamma)$ 与
$f':(T',\mathcal J',\delta')\to(S,\mathcal I,\gamma)$ 是除幂概形态射。
则存在一个除幂概形 $(T'',\mathcal J'',\delta'')$ 以及笛卡尔图
$$
\xymatrix{
T \ar[d]_f & T'' \ar[d] \ar[l] \\
S & T' \ar[l]_{f'}
}
$$
于除幂概形范畴中成立。态射 $T''\to T\times_S T'$ 是闭浸入，
而态射 $T''_0\to T_0\times_{S_0}T'_0$ 是同构。
\end{lemma}

\begin{proof}
证明梗概。注意，借助 $\Spec(-)$，最后两个断言与
“除幂代数”注 \ref{dpa-remark-pushout}
中关于除幂代数推出的结论相容。因此我们如下把 $T''$ 构造成
$T\times_S T'$ 的闭子概形：对任意仿射开集 $U\subset S$、$V\subset T$、
$V'\subset T'$，若 $f(V),f'(V')\subset U$，便考虑由
“除幂代数”注 \ref{dpa-remark-pushout}
中的构造所确定的 $V\times_U V'$ 的闭子概形。由于各
$V\times_U V'$ 构成 $T\times_S T'$ 的一个开覆盖，可以依次进行：
(1) 证明这些闭子概形能够粘合；(2) 证明所得除幂结构能够粘合；
(3) 证明粘合所得对象是除幂概形范畴中的纤维积。
断言 (1) 等价于证明除幂环范畴中的推出构造与局部化交换
（适当表述后）；这由“除幂代数”引理
\ref{dpa-lemma-gamma-extends} 得出，因为该结果说明除幂可由平坦性
延拓到局部化上。
\end{proof}

\noindent
我们作如下观察。设 $(U,T,\gamma)$ 是除幂概形，$T$ 是
$\mathbf Z_{(p)}$ 上的概形，并且 $p$ 在 $U$ 上局部幂零。则
\begin{enumerate}
\item $p$ 在 $T$ 上局部幂零 $\Leftrightarrow U\to T$ 是加厚
（见“除幂代数”引理 \ref{dpa-lemma-nil}）；以及
\item 当 $e\gg0$ 时，$p^e\mathcal O_T$ 在 $T$ 上局部地被 $\gamma$ 保持
（见“除幂代数”引理
\ref{dpa-lemma-extend-to-completion}）。
\end{enumerate}
这表明，在以下假设下可以得到关于除幂加厚的良好结果。

\begin{situation}
\label{crystalline-situation-global}
这里 $p$ 是素数，$(S,\mathcal I,\gamma)$ 是 $\mathbf Z_{(p)}$ 上的
除幂概形。置 $S_0=V(\mathcal I)\subset S$。最后，$X\to S_0$
是概形态射，并且 $p$ 在 $X$ 上局部幂零。
\end{situation}

\noindent
我们将在此情形下定义大、小晶体位点。









\section{大晶体位点}
\label{crystalline-section-big-site}

\noindent
我们先定义大位点。给定除幂概形 $(S,\mathcal I,\gamma)$，若 $T$ 配备
一个除幂概形态射 $T\to S$，便称 $(T,\mathcal J,\delta)$ 是
$(S,\mathcal I,\gamma)$ 上的除幂概形。类似地，若 $T$ 配备一个
除幂概形态射 $T\to S$，便称除幂加厚 $(U,T,\delta)$ 是
$(S,\mathcal I,\gamma)$ 上的除幂加厚。

\begin{definition}
\label{crystalline-definition-divided-power-thickening-X}
在 情形 \ref{crystalline-situation-global} 中：
\begin{enumerate}
\item 一个{$X$ 相对于 $(S,\mathcal I,\gamma)$ 的\it 除幂加厚}，
由 $(S,\mathcal I,\gamma)$ 上的除幂加厚 $(U,T,\delta)$ 与一个
$S$-态射 $U\to X$ 给出。
\item {$X$ 相对于 $(S,\mathcal I,\gamma)$ 的\it 除幂加厚态射}
按显然方式定义。
\end{enumerate}
$X$ 相对于 $(S,\mathcal I,\gamma)$ 的除幂加厚所成的范畴记为
$\text{CRIS}(X/S,\mathcal I,\gamma)$，或简记为 $\text{CRIS}(X/S)$。
\end{definition}

\noindent
对 $\text{CRIS}(X/S)$ 中任意 $(U,T,\delta)$，$p$ 在 $T$ 上局部幂零；
见情形 \ref{crystalline-situation-global} 之前的讨论。与 $(U,T,\delta)$
相伴的全部数据可由下列交换图清楚表示：
$$
\xymatrix{
T \ar[dd] & U \ar[l] \ar[d] \\
& X \ar[d] \\
S & S_0 \ar[l]
}
$$
其中 $S_0=V(\mathcal I)\subset S$。$\text{CRIS}(X/S)$ 的态射也可类似地
表示成较大的交换图。特别地，有典范遗忘函子
\begin{equation}
\label{crystalline-equation-forget}
\text{CRIS}(X/S) \longrightarrow \Sch/X,\quad
(U, T, \delta) \longmapsto U
\end{equation}
以及它的单侧逆（同时也是左伴随）
\begin{equation}
\label{crystalline-equation-endow-trivial}
\Sch/X \longrightarrow \text{CRIS}(X/S),\quad
U \longmapsto (U, U, \emptyset)
\end{equation}
后者有时很有用。

\begin{lemma}
\label{crystalline-lemma-divided-power-thickening-fibre-products}
在情形 \ref{crystalline-situation-global} 中，范畴 $\text{CRIS}(X/S)$
具有一切非空有限极限，特别是二元积与纤维积。
函子 (\ref{crystalline-equation-forget}) 与这些极限交换。
\end{lemma}

\begin{proof}
略。提示：仿射情形见引理 \ref{crystalline-lemma-affine-thickenings-colimits}；
另见“除幂代数”注 \ref{dpa-remark-pushout}。
\end{proof}

\begin{lemma}
\label{crystalline-lemma-divided-power-thickening-base-change-flat}
在情形 \ref{crystalline-situation-global} 中，设
$$
\xymatrix{
(U_3, T_3, \delta_3) \ar[d] \ar[r] & (U_2, T_2, \delta_2) \ar[d] \\
(U_1, T_1, \delta_1) \ar[r] & (U, T, \delta)
}
$$
是 $X$ 相对于 $(S,\mathcal I,\gamma)$ 的除幂加厚范畴中的纤维方块。
若 $T_2\to T$ 平坦且 $U_2=T_2\times_TU$，则作为概形有
$T_3=T_1\times_TT_2$。
\end{lemma}

\begin{proof}
这是因为除幂结构沿平坦环同态可唯一延拓。见
“除幂代数”引理 \ref{dpa-lemma-gamma-extends}。
\end{proof}

\noindent
上述引理表明，除幂加厚的平坦态射经过基变换后仍是平坦态射，
并且事实上就是该态射的“通常”基变换。因此下述定义是合理的。

\begin{definition}
\label{crystalline-definition-big-crystalline-site}
在 情形 \ref{crystalline-situation-global} 中：
\begin{enumerate}
\item $X/S$ 的除幂加厚态射族
$\{(U_i,T_i,\delta_i)\to(U,T,\delta)\}$ 称为一个
{\it Zariski、\'etale、光滑、syntomic 或 fppf 覆盖}，当且仅当
\begin{enumerate}
\item 对所有 $i$ 都有 $U_i=U\times_TT_i$；并且
\item $\{T_i\to T\}$ 是 Zariski、\'etale、光滑、syntomic 或 fppf 覆盖。
\end{enumerate}
\item $X$ 在 $(S,\mathcal I,\gamma)$ 上的{\it 大晶体位点}，
是赋予 Zariski 拓扑的范畴 $\text{CRIS}(X/S)$。
\item $\text{CRIS}(X/S)$ 上的层拓扑斯记为
$(X/S)_{\text{CRIS}}$，有时也记为
$(X/S, \mathcal{I}, \gamma)_{\text{CRIS}}$\footnote{这与我们用
$\Sh(\mathcal C)$ 表示位点 $\mathcal C$ 所对应拓扑斯的约定冲突。}。
\end{enumerate}
\end{definition}

\noindent
这些拓扑斯具有一些显然的函子性。

\begin{remark}[函子性]
\label{crystalline-remark-functoriality-big-cris}
设 $p$ 是素数，并设
$(S,\mathcal I,\gamma)\to(S',\mathcal I',\gamma')$
是 $\mathbf Z_{(p)}$ 上的除幂概形态射。置
$S_0=V(\mathcal I)$ 与 $S'_0=V(\mathcal I')$。设
$$
\xymatrix{
X \ar[r]_f \ar[d] & Y \ar[d] \\
S_0 \ar[r] & S'_0
}
$$
是概形态射的交换图，并假设 $p$ 在 $X$ 与 $Y$ 上局部幂零。
则得到连续且余连续的函子
$$
\text{CRIS}(X/S) \longrightarrow \text{CRIS}(Y/S')
$$
它把 $(U,T,\delta)$ 对应到同一个 $(U,T,\delta)$，并以
$U\to X\to Y$ 作为从 $U$ 到 $Y$ 的 $S'$-态射。因此得到拓扑斯态射
$$
f_{\text{CRIS}} : (X/S)_{\text{CRIS}} \longrightarrow (Y/S')_{\text{CRIS}}
$$
见“位点论”第 \ref{sites-section-cocontinuous-morphism-topoi} 节。
\end{remark}

\begin{remark}[与 Zariski 位点的比较]
\label{crystalline-remark-compare-big-zariski}
在 情形 \ref{crystalline-situation-global} 中，函子 (\ref{crystalline-equation-forget})
是余连续的（细节略），并且与积和纤维积交换
（引理 \ref{crystalline-lemma-divided-power-thickening-fibre-products}）。
因此得到拓扑斯态射
$$
U_{X/S} : (X/S)_{\text{CRIS}} \longrightarrow \Sh((\Sch/X)_{Zar})
$$
从 $X/S$ 的大晶体拓扑斯指向 $X$ 的大 Zariski 拓扑斯。
见“位点论”第 \ref{sites-section-cocontinuous-morphism-topoi} 节。
\end{remark}

\begin{remark}[结构态射]
\label{crystalline-remark-big-structure-morphism}
在 情形 \ref{crystalline-situation-global} 中，考虑闭子概形
$S_0=V(\mathcal I)\subset S$。若假设 $p$ 在 $S_0$ 上局部幂零
（实际中总是如此），则以 $S_0$ 代替 $X$，便得到 定义
\ref{crystalline-definition-divided-power-thickening-X} 中的情形，从而得到位点
$\text{CRIS}(S_0/S)$。若 $f:X\to S_0$ 是 $X$ 在 $S$ 上的结构态射，
则得到环化拓扑斯态射的交换图
$$
\xymatrix{
(X/S)_{\text{CRIS}}
\ar[r]_{f_{\text{CRIS}}} \ar[d]_{U_{X/S}} &
(S_0/S)_{\text{CRIS}} \ar[d]^{U_{S_0/S}} \\
\Sh((\Sch/X)_{Zar}) \ar[r]^{f_{big}} & \Sh((\Sch/S_0)_{Zar}) \ar[rd] \\
& & \Sh((\Sch/S)_{Zar})
}
$$
这来自 注 \ref{crystalline-remark-functoriality-big-cris}。我们把复合
$(X/S)_{\text{CRIS}}\to\Sh((\Sch/S)_{Zar})$ 看作大晶体位点的结构态射。
即使 $p$ 在 $S_0$ 上并非局部幂零，结构态射
$$
(X/S)_{\text{CRIS}} \longrightarrow \Sh((\Sch/S)_{Zar})
$$
仍有定义，因为可以沿上图的下方路径得到它。因此，它是与下述余连续函子
$\text{CRIS}(X/S)\to(\Sch/S)_{Zar}$ 对应的拓扑斯态射：
$(U,T,\delta)/S\mapsto U/S$；见
“位点论”第 \ref{sites-section-cocontinuous-morphism-topoi} 节。
\end{remark}

\begin{remark}[相容性]
\label{crystalline-remark-compatibilities-big-cris}
上述态射满足许多相容性。例如，在 注
\ref{crystalline-remark-functoriality-big-cris} 的情形中，得到环化拓扑斯的交换图
$$
\xymatrix{
(X/S)_{\text{CRIS}} \ar[d] \ar[r] & (Y/S')_{\text{CRIS}} \ar[d] \\
\Sh((\Sch/S)_{Zar}) \ar[r] & \Sh((\Sch/S')_{Zar})
}
$$
其中竖直箭头为结构态射。
\end{remark}




\section{晶体位点}
\label{crystalline-section-site}

\noindent
由于 (\ref{crystalline-equation-forget}) 与积和纤维积交换，考察满足 $U\to X$
为开浸入的那些 $(U,T,\delta)$，便定义出一个在纤维积（更一般地，
非空有限极限）下保持的全子范畴。因此下述定义是合理的。

\begin{definition}
\label{crystalline-definition-crystalline-site}
在 情形 \ref{crystalline-situation-global} 中：
\begin{enumerate}
\item $X$ 在 $(S,\mathcal I,\gamma)$ 上的（小）{\it 晶体位点}，
记为 $\text{Cris}(X/S,\mathcal I,\gamma)$，或简记为
$\text{Cris}(X/S)$，是 $\text{CRIS}(X/S)$ 中由所有满足 $U\to X$
为开浸入的 $(U,T,\delta)$ 所成的全子范畴，并赋予 Zariski 拓扑。
\item $\text{Cris}(X/S)$ 上的层拓扑斯记为 $(X/S)_{\text{cris}}$，
有时也记为 $(X/S,\mathcal I,\gamma)_{\text{cris}}$\footnote{这与我们用
$\Sh(\mathcal C)$ 表示位点 $\mathcal C$ 所对应拓扑斯的约定冲突。}。
\end{enumerate}
\end{definition}

\noindent
对 $\text{Cris}(X/S)$ 中任意 $(U,T,\delta)$，态射 $U\to X$
定义了 $X$ 的小 Zariski 位点 $X_{Zar}$ 的一个对象。因此有典范遗忘函子
\begin{equation}
\label{crystalline-equation-forget-small}
\text{Cris}(X/S) \longrightarrow X_{Zar},\quad
(U, T, \delta) \longmapsto U
\end{equation}
以及左伴随
\begin{equation}
\label{crystalline-equation-endow-trivial-small}
X_{Zar} \longrightarrow \text{Cris}(X/S),\quad
U \longmapsto (U, U, \emptyset)
\end{equation}
后者有时很有用。

\medskip\noindent
正如可以比较一个概形的小、大 Zariski 位点一样，我们也可以比较
小、大晶体位点；见 “拓扑” 引理
\ref{topologies-lemma-at-the-bottom}。

\begin{lemma}
\label{crystalline-lemma-compare-big-small}
假设同定义 \ref{crystalline-definition-divided-power-thickening-X}。包含函子
$$
\text{Cris}(X/S) \to \text{CRIS}(X/S)
$$
与非空有限极限交换，并且是全忠实、连续且余连续的。存在拓扑斯态射
$$
(X/S)_{\text{cris}} \xrightarrow{i} (X/S)_{\text{CRIS}}
\xrightarrow{\pi} (X/S)_{\text{cris}}
$$
其复合为恒等态射，其中第一个由包含函子诱导。此外，$\pi_*=i^{-1}$。
\end{lemma}

\begin{proof}
第一个断言见引理 \ref{crystalline-lemma-divided-power-thickening-fibre-products}。
由此得到拓扑斯态射
$i:(X/S)_{\text{cris}}\to(X/S)_{\text{CRIS}}$，以及左伴随 $i_!$，
满足 $i^{-1}i_!=i^{-1}i_*=\text{id}$；见 “位点论” 引理
\ref{sites-lemma-when-shriek}、\ref{sites-lemma-preserve-equalizers}
及 \ref{sites-lemma-back-and-forth}。
我们断言 $i_!$ 是正合的。若此断言成立，便可由
$\pi^{-1}=i_!$ 与 $\pi_*=i^{-1}$ 定义 $\pi$，其余皆明。
为证明该断言，注意我们已知 $i_!$ 右正合且保持纤维积（见上述引文）。
因此只需证明 $i_!*=*$，其中 $*$ 表示集合层范畴的终对象。
为此，只需给出 $\text{Cris}(X/S)$ 的一族对象
$(U_i,T_i,\delta_i)$，$i\in I$，使得
$$
\coprod\nolimits_{i \in I} h_{(U_i, T_i, \delta_i)} \to *
$$
在 $(X/S)_{\text{CRIS}}$ 中满射（细节略；提示：利用
$\text{Cris}(X/S)$ 具有积，且函子
$\text{Cris}(X/S)\to\text{CRIS}(X/S)$ 与积交换）。
在仿射情形，这由引理 \ref{crystalline-lemma-set-generators} 得出。
一般情形的证明从略。
\end{proof}

\begin{remark}[函子性]
\label{crystalline-remark-functoriality-cris}
设 $p$ 是素数，并设
$(S,\mathcal I,\gamma)\to(S',\mathcal I',\gamma')$
是 $\mathbf Z_{(p)}$ 上的除幂概形态射。设
$$
\xymatrix{
X \ar[r]_f \ar[d] & Y \ar[d] \\
S_0 \ar[r] & S'_0
}
$$
是概形态射的交换图，并假设 $p$ 在 $X$ 与 $Y$ 上局部幂零。
仿照 “拓扑” 引理 \ref{topologies-lemma-morphism-big-small}，定义
$$
f_{\text{cris}} : (X/S)_{\text{cris}} \longrightarrow (Y/S')_{\text{cris}}
$$
公式为 $f_{\text{cris}}=\pi_Y\circ f_{\text{CRIS}}\circ i_X$；
其中 $i_X$ 与 $\pi_Y$ 分别是 引理 \ref{crystalline-lemma-compare-big-small}
中关于 $X$ 与 $Y$ 的态射，而 $f_{\text{CRIS}}$ 如
注 \ref{crystalline-remark-functoriality-big-cris} 所定义。
\end{remark}

\begin{remark}[与 Zariski 位点的比较]
\label{crystalline-remark-compare-zariski}
在 情形 \ref{crystalline-situation-global} 中，函子
(\ref{crystalline-equation-forget-small}) 连续、余连续，并且与积和纤维积交换。
因此得到拓扑斯态射
$$
u_{X/S} : (X/S)_{\text{cris}} \longrightarrow \Sh(X_{Zar})
$$
它联系 $X/S$ 的小晶体拓扑斯与 $X$ 的小 Zariski 拓扑斯。
见“位点论”第 \ref{sites-section-cocontinuous-morphism-topoi} 节。
\end{remark}

\begin{lemma}
\label{crystalline-lemma-localize}
在 情形 \ref{crystalline-situation-global} 中，设 $X'\subset X$ 与
$S'\subset S$ 是开子概形，并且 $X'$ 映入 $S'$。则存在全忠实函子
$\text{Cris}(X'/S')\to\text{Cris}(X/S)$，它诱导一个拓扑斯态射，
使下图交换：
$$
\xymatrix{
(X'/S')_{\text{cris}} \ar[r] \ar[d]_{u_{X'/S'}} &
(X/S)_{\text{cris}} \ar[d]^{u_{X/S}} \\
\Sh(X'_{Zar}) \ar[r] & \Sh(X_{Zar})
}
$$
此外，此图是 “位点论” 引理 \ref{sites-lemma-localize-morphism-topoi}
所述拓扑斯态射局部化的一个例子。
\end{lemma}

\begin{proof}
把 $\text{Cris}(X'/S')$ 的对象看作 $X$ 的除幂加厚
$(U,T,\delta)$，其中 $U\to X$ 通过 $X'\subset X$ 分解，便得到上述
全忠实函子（此时 $T\to S$ 自动通过 $S'$ 分解）。该函子显然余连续，
因而得到所示拓扑斯态射。把 $X'$ 看成 $X_{Zar}$ 的对象，并记其相伴的
可表层为 $h_{X'}\in\Sh(X_{Zar})$。显然 $\Sh(X'_{Zar})$ 是局部化
$\Sh(X_{Zar})/h_{X'}$。另一方面，借助上述函子，范畴
$\text{Cris}(X/S)/u_{X/S}^{-1}h_{X'}$
（见 “位点论” 引理 \ref{sites-lemma-localize-topos-site}）
与 $\text{Cris}(X'/S')$ 典范等同。证毕。
\end{proof}

\begin{remark}[结构态射]
\label{crystalline-remark-structure-morphism}
在 情形 \ref{crystalline-situation-global} 中，考虑闭子概形
$S_0=V(\mathcal I)\subset S$。若假设 $p$ 在 $S_0$ 上局部幂零
（实际中总是如此），则以 $S_0$ 代替 $X$，便得到 定义
\ref{crystalline-definition-divided-power-thickening-X} 中的情形，从而得到位点
$\text{Cris}(S_0/S)$。若 $f:X\to S_0$ 是 $X$ 在 $S$ 上的结构态射，
则得到拓扑斯的交换图
$$
\xymatrix{
(X/S)_{\text{cris}}
\ar[r]_{f_{\text{cris}}} \ar[d]_{u_{X/S}} &
(S_0/S)_{\text{cris}} \ar[d]^{u_{S_0/S}} \\
\Sh(X_{Zar}) \ar[r]^{f_{small}} & \Sh(S_{0, Zar}) \ar[rd] \\
& & \Sh(S_{Zar})
}
$$
见注 \ref{crystalline-remark-functoriality-cris}。我们把复合
$(X/S)_{\text{cris}}\to\Sh(S_{Zar})$ 看作晶体位点的结构态射。
即使 $p$ 在 $S_0$ 上并非局部幂零，结构态射
$$
\tau_{X/S} : (X/S)_{\text{cris}} \longrightarrow \Sh(S_{Zar})
$$
仍有定义，因为可以沿上图的下方路径得到它。
\end{remark}

\begin{remark}[相容性]
\label{crystalline-remark-compatibilities}
上述态射满足许多相容性。例如，在 注
\ref{crystalline-remark-functoriality-cris} 的情形中，得到环化拓扑斯的交换图
$$
\xymatrix{
(X/S)_{\text{cris}} \ar[d] \ar[r] & (Y/S')_{\text{cris}} \ar[d] \\
\Sh((\Sch/S)_{Zar}) \ar[r] & \Sh((\Sch/S')_{Zar})
}
$$
其中竖直箭头为结构态射。
\end{remark}



\section{晶体位点上的层}
\label{crystalline-section-sheaves}

\noindent
记号与假设同情形 \ref{crystalline-situation-global}。为了在本节同时讨论
$X/S$ 的小、大晶体位点，令
$$
\mathcal{C} = \text{CRIS}(X/S)
\quad\text{或}\quad
\mathcal{C} = \text{Cris}(X/S).
$$
$\mathcal C$ 上的层 $\mathcal F$ 对 $\mathcal C$ 的每个对象
$(U,T,\delta)$ 给出一个{\it 限制} $\mathcal F_T$。具体地，
$\mathcal F_T$ 是概形 $T$ 上由下式定义的 Zariski 层：
$$
\mathcal{F}_T(W) = \mathcal{F}(U \cap W, W, \delta|_W)
$$
其中 $W\subset T$ 为开集。此外，若 $f:T\to T'$ 是 $\mathcal C$
的对象 $(U,T,\delta)$ 与 $(U',T',\delta')$ 之间的态射，则有典范
{\it 比较}映射
\begin{equation}
\label{crystalline-equation-comparison}
c_f : f^{-1}\mathcal{F}_{T'} \longrightarrow \mathcal{F}_T.
\end{equation}
事实上，若 $W'\subset T'$ 为开集，则 $f$ 诱导 $\mathcal C$ 中的态射
$$
f|_{f^{-1}W'} :
(U \cap f^{-1}(W'), f^{-1}W', \delta|_{f^{-1}W'})
\longrightarrow
(U' \cap W', W', \delta|_{W'})
$$
因而可用 $\mathcal F$ 的限制映射 $(f|_{f^{-1}W'})^*$ 定义映射
$\mathcal F_{T'}(W')\to\mathcal F_T(f^{-1}W')$。这些映射显然与进一步
限制相容，故定义了从 $\mathcal F_{T'}$ 到 $\mathcal F_T$ 的一个
$f$-映射（见“层”第 \ref{sheaves-section-presheaves-functorial} 节，
尤其是 “层” 定义 \ref{sheaves-definition-f-map}），
从而得到 (\ref{crystalline-equation-comparison}) 中的映射 $c_f$。
注意，若 $f$ 是开浸入，则 $c_f$ 是同构，因为此时 $\mathcal F_T$
正是 $\mathcal F_{T'}$ 在 $T$ 上的限制。

\medskip\noindent
反过来，若对 $\mathcal C$ 的每个对象 $(U,T,\delta)$ 给定 Zariski 层
$\mathcal F_T$，并给定上述比较映射 $c_f$，且它们 (a) 对开浸入为同构，
(b) 满足适当的上闭链条件，则得到 $\mathcal C$ 上的一个层。其证明与
“拓扑” 引理 \ref{topologies-lemma-characterize-sheaf-big} 完全相同。

\medskip\noindent
$\mathcal C$ 上的{\it 结构层}是由下式定义的层 $\mathcal O_{X/S}$：
$$
\mathcal{O}_{X/S} :
(U, T, \delta)
\longmapsto
\Gamma(T, \mathcal{O}_T)
$$
由 $\mathcal C$ 中覆盖的定义可知这是一个层。设 $\mathcal F$ 是
$\mathcal O_{X/S}$-模层。此时比较映射 (\ref{crystalline-equation-comparison})
定义一个 $\mathcal O_T$-模的比较映射
\begin{equation}
\label{crystalline-equation-comparison-modules}
c_f : f^*\mathcal{F}_{T'} \longrightarrow \mathcal{F}_T
\end{equation}

\medskip\noindent
另一类例子从 $(\Sch/X)_{Zar}$ 或 $X_{Zar}$ 上的层 $\mathcal G$ 出发
（取决于 $\mathcal C=\text{CRIS}(X/S)$ 还是
$\mathcal C=\text{Cris}(X/S)$）。则由下式定义的
$\underline{\mathcal G}$ 是 $\mathcal C$ 上的层：
$$
\underline{\mathcal{G}} :
(U, T, \delta)
\longmapsto
\mathcal{G}(U)
$$
特别地，取 $\mathcal G=\mathbf G_a=\mathcal O_X$，得到
$$
\underline{\mathbf{G}_a} :
(U, T, \delta)
\longmapsto
\Gamma(U, \mathcal{O}_U)
$$
对每个对象 $(U,T,\delta)$，典范映射
$\Gamma(T,\mathcal O_T)\to\Gamma(U,\mathcal O_U)$ 定义层的满射
$\mathcal O_{X/S}\to\underline{\mathbf G_a}$。记其核为
$\mathcal J_{X/S}$，从而得到短正合列
$$
0 \to
\mathcal{J}_{X/S} \to
\mathcal{O}_{X/S} \to
\underline{\mathbf{G}_a} \to 0
$$
注意，$\mathcal J_{X/S}$ 配备典范除幂结构。事实上，对每个对象
$(U,T,\delta)$，第三个分量 $\delta$ 本来就是
$\mathcal O_T\to\mathcal O_U$ 之核上的除幂结构。因此（大）晶体拓扑斯
是除幂拓扑斯。





\section{模中的晶体}
\label{crystalline-section-crystals}

\noindent
晶体原来是一种非常一般的对象。不过其定义可能稍难理解，因此我们先在
晶体位点上的模这一情形中给出定义。

\begin{definition}
\label{crystalline-definition-modules}
在情形 \ref{crystalline-situation-global} 中，令
$\mathcal C=\text{CRIS}(X/S)$ 或 $\mathcal C=\text{Cris}(X/S)$，
并令 $\mathcal F$ 是 $\mathcal C$ 上的 $\mathcal O_{X/S}$-模层。
\begin{enumerate}
\item 若对 $\mathcal C$ 的每个对象 $(U,T,\delta)$，限制 $\mathcal F_T$
都是拟凝聚 $\mathcal O_T$-模，则称 $\mathcal F$ {\it 局部拟凝聚}。
\item 若 $\mathcal F$ 在“位点上的模”定义
\ref{sites-modules-definition-site-local} 的意义下拟凝聚，则称
$\mathcal F$ {\it 拟凝聚}。
\item 若所有比较映射 (\ref{crystalline-equation-comparison-modules}) 都是同构，
则称 $\mathcal F$ 是一个{$\mathcal O_{X/S}$-模中的\it 晶体}。
\end{enumerate}
\end{definition}

\noindent
这些概念之间有如下关系。

\begin{lemma}
\label{crystalline-lemma-crystal-quasi-coherent-modules}
采用 定义 \ref{crystalline-definition-modules} 中的记号
$X/S,\mathcal I,\gamma,\mathcal C,\mathcal F$。以下条件等价：
\begin{enumerate}
\item $\mathcal F$ 拟凝聚；
\item $\mathcal F$ 局部拟凝聚，并且是 $\mathcal O_{X/S}$-模中的晶体。
\end{enumerate}
\end{lemma}

\begin{proof}
假设 (1)。设 $f:(U',T',\delta')\to(U,T,\delta)$ 是 $\mathcal C$ 中的态射。
我们要证明：(a) $\mathcal F_T$ 是拟凝聚 $\mathcal O_T$-模；
(b) $c_f:f^*\mathcal F_T\to\mathcal F_{T'}$ 是同构。该假设意味着可以
找到一个覆盖 $\{(T_i,U_i,\delta_i)\to(T,U,\delta)\}$，使得对每个 $i$，
$\mathcal F$ 在 $\mathcal C/(T_i,U_i,\delta_i)$ 上的限制具有全局表现。
由于只需在 Zariski 局部证明 (a) 与 (b)，可以把
$f:(T',U',\delta')\to(T,U,\delta)$ 换成到 $(T_i,U_i,\delta_i)$ 的基变换，
并假设 $\mathcal F$ 在 $\mathcal C/(T,U,\delta)$ 上的限制具有全局表现
$$
\bigoplus\nolimits_{j \in J}
\mathcal{O}_{X/S}|_{\mathcal{C}/(U, T, \delta)} \longrightarrow
\bigoplus\nolimits_{i \in I}
\mathcal{O}_{X/S}|_{\mathcal{C}/(U, T, \delta)} \longrightarrow
\mathcal{F}|_{\mathcal{C}/(U, T, \delta)}
\longrightarrow 0
$$
显然，这给出表现
$$
\bigoplus\nolimits_{j \in J} \mathcal{O}_T \longrightarrow
\bigoplus\nolimits_{i \in I} \mathcal{O}_T \longrightarrow
\mathcal{F}_T
\longrightarrow 0
$$
故 (a) 成立。此外，该表现限制到 $T'$ 后给出 $\mathcal F_{T'}$
的同类表现，故 (b) 成立。

\medskip\noindent
假设 (2)。设 $(U,T,\delta)$ 是 $\mathcal C$ 的对象。我们要找
$(U,T,\delta)$ 的一个覆盖，使 $\mathcal F$ 限制到 $\mathcal C$
在覆盖各成员处的局部化后具有全局表现。因此可以假设 $T$ 仿射。
此时可选取表现
$$
\bigoplus\nolimits_{j \in J} \mathcal{O}_T \longrightarrow
\bigoplus\nolimits_{i \in I} \mathcal{O}_T \longrightarrow
\mathcal{F}_T
\longrightarrow 0
$$
因为已假设 $\mathcal F_T$ 是拟凝聚 $\mathcal O_T$-模。由
$\mathcal F$ 的晶体性质，对 $\mathcal C$ 的任意态射
$f:(U',T',\delta')\to(U,T,\delta)$，此表现拉回为 $\mathcal F_{T'}$
的表现。于是得到 $\mathcal F|_{\mathcal C/(U,T,\delta)}$ 所需的表现。
\end{proof}

\begin{definition}
\label{crystalline-definition-crystal-quasi-coherent-modules}
若 $\mathcal F$ 满足 引理 \ref{crystalline-lemma-crystal-quasi-coherent-modules}
中的等价条件，则称 $\mathcal F$ 是一个{\it 拟凝聚模中的晶体}。
若此外 $\mathcal F$ 还是有限局部自由的，则称它是一个
{\it 有限局部自由模中的晶体}。
\end{definition}

\noindent
当然，引理 \ref{crystalline-lemma-crystal-quasi-coherent-modules} 表明这种说法略显
累赘，因为拟凝聚模总是晶体。不过这是文献中的标准术语。

\begin{remark}
\label{crystalline-remark-crystal}
为表述晶体的一般概念，我们采用栈与强笛卡尔态射的语言；见
“栈” 定义 \ref{stacks-definition-stack} 及 “范畴”
定义 \ref{categories-definition-cartesian-over-C}。
在 情形 \ref{crystalline-situation-global} 中，设
$p:\mathcal C\to\text{Cris}(X/S)$ 是一个栈。一个
{$X$ 上相对于 $S$ 的、取值于 $\mathcal C$ 对象的\it 晶体}，
是一个{\it 笛卡尔截面} $\sigma:\text{Cris}(X/S)\to\mathcal C$；
亦即函子 $\sigma$ 满足 $p\circ\sigma=\text{id}$，且对
$\text{Cris}(X/S)$ 的每个态射 $f$，$\sigma(f)$ 都是强笛卡尔的。
大晶体位点的情形类似。
\end{remark}





\section{微分层}
\label{crystalline-section-differentials-sheaf}

\noindent
本节只使用（小）晶体位点，因为这似乎更自然。我们如下整体化
定义 \ref{crystalline-definition-derivation}。

\begin{definition}
\label{crystalline-definition-global-derivation}
在 情形 \ref{crystalline-situation-global} 中，设 $\mathcal F$ 是
$\text{Cris}(X/S)$ 上的 $\mathcal O_{X/S}$-模层。一个
{$S$-{\it 导子} $D:\mathcal O_{X/S}\to\mathcal F$} 是层的映射，
使得对 $\text{Cris}(X/S)$ 的每个对象 $(U,T,\delta)$，映射
$$
D : \Gamma(T, \mathcal{O}_T) \longrightarrow \Gamma(T, \mathcal{F})
$$
是除幂 $\Gamma(V,\mathcal O_V)$-导子，其中 $V\subset S$ 是任意开集，
使得 $T\to S$ 通过 $V$ 分解。
\end{definition}

\noindent
这意味着 $D$ 可加，满足 Leibniz 法则，消去来自 $S$ 的函数，并且对
除幂理想 $\mathcal J_{X/S}$ 的局部截面 $f$ 满足
$D(f^{[n]})=f^{[n-1]}D(f)$。这是下面将描述的一个非常一般概念的特例。

\medskip\noindent
请把以下讨论与“位点上的模”第
\ref{sites-modules-section-differentials} 节比较。设 $\mathcal C$ 是位点，
$\mathcal A\to\mathcal B$ 是 $\mathcal C$ 上环层的映射，
$\mathcal J\subset\mathcal B$ 是理想层，$\delta$ 是 $\mathcal J$
上的除幂结构，而 $\mathcal F$ 是 $\mathcal B$-模层。此时可以定义
{\it 除幂 $\mathcal A$-导子} $D:\mathcal B\to\mathcal F$：这意味着
$D$ 是 $\mathcal A$-线性的，满足 Leibniz 法则，并且对 $\mathcal J$
的局部截面 $x$ 满足 $D(\delta_n(x))=\delta_{n-1}(x)D(x)$。
在这种情形中，存在一个{\it 泛除幂 $\mathcal A$-导子}
$$
\text{d}_{\mathcal{B}/\mathcal{A}, \delta} :
\mathcal{B}
\longrightarrow
\Omega_{\mathcal{B}/\mathcal{A}, \delta}
$$
此外，$\text{d}_{\mathcal B/\mathcal A,\delta}$ 是复合
$$
\mathcal{B}
\longrightarrow
\Omega_{\mathcal{B}/\mathcal{A}}
\longrightarrow
\Omega_{\mathcal{B}/\mathcal{A}, \delta}
$$
其中第一个映射是在“位点上的模”引理
\ref{sites-modules-lemma-universal-module} 的证明中构造的泛导子，
第二个箭头则是对由下列局部截面生成的子模取商：
$\text{d}_{\mathcal{B}/\mathcal{A}}(\delta_n(x)) -
\delta_{n - 1}(x)\text{d}_{\mathcal{B}/\mathcal{A}}(x)$.

\medskip\noindent
我们把它转化为如下相对概念。设
$(f,f^\sharp):(\Sh(\mathcal C),\mathcal O)\to
(\Sh(\mathcal C'),\mathcal O')$ 是环化拓扑斯态射，
$\mathcal J\subset\mathcal O$ 是理想层，$\delta$ 是 $\mathcal J$
上的除幂结构，而 $\mathcal F$ 是 $\mathcal O$-模层。若 $D$ 是上述意义下
的除幂 $f^{-1}\mathcal O'$-导子，便称 $D:\mathcal O\to\mathcal F$
为除幂 $\mathcal O'$-导子。此外，记
$$
\Omega_{\mathcal{O}/\mathcal{O}', \delta} =
\Omega_{\mathcal{O}/f^{-1}\mathcal{O}', \delta}
$$
它是泛除幂 $\mathcal O'$-导子的取值模。

\medskip\noindent
把这应用于结构态射
$$
(X/S)_{\text{Cris}} \longrightarrow \Sh(S_{Zar})
$$
（见注 \ref{crystalline-remark-structure-morphism}），便恢复上述 定义
\ref{crystalline-definition-global-derivation} 的概念。特别地，存在泛除幂导子
$$
d_{X/S} : \mathcal{O}_{X/S} \to \Omega_{X/S}
$$
注意，我们在记号中省略了表示微分模与除幂相容的修饰符（似乎不太可能
有人会考虑晶体位点结构层的通常微分模）。

\begin{lemma}
\label{crystalline-lemma-module-differentials-divided-power-scheme}
设 $(T,\mathcal J,\delta)$ 是除幂概形，并设 $T\to S$ 是概形态射。
上述商 $\Omega_{T/S}\to\Omega_{T/S,\delta}$ 是拟凝聚
$\mathcal O_T$-模。对映入仿射开集 $V\subset S$ 的仿射开集
$W\subset T$，有
$$
\Gamma(W, \Omega_{T/S, \delta}) =
\Omega_{\Gamma(W, \mathcal{O}_W)/\Gamma(V, \mathcal{O}_V), \delta}
$$
其中右端如第 \ref{crystalline-section-differentials} 节中所构造。
\end{lemma}

\begin{proof}
略。
\end{proof}

\begin{lemma}
\label{crystalline-lemma-module-of-differentials}
在 情形 \ref{crystalline-situation-global} 中，对
$\text{Cris}(X/S)$ 中的 $(U,T,\delta)$，$\Omega_{X/S}$ 在 $T$ 上的限制
$(\Omega_{X/S})_T$ 等于 $\Omega_{T/S,\delta}$，而限制
$\text{d}_{X/S}|_T$ 等于 $\text{d}_{T/S,\delta}$。
\end{lemma}

\begin{proof}
略。
\end{proof}

\begin{lemma}
\label{crystalline-lemma-module-of-differentials-on-affine}
在 情形 \ref{crystalline-situation-global} 中，对 $\text{Cris}(X/S)$
中映入仿射开集 $V\subset S$ 的任意仿射对象 $(U,T,\delta)$，有
$$
\Gamma((U, T, \delta), \Omega_{X/S}) =
\Omega_{\Gamma(T, \mathcal{O}_T)/\Gamma(V, \mathcal{O}_V), \delta}
$$
其中右端如第 \ref{crystalline-section-differentials} 节中所构造。
\end{lemma}

\begin{proof}
合用 引理 \ref{crystalline-lemma-module-differentials-divided-power-scheme}
与 \ref{crystalline-lemma-module-of-differentials}。
\end{proof}

\begin{lemma}
\label{crystalline-lemma-describe-omega-small}
在 情形 \ref{crystalline-situation-global} 中，设 $(U,T,\delta)$ 是
$\text{Cris}(X/S)$ 的对象。令
$$
(U(1), T(1), \delta(1)) = (U, T, \delta) \times (U, T, \delta)
$$
为 $\text{Cris}(X/S)$ 中的积。令
$\mathcal K\subset\mathcal O_{T(1)}$ 是与闭浸入
$\Delta:T\to T(1)$ 对应的拟凝聚理想层。则
$\mathcal K\subset\mathcal J_{T(1)}$ 被 $\mathcal J_{T(1)}$
上的除幂结构保持，并且
$$
(\Omega_{X/S})_T = \mathcal{K}/\mathcal{K}^{[2]}
$$
\end{lemma}

\begin{proof}
由于 $U\to X$ 是开浸入，且 (\ref{crystalline-equation-forget-small}) 与积交换，
故 $U=U(1)$，从而 $\mathcal K\subset\mathcal J_{T(1)}$。有了这一事实，
在 $T$ 上仿射局部地工作，并使用 引理
\ref{crystalline-lemma-module-of-differentials-on-affine} 与
\ref{crystalline-lemma-diagonal-and-differentials-affine-site}，即得结论。
\end{proof}

\noindent
$\Omega_{X/S}$ 并不是拟凝聚 $\mathcal O_{X/S}$-模中的晶体，
但它满足两个紧密相关的性质（与 引理
\ref{crystalline-lemma-crystal-quasi-coherent-modules} 比较）。

\begin{lemma}
\label{crystalline-lemma-omega-locally-quasi-coherent}
在 情形 \ref{crystalline-situation-global} 中，微分层 $\Omega_{X/S}$
具有以下两个性质：
\begin{enumerate}
\item $\Omega_{X/S}$ 局部拟凝聚；
\item 对 $\text{Cris}(X/S)$ 中任意态射
$(U,T,\delta)\to(U',T',\delta')$，若 $f:T\to T'$ 是闭浸入，
则映射 $c_f:f^*(\Omega_{X/S})_{T'}\to(\Omega_{X/S})_T$ 满射。
\end{enumerate}
\end{lemma}

\begin{proof}
(1) 由引理 \ref{crystalline-lemma-module-differentials-divided-power-scheme}
与 \ref{crystalline-lemma-module-of-differentials} 合用得出。
(2) 由以下事实得出：$(\Omega_{X/S})_T=\Omega_{T/S,\delta}$
是 $\Omega_{T/S}$ 的商，且
$f^*\Omega_{T'/S}\to\Omega_{T/S}$ 满射。
\end{proof}







\section{两个泛加厚}
\label{crystalline-section-universal-thickenings}

\noindent
本节的构造将帮助我们在晶体位点上的模晶体上定义联络。在某种意义下，
这里的构造是在第 \ref{crystalline-section-explicit-thickenings} 节中给出的
“层化、泛化”版本。

\begin{remark}
\label{crystalline-remark-first-order-thickening}
在 情形 \ref{crystalline-situation-global} 中，设 $(U,T,\delta)$ 是
$\text{Cris}(X/S)$ 的对象。记
$\Omega_{T/S,\delta}=(\Omega_{X/S})_T$；见
引理 \ref{crystalline-lemma-module-of-differentials}。我们显式描述 $T$
的一个一阶加厚 $T'$。具体地，置
$$
\mathcal{O}_{T'} = \mathcal{O}_T \oplus \Omega_{T/S, \delta}
$$
并赋予代数结构，使 $\Omega_{T/S,\delta}$ 是平方为零的理想。
令 $\mathcal J\subset\mathcal O_T$ 是闭浸入 $U\to T$ 的理想层，置
$\mathcal J'=\mathcal J\oplus\Omega_{T/S,\delta}$。由下式在
$\mathcal J'$ 上定义除幂结构：
$$
\delta_n'(f, \omega) = (\delta_n(f), \delta_{n - 1}(f)\omega),
$$
见引理 \ref{crystalline-lemma-divided-power-first-order-thickening}。有两个环映射
$$
p_0, p_1 : \mathcal{O}_T \to \mathcal{O}_{T'}
$$
第一个由 $f\mapsto(f,0)$ 给出，第二个由
$f\mapsto(f,\text{d}_{T/S,\delta}f)$ 给出。注意二者都与
$\mathcal J$ 和 $\mathcal J'$ 上的除幂结构相容，商映射
$\mathcal O_{T'}\to\mathcal O_T$ 亦然。因此得到
$\text{Cris}(X/S)$ 的对象 $(U,T',\delta')$ 及交换图
$$
\xymatrix{
& T \ar[ld]_{\text{id}} \ar[d]^i \ar[rd]^{\text{id}} \\
T & T' \ar[l]_{p_0} \ar[r]^{p_1} & T
}
$$
于 $\text{Cris}(X/S)$ 中成立；其中 $i$ 是一阶加厚，其理想层与
$\Omega_{T/S,\delta}$ 等同，并且
$p_1-p_0:\mathcal O_T\to\mathcal O_{T'}$ 与泛导子
$\text{d}_{T/S,\delta}$ 复合包含
$\Omega_{T/S,\delta}\to\mathcal O_{T'}$ 后的映射等同。
\end{remark}

\begin{remark}
\label{crystalline-remark-second-order-thickening}
在 情形 \ref{crystalline-situation-global} 中，设 $(U,T,\delta)$ 是
$\text{Cris}(X/S)$ 的对象。记
$\Omega_{T/S,\delta}=(\Omega_{X/S})_T$；见
引理 \ref{crystalline-lemma-module-of-differentials}。并以
$\Omega^2_{T/S,\delta}$ 表示其二次外幂。我们显式描述 $T$
的一个二阶加厚 $T''$。具体地，置
$$
\mathcal{O}_{T''} =
\mathcal{O}_T \oplus \Omega_{T/S, \delta} \oplus \Omega_{T/S, \delta}
\oplus \Omega^2_{T/S, \delta}
$$
其代数结构定义如下：
$$
(f, \omega_1, \omega_2, \eta) \cdot
(f', \omega_1', \omega_2', \eta') =
(ff', f\omega_1' + f'\omega_1, f\omega_2' + f'\omega_2,
f\eta' + f'\eta + \omega_1 \wedge \omega_2' + \omega_1' \wedge \omega_2).
$$
令 $\mathcal J\subset\mathcal O_T$ 是闭浸入 $U\to T$ 的理想层，
并令 $\mathcal J''$ 是 $\mathcal J$ 在投影
$\mathcal O_{T''}\to\mathcal O_T$ 下的逆像。由下式在
$\mathcal J''$ 上定义除幂结构：
$$
\delta_n''(f, \omega_1, \omega_2, \eta) =
(\delta_n(f), \delta_{n - 1}(f)\omega_1, \delta_{n - 1}(f)\omega_2,
\delta_{n - 1}(f)\eta + \delta_{n - 2}(f)\omega_1 \wedge \omega_2)
$$
见引理 \ref{crystalline-lemma-divided-power-second-order-thickening}。有三个环映射
$q_0,q_1,q_2:\mathcal O_T\to\mathcal O_{T''}$，定义为
\begin{align*}
q_0(f) & = (f, 0, 0, 0), \\
q_1(f) & = (f, \text{d}f, 0, 0), \\
q_2(f) & = (f, \text{d}f, \text{d}f, 0)
\end{align*}
其中 $\text d=\text d_{T/S,\delta}$。注意三者都与 $\mathcal J$
及 $\mathcal J''$ 上的除幂结构相容。另有三个环映射
$q_{01},q_{12},q_{02}:\mathcal O_{T'}\to\mathcal O_{T''}$，
其中 $\mathcal O_{T'}$ 如 注 \ref{crystalline-remark-first-order-thickening}
所定义。具体地，置
\begin{align*}
q_{01}(f, \omega) & = (f, \omega, 0, 0), \\
q_{12}(f, \omega) & =
(f, \text{d}f, \omega, \text{d}\omega), \\
q_{02}(f, \omega) & = (f, \omega, \omega, 0)
\end{align*}
这些映射也与给定的除幂结构相容。我们验证 $q_{12}$：注意
$q_{12}$ 是环同态，因为
\begin{align*}
q_{12}(f, \omega)q_{12}(g, \eta) & =
(f, \text{d}f, \omega, \text{d}\omega)(g, \text{d}g, \eta, \text{d}\eta) \\
& =
(fg, f\text{d}g + g \text{d}f, f\eta + g\omega,
f\text{d}\eta + g\text{d}\omega + \text{d}f \wedge \eta +
\text{d}g \wedge \omega) \\
& = q_{12}(fg, f\eta + g\omega) = q_{12}((f, \omega)(g, \eta))
\end{align*}
并且 $q_{12}$ 与除幂相容，因为
\begin{align*}
\delta_n''(q_{12}(f, \omega)) & =
\delta_n''((f, \text{d}f, \omega, \text{d}\omega)) \\
& =
(\delta_n(f), \delta_{n - 1}(f)\text{d}f, \delta_{n - 1}(f)\omega,
\delta_{n - 1}(f)\text{d}\omega + \delta_{n - 2}(f)\text{d}(f) \wedge \omega)
\\
& = q_{12}((\delta_n(f), \delta_{n - 1}(f)\omega)) =
q_{12}(\delta'_n(f, \omega))
\end{align*}
对 $q_{01}$ 与 $q_{02}$ 的验证更为容易。注意
$q_0=q_{01}\circ p_0$、$q_1=q_{01}\circ p_1$、
$q_1=q_{12}\circ p_0$、$q_2=q_{12}\circ p_1$、
$q_0=q_{02}\circ p_0$，以及 $q_2=q_{02}\circ p_1$。
因此 $(U,T'',\delta'')$ 是 $\text{Cris}(X/S)$ 的对象，并得到态射
$$
\xymatrix{
T''
\ar@<2ex>[r]
\ar@<0ex>[r]
\ar@<-2ex>[r]
&
T'
\ar@<1ex>[r]
\ar@<-1ex>[r]
&
T
}
$$
于 $\text{Cris}(X/S)$ 中满足上述关系。在应用中，我们用
$q_i:T''\to T$ 与 $q_{ij}:T''\to T'$ 表示与上述环映射相伴的态射。
\end{remark}






\section{德拉姆复形}
\label{crystalline-section-de-Rham}

\noindent
在 情形 \ref{crystalline-situation-global} 中，我们在（小）晶体位点上工作，
对 $i\geq0$ 定义
$\Omega^i_{X/S}=\wedge^i_{\mathcal O_{X/S}}\Omega_{X/S}$。
泛 $S$-导子 $\text d_{X/S}$ 给出 $\text{Cris}(X/S)$ 上的
{\it 德拉姆复形}
$$
\mathcal{O}_{X/S} \to \Omega^1_{X/S} \to \Omega^2_{X/S} \to \ldots
$$
见引理 \ref{crystalline-lemma-module-of-differentials-on-affine} 与
注 \ref{crystalline-remark-divided-powers-de-rham-complex}。


\section{联络}
\label{crystalline-section-connections}

\noindent
在 情形 \ref{crystalline-situation-global} 中，给定 $\text{Cris}(X/S)$
上的 $\mathcal O_{X/S}$-模 $\mathcal F$，一个{\it 联络}是交换群层的映射
$$
\nabla :
\mathcal{F}
\longrightarrow
\mathcal{F} \otimes_{\mathcal{O}_{X/S}} \Omega_{X/S}
$$
使得对 $\mathcal F$ 与 $\mathcal O_{X/S}$ 的局部截面 $s,f$，有
$\nabla(fs)=f\nabla(s)+s\otimes\text df$。给定联络，存在典范映射
$
\nabla :
\mathcal{F} \otimes_{\mathcal{O}_{X/S}} \Omega^i_{X/S}
\longrightarrow
\mathcal{F} \otimes_{\mathcal{O}_{X/S}} \Omega^{i + 1}_{X/S}
$
其定义为 $\nabla(s\otimes\omega)=\nabla(s)\wedge\omega+
s\otimes\text d\omega$，同注 \ref{crystalline-remark-connection}。
若 $\nabla\circ\nabla=0$，则称该联络{\it 可积}。若 $\nabla$ 可积，
便得到 $\text{Cris}(X/S)$ 上的{\it 德拉姆复形}
$$
\mathcal{F} \to
\mathcal{F} \otimes_{\mathcal{O}_{X/S}} \Omega^1_{X/S} \to
\mathcal{F} \otimes_{\mathcal{O}_{X/S}} \Omega^2_{X/S} \to \ldots
$$
任意 $\mathcal O_{X/S}$-模中的晶体都典范地配备一个可积联络。

\begin{lemma}
\label{crystalline-lemma-automatic-connection}
在 情形 \ref{crystalline-situation-global} 中，设 $\mathcal F$ 是
$\text{Cris}(X/S)$ 上的 $\mathcal O_{X/S}$-模中的晶体。
则 $\mathcal F$ 典范地配备一个可积联络。
\end{lemma}

\begin{proof}
设 $(U,T,\delta)$ 是 $\text{Cris}(X/S)$ 的对象。令
$(U,T',\delta')$ 是 注 \ref{crystalline-remark-first-order-thickening}
中用 $(\Omega_{X/S})_T=\Omega_{T/S,\delta}$ 构造的 $T$ 的无穷小加厚。
它配备投影 $p_0,p_1:T'\to T$ 与对角映射 $i:T\to T'$。
由假设得到 $\mathcal O_{T'}$-模的同构
$$
p_0^*\mathcal{F}_T \xrightarrow{c_0}
\mathcal{F}_{T'} \xleftarrow{c_1}
p_1^*\mathcal{F}_T
$$
把 $c=c_1^{-1}\circ c_0$ 沿 $i$ 拉回到 $T$，得到 $\mathcal F_T$
的恒等映射。因此，若 $s\in\Gamma(T,\mathcal F_T)$，则
$\nabla(s)=p_1^*s-c(p_0^*s)$ 是 $p_1^*\mathcal F_T$ 的截面，沿 $i$
拉回后为零。所以 $\nabla(s)$ 是下列模的截面：
$$
\mathcal{F}_T
\otimes_{\mathcal{O}_T}
\Omega_{T/S, \delta}
$$
因为按构造 $\mathcal O_{T'}=\mathcal O_T\oplus\Omega_{T/S,\delta}$，
这正是 $p_1^*\mathcal F_T\to\mathcal F_T$ 的核。利用注
\ref{crystalline-remark-first-order-thickening} 中对 $\text d$ 的描述，容易验证
$\nabla(fs)=f\nabla(s)+s\otimes\text d(f)$。

\medskip\noindent
由此得到的映射族
$$
\nabla : \Gamma(T, \mathcal{F}_T) \to
\Gamma(T, \mathcal{F}_T \otimes_{\mathcal{O}_T} \Omega_{T/S, \delta})
$$
关于 $T$ 具有函子性，因为 $T'$ 的构造关于 $T$ 具有函子性。
因此得到一个联络。

\medskip\noindent
为证明该联络可积，考虑 注 \ref{crystalline-remark-second-order-thickening}
中构造的对象 $(U,T'',\delta'')$。由于 $\mathcal F$ 是层，可知
$$
\xymatrix{
q_0^*\mathcal{F}_T \ar[rr]_{q_{01}^*c} \ar[rd]_{q_{02}^*c} & &
q_1^*\mathcal{F}_T \ar[ld]^{q_{12}^*c} \\
& q_2^*\mathcal{F}_T
}
$$
是 $\mathcal O_{T''}$-模的交换图。对
$s\in\Gamma(T,\mathcal F_T)$，有
$c(p_0^*s)=p_1^*s-\nabla(s)$。写成
$\nabla(s)=\sum p_1^*s_i\cdot\omega_i$，其中 $s_i$ 是
$\mathcal F_T$ 的局部截面，$\omega_i$ 是 $\Omega_{T/S,\delta}$
的局部截面。把 $\omega_i$ 看作 $\mathcal O_{T'}$ 的结构层的局部截面，
所以这里写乘积而非张量积。一方面
\begin{align*}
q_{12}^*c \circ q_{01}^*c(q_0^*s) & = 
q_{12}^*c(q_1^*s - \sum q_1^*s_i \cdot q_{01}^*\omega_i) \\
& =
q_2^*s - \sum q_2^*s_i \cdot q_{12}^*\omega_i -
\sum q_2^*s_i \cdot q_{01}^*\omega_i +
\sum q_{12}^*\nabla(s_i) \cdot q_{01}^*\omega_i
\end{align*}
另一方面
$$
q_{02}^*c(q_0^*s) = q_2^*s - \sum q_2^*s_i \cdot q_{02}^*\omega_i.
$$
由注 \ref{crystalline-remark-second-order-thickening} 中的公式可知
$q_{01}^*\omega_i + q_{12}^*\omega_i - q_{02}^*\omega_i = \text{d}\omega_i$.
因此上述两个表达式之差为
$$
\sum q_2^*s_i \cdot \text{d}\omega_i -
\sum q_{12}^*\nabla(s_i) \cdot q_{01}^*\omega_i
$$
由 $\mathcal O_{T''}$ 上乘法的定义，注意
$q_{12}^*\omega\cdot q_{01}^*\omega'=\omega'\wedge\omega
=-\omega\wedge\omega'$。因此，上式正是 $\nabla^2(s)$，看作
$q_2^*\mathcal F$ 的子层
$\mathcal F_T\otimes\Omega^2_{T/S,\delta}$ 的截面。
交换性遂给出可积条件。
\end{proof}




\section{余单纯代数}
\label{crystalline-section-cosimplicial}

\noindent
本节应移至别处。一个{\it 余单纯环}是环范畴中的余单纯对象。
给定环 $R$，一个{\it 余单纯 $R$-代数}是 $R$-代数范畴中的余单纯对象。
余单纯环 $A_*$ 中的一个{\it 余单纯理想}，是对每个 $n$ 给定理想
$I_n\subset A_n$，使得对 $\Delta$ 中所有 $f:[n]\to[m]$，都有
$A(f)(I_n)\subset I_m$。

\medskip\noindent
设 $A_*$ 是余单纯环。令 $\mathcal C$ 为如下偶对 $(A,M)$ 所成的范畴：
$A$ 是环，$M$ 是 $A$-模。态射 $(A,M)\to(A',M')$ 由环映射
$A\to A'$ 与 $A$-模映射 $M\to M'$ 组成；这里借助 $A\to A'$ 及
$M'$ 的 $A'$-模结构，把 $M'$ 看作 $A$-模。于是可以把
{$A_*$ 上的\it 余单纯模 $M_*$} 定义为 $\mathcal C$ 的余单纯对象
$(A_*,M_*)$，其第一分量等于 $A_*$。{$A_*$ 上余单纯模的一个\it 同态
$\varphi_*:M_*\to N_*$}，是 $\mathcal C$ 中余单纯对象的态射
$(A_*,M_*)\to(A_*,N_*)$，其第一分量为 $1_{A_*}$。

\medskip\noindent
在 $A_*$ 上余单纯模同态 $\varphi_*,\psi_*:M_*\to N_*$ 之间的
一个{\it 同伦}，是相伴映射 $(A_*,M_*)\to(A_*,N_*)$ 之间的同伦，
其第一分量为平凡同伦（与 “单纯方法” 例
\ref{simplicial-example-trivial-homotopy} 对偶）。具体而言，这样的同伦是
$$
h : M_* \longrightarrow \Hom(\Delta[1], N_*)
$$
作为余单纯交换群同态的 $\varphi_*$ 与 $\psi_*$ 之间的同伦，
并且对每个 $n$，映射
$h_n:M_n\to\prod_{\alpha\in\Delta[1]_n}N_n$ 是 $A_n$-线性的。
下述引理是 “单纯方法” 引理
\ref{simplicial-lemma-functorial-homotopy} 对余单纯模的版本。

\begin{lemma}
\label{crystalline-lemma-homotopy-tensor}
设 $A_*$ 是余单纯环，并设 $\varphi_*,\psi_*:K_*\to M_*$
是余单纯 $A_*$-模的同态。
\begin{enumerate}
\item
\label{crystalline-item-tensor}
若 $\varphi_*$ 与 $\psi_*$ 同伦，则对任意余单纯 $A_*$-模 $L_*$，
$$
\varphi_* \otimes 1, \psi_* \otimes 1 :
K_* \otimes_{A_*} L_* \longrightarrow M_* \otimes_{A_*} L_*
$$
同伦。
\item
\label{crystalline-item-wedge}
若 $\varphi_*$ 与 $\psi_*$ 同伦，则
$$
\wedge^i(\varphi_*), \wedge^i(\psi_*) :
\wedge^i(K_*) \longrightarrow \wedge^i(M_*)
$$
同伦。
\item
\label{crystalline-item-base-change}
若 $\varphi_*$ 与 $\psi_*$ 同伦，且 $A_*\to B_*$ 是余单纯环同态，则
$$
\varphi_* \otimes 1, \psi_* \otimes 1 :
K_* \otimes_{A_*} B_* \longrightarrow M_* \otimes_{A_*} B_*
$$
作为余单纯 $B_*$-模同态是同伦的。
\item
\label{crystalline-item-completion}
若 $I_*\subset A_*$ 是余单纯理想，则完备化之间的诱导映射
$$
\varphi^\wedge_*, \psi^\wedge_* :
K_*^\wedge \longrightarrow M_*^\wedge
$$
同伦。
\item 可按需在此添加更多情形，例如对称幂。
\end{enumerate}
\end{lemma}

\begin{proof}
设 $h:M_*\longrightarrow\Hom(\Delta[1],N_*)$ 是给定同伦。在次数 $n$，有
$$
h_n = (h_{n, \alpha}) :
K_n \longrightarrow
\prod\nolimits_{\alpha \in \Delta[1]_n} K_n
$$
见“单纯方法”第 \ref{simplicial-section-homotopy-cosimplicial} 节。
一族 $h_{n,\alpha}$ 构成同伦的充要条件是：对每个
$f:[n]\to[m]$，都有
$$
h_{m, \alpha} \circ M_*(f) = N_*(f) \circ h_{n, \alpha \circ f}
$$
见 “单纯方法” Equation
(\ref{simplicial-equation-property-homotopy-cosimplicial})。
此外还应有 $\psi_n=h_{n,0:[n]\to[1]}$ 与
$\varphi_n=h_{n,1:[n]\to[1]}$。

\medskip\noindent
在引理的各情形中，我们都可给出相应映射。
情形 (\ref{crystalline-item-tensor})：使用同伦 $h\otimes1$，其次数 $n$ 分量定义为
$$
(h \otimes 1)_{n, \alpha} = h_{n, \alpha} \otimes 1_{L_n} :
K_n \otimes_{A_n} L_n
\longrightarrow
M_n \otimes_{A_n} L_n.
$$
情形 (\ref{crystalline-item-wedge})：使用同伦 $\wedge^ih$，其次数 $n$ 分量定义为
$$
\wedge^i(h)_{n, \alpha} = \wedge^i(h_{n, \alpha}) :
\wedge_{A_n}(K_n)
\longrightarrow
\wedge^i_{A_n}(M_n).
$$
情形 (\ref{crystalline-item-base-change})：使用同伦 $h\otimes1$，其次数 $n$ 分量定义为
$$
(h \otimes 1)_{n, \alpha} = h_{n, \alpha} \otimes 1 :
K_n \otimes_{A_n} B_n
\longrightarrow
M_n \otimes_{A_n} B_n.
$$
情形 (\ref{crystalline-item-completion})：使用同伦 $h^\wedge$，其次数 $n$ 分量定义为
$$
(h^\wedge)_{n, \alpha} = h_{n, \alpha}^\wedge :
K_n^\wedge
\longrightarrow
M_n^\wedge.
$$
这是成立的，因为每个 $h_{n,\alpha}$ 都是 $A_n$-线性的。
\end{proof}






\section{拟凝聚模中的晶体}
\label{crystalline-section-quasi-coherent-crystals}

\noindent
在 情形 \ref{crystalline-situation-affine} 中，置 $X=\Spec(C)$ 与
$S=\Spec(A)$。我们将分类 $\text{Cris}(X/S)$ 上拟凝聚模中的晶体。
在此之前先固定一些记号。

\medskip\noindent
选取 $A$ 上多项式环 $P=A[x_i]$ 及 $A$-代数满射 $P\to C$，其核为
$J=\Ker(P\to C)$。置
\begin{equation}
\label{crystalline-equation-D}
D = \lim_e D_{P, \gamma}(J) / p^eD_{P, \gamma}(J)
\end{equation}
这是 $p$-进完备化的除幂包络。该环配备除幂理想 $\bar J$ 与除幂结构
$\bar\gamma$；见引理 \ref{crystalline-lemma-list-properties}。置
$D_e=D/p^eD$，并以 $\bar J_e$ 表示 $\bar J$ 在 $D_e$ 中的像。
我们用简记
\begin{equation}
\label{crystalline-equation-omega-D}
\Omega_D = \lim_e \Omega_{D_e/A, \bar\gamma} =
\lim_e \Omega_{D/A, \bar\gamma}/p^e\Omega_{D/A, \bar\gamma}
\end{equation}
表示除幂微分模的 $p$-进完备化；见引理
\ref{crystalline-lemma-differentials-completion}。它也是
$\Omega_{D_{P,\gamma}(J)/A,\bar\gamma}$ 的 $p$-进完备化，
而后者以 $\text d x_i$ 为自由基；见引理
\ref{crystalline-lemma-module-differentials-divided-power-envelope}。
因此，$\Omega_D$ 的任意元素可唯一写成和
$\sum f_i\text d x_i$，并且对每个 $e$，只有有限多个 $f_i$
不属于 $p^eD$。此外，各映射
$\text d_{D_e/A,\bar\gamma}:D_e\to\Omega_{D_e/A,\bar\gamma}$
拼合成 $p$-进完备化上的一个除幂 $A$-导子
\begin{equation}
\label{crystalline-equation-derivation-D}
\text{d} : D \longrightarrow \Omega_D
\end{equation}

\medskip\noindent
我们还需要“$\Spec(D)$ 的积 $\Spec(D(n))$”；解释见命题
\ref{crystalline-proposition-compute-cohomology} 及其证明。其形式定义如下。
对 $n\geq0$，令
$J(n)=\Ker(P\otimes_A\ldots\otimes_AP\to C)$，其中张量积有
$n+1$ 个因子。置
\begin{equation}
\label{crystalline-equation-Dn}
D(n) = \lim_e
D_{P \otimes_A \ldots \otimes_A P, \gamma}(J(n))/
p^eD_{P \otimes_A \ldots \otimes_A P, \gamma}(J(n))
\end{equation}
即除幂包络的 $p$-进完备化。以 $\bar J(n)$ 表示其除幂理想，
以 $\bar\gamma(n)$ 表示其除幂。另引入
$D(n)_e=D(n)/p^eD(n)$、$p$-进完备化的微分模
\begin{equation}
\label{crystalline-equation-omega-Dn}
\Omega_{D(n)} = \lim_e \Omega_{D(n)_e/A, \bar\gamma} =
\lim_e \Omega_{D(n)/A, \bar\gamma}/p^e\Omega_{D(n)/A, \bar\gamma}
\end{equation}
以及导子
\begin{equation}
\label{crystalline-equation-derivation-Dn}
\text{d} : D(n) \longrightarrow \Omega_{D(n)}
\end{equation}
当然 $D=D(0)$。注意，各环 $D(0),D(1),D(2),\ldots$
构成除幂环范畴中的一个余单纯对象。

\begin{lemma}
\label{crystalline-lemma-structure-Dn}
设 $D$ 与 $D(n)$ 如 (\ref{crystalline-equation-D}) 和 (\ref{crystalline-equation-Dn})。
余积典范映射 $P\to P\otimes_A\ldots\otimes_AP$，
$f\mapsto f\otimes1\otimes\ldots\otimes1$，诱导 $D$-代数同构
\begin{equation}
\label{crystalline-equation-structure-Dn}
D(n) = \lim_e D\langle \xi_i(j) \rangle/p^eD\langle \xi_i(j) \rangle
\end{equation}
其中
$$
\xi_i(j) = x_i \otimes 1 \otimes \ldots \otimes 1 -
1 \otimes \ldots \otimes 1 \otimes x_i \otimes 1 \otimes \ldots \otimes 1
$$
对 $j=1,\ldots,n$，第二个 $x_i$ 位于第 $j+1$ 个位置；回忆 $D(n)$
从 $P$ 在 $A$ 上的 $n+1$ 重张量积出发构造。
\end{lemma}

\begin{proof}
我们有
$$
P \otimes_A \ldots \otimes_A P = P[\xi_i(j)]
$$
而 $J(n)$ 由 $J$ 与各元素 $\xi_i(j)$ 生成。因此引理由
引理 \ref{crystalline-lemma-divided-power-envelope-add-variables} 得出。
\end{proof}

\begin{lemma}
\label{crystalline-lemma-property-Dn}
设 $D$ 与 $D(n)$ 如 (\ref{crystalline-equation-D}) 和 (\ref{crystalline-equation-Dn})。
则 $(D,\bar J,\bar\gamma)$ 与
$(D(n),\bar J(n),\bar\gamma(n))$ 是 $\text{Cris}^\wedge(C/A)$
的对象（见注 \ref{crystalline-remark-completed-affine-site}），并且
$$
D(n) = \coprod\nolimits_{j = 0, \ldots, n} D
$$
为 $\text{Cris}^\wedge(C/A)$ 中的余积。
\end{lemma}

\begin{proof}
第一个断言显然。对于第二个断言，若 $(B\to C,\delta)$ 是
$\text{Cris}^\wedge(C/A)$ 的对象，则
$$
\Mor_{\text{Cris}^\wedge(C/A)}(D, B) = 
\Hom_A((P, J), (B, \Ker(B \to C)))
$$
对 $D(n)$ 也有类似等式，其中以
$(P\otimes_A\ldots\otimes_AP,J(n))$ 代替 $(P,J)$。
由于 $P\otimes_A\ldots\otimes_AP$ 是余积，所述余积性质随即得出。
\end{proof}

\noindent
在下述引理中，我们考虑满足以下条件的偶对 $(M,\nabla)$：
\begin{enumerate}
\item
\label{crystalline-item-complete}
$M$ 是 $p$-进完备的 $D$-模；
\item
\label{crystalline-item-connection}
$\nabla:M\to M\otimes^\wedge_D\Omega_D$ 是联络，亦即
$\nabla(fm)=m\otimes\text df+f\nabla(m)$；
\item
\label{crystalline-item-integrable}
$\nabla$ 可积（见注 \ref{crystalline-remark-connection}）；
\item
\label{crystalline-item-topologically-quasi-nilpotent}
$\nabla$ {\it 拓扑拟幂零}：若对某些算子 $\theta_i:M\to M$ 写成
$\nabla(m)=\sum\theta_i(m)\text d x_i$，则对任意 $m\in M$，
只有有限多个偶对 $(i,k)$ 满足 $\theta_i^k(m)\notin pM$。
\end{enumerate}
文献中有时把算子 $\theta_i$ 记为 $\nabla_{\partial/\partial x_i}$。
下述引理构造从 $\text{Cris}(X/S)$ 上拟凝聚模中的晶体到这类偶对所成
范畴的函子。我们将在 命题 \ref{crystalline-proposition-crystals-on-affine}
中证明此函子是等价。

\begin{lemma}
\label{crystalline-lemma-crystals-on-affine}
在上述情形中，存在函子
$$
\begin{matrix}
\text{Cris}(X/S)\text{ 上拟凝聚} \\
\mathcal{O}_{X/S}\text{-模中的晶体}
\end{matrix}
\longrightarrow
\begin{matrix}
\text{满足 (\ref{crystalline-item-complete})、(\ref{crystalline-item-connection})、} \\
\text{(\ref{crystalline-item-integrable}) 及
(\ref{crystalline-item-topologically-quasi-nilpotent}) 的偶对 }(M,\nabla)
\end{matrix}
$$
\end{lemma}

\begin{proof}
设 $\mathcal F$ 是 $X/S$ 上拟凝聚模中的晶体。置
$T_e=\Spec(D_e)$，则当 $e\gg0$ 时，$(X,T_e,\bar\gamma)$ 是
$\text{Cris}(X/S)$ 的对象。我们有闭浸入
$$
(X, T_e, \bar\gamma) \to (X, T_{e + 1}, \bar\gamma) \to \ldots
$$
置
$$
M =
\lim_e \Gamma((X, T_e, \bar\gamma), \mathcal{F}) =
\lim_e \Gamma(T_e, \mathcal{F}_{T_e}) = \lim_e M_e
$$
由于 $\mathcal F$ 局部拟凝聚，有
$\mathcal F_{T_e}=\widetilde{M_e}$。由于 $\mathcal F$ 是晶体，有
$M_e=M_{e+1}/p^eM_{e+1}$。故 $M_e=M/p^eM$，且 $M$ 是 $p$-进完备的；
见 “代数” 引理 \ref{algebra-lemma-limit-complete}。

\medskip\noindent
由引理 \ref{crystalline-lemma-automatic-connection}，$\mathcal F$ 配备典范可积联络
$\nabla:\mathcal F\to\mathcal F\otimes\Omega_{X/S}$。
在上述对象 $T_e$ 上取此联络，得到典范可积联络
$$
\nabla : M \longrightarrow M \otimes^\wedge_D \Omega_D
$$
为看出它拓扑幂零，我们具体解释此条件。

\medskip\noindent
现在对各环 $D(n)$ 作同样处理，得到 $p$-进完备的 $D(n)$-模 $M(n)$。
再次利用 $\mathcal F$ 的晶体性质，得到同构
$$
M \otimes^\wedge_{D, p_0} D(1) \rightarrow M(1)
\leftarrow M \otimes^\wedge_{D, p_1} D(1)
$$
与 引理 \ref{crystalline-lemma-automatic-connection} 的证明比较。以 $c$
表示从左到右的复合。取 $m\in M$，并记
$\xi_i=x_i\otimes1-1\otimes x_i$。利用
(\ref{crystalline-equation-structure-Dn})，可唯一写成
$$
c(m \otimes 1) = \sum\nolimits_K \theta_K(m) \otimes \prod \xi_i^{[k_i]}
$$
其中 $\theta_K(m)\in M$，求和遍历满足 $k_i\geq0$ 且
$\sum k_i<\infty$ 的多重指标 $K=(k_i)$。当 $K$ 在第 $i$ 个位置为 $1$、
其余位置为零时，置 $\theta_i=\theta_K$。我们有
$$
\nabla(m) = \sum \theta_i(m) \text{d}x_i.
$$
这可由与 $\nabla$ 的定义比较看出。具体地，在 引理
\ref{crystalline-lemma-automatic-connection} 中定义等式为
$p_1^*m=\nabla(m)-c(p_0^*m)$；而符号仍相符，因为 Stacks Project
始终在模去对角理想平方后采用
$\text df=p_1(f)-p_0(f)$，所以
$\xi_i=x_i\otimes1-1\otimes x_i$ 模去对角理想平方后映到
$-\text d x_i$。

\medskip\noindent
以 $q_i:D\to D(2)$ 与 $q_{ij}:D(1)\to D(2)$ 表示对应于指标 $i,j$
的余积典范映射。与 引理 \ref{crystalline-lemma-automatic-connection}
证明的最后一段相同，可知
$$
q_{02}^*c = q_{12}^*c \circ q_{01}^*c.
$$
这意味着
$$
\sum\nolimits_{K''} \theta_{K''}(m) \otimes \prod {\zeta''_i}^{[k''_i]}
=
\sum\nolimits_{K', K} \theta_{K'}(\theta_K(m))
\otimes \prod {\zeta'_i}^{[k'_i]} \prod \zeta_i^{[k_i]}
$$
在 $M\otimes^\wedge_{D,q_2}D(2)$ 中成立，其中
\begin{align*}
\zeta_i & = x_i \otimes 1 \otimes 1 - 1 \otimes x_i \otimes 1,\\
\zeta'_i & = 1 \otimes x_i \otimes 1 - 1 \otimes 1 \otimes x_i,\\
\zeta''_i & = x_i \otimes 1 \otimes 1 - 1 \otimes 1 \otimes x_i.
\end{align*}
特别地，$\zeta''_i=\zeta_i+\zeta'_i$，并且由引理
\ref{crystalline-lemma-structure-Dn}，$D(2)$ 是 $q_2(D)$ 上以
$\zeta_i,\zeta'_i$ 为变量的除幂多项式环的 $p$-进完备化。
比较上式系数，立即得到
$\theta_i\circ\theta_j=\theta_j\circ\theta_i$
（这也给出 $\nabla$ 可积性的另一证明），以及
$$
\theta_K(m) = (\prod \theta_i^{k_i})(m).
$$
特别地，由于上述表示 $c(m\otimes1)$ 的和必须 $p$-进收敛，
可知对每个 $i$ 与每个 $m\in M$，只有有限多个 $\theta_i^k(m)$
模 $p$ 非零。
\end{proof}

\begin{proposition}
\label{crystalline-proposition-crystals-on-affine}
函子
$$
\begin{matrix}
\text{Cris}(X/S)\text{ 上拟凝聚} \\
\mathcal{O}_{X/S}\text{-模中的晶体}
\end{matrix}
\longrightarrow
\begin{matrix}
\text{满足 (\ref{crystalline-item-complete})、(\ref{crystalline-item-connection})、} \\
\text{(\ref{crystalline-item-integrable}) 及
(\ref{crystalline-item-topologically-quasi-nilpotent}) 的偶对 }(M,\nabla)
\end{matrix}
$$
是 引理 \ref{crystalline-lemma-crystals-on-affine} 中的范畴等价。
\end{proposition}

\begin{proof}
给定 $(M,\nabla)$。我们将构造拟凝聚模中的晶体 $\mathcal F$。
写成 $\nabla(m)=\sum\theta_i(m)\text d x_i$。则
$\theta_i\circ\theta_j=\theta_j\circ\theta_i$；对任意满足
$k_i\geq0$ 且 $\sum k_i<\infty$ 的多重指标 $K=(k_i)$，可置
$\theta_K(m)=(\prod\theta_i^{k_i})(m)$。

\medskip\noindent
设 $(U,T,\delta)$ 是 $\text{Cris}(X/S)$ 中 $T$ 仿射的任意对象。
写 $T=\Spec(B)$，并设 $U\to T$ 的理想为 $J_B\subset B$。
由引理 \ref{crystalline-lemma-set-generators}，存在整数 $e$ 及态射
$$
f : (U, T, \delta) \longrightarrow (X, T_e, \bar\gamma)
$$
其中 $T_e=\Spec(D_e)$，同引理 \ref{crystalline-lemma-crystals-on-affine} 的证明。
选取这样的 $e$ 与 $f$；相应的除幂 $A$-代数映射也记为 $f:D\to B$。
我们把 $\mathcal F_T$ 定义为与下列 $B$-模相伴的拟凝聚
$\mathcal O_T$-模层：
$$
M \otimes_{D, f} B.
$$
但必须证明它与 $f$ 的选取无关。设 $g:D\to B$ 是另一个这样的态射。
由于 $f$ 与 $g$ 是 $\text{Cris}(X/S)$ 中的态射，映射
$f-g:D\to B$ 的像包含在除幂理想 $J_B$ 中。记
$\xi_i=f(x_i)-g(x_i)\in J_B$。仿照引理
\ref{crystalline-lemma-crystals-on-affine} 的证明，由下式定义同构
$$
c_{f, g} : M \otimes_{D, f} B \longrightarrow M \otimes_{D, g} B
$$
$$
m \otimes 1 \longmapsto 
\sum\nolimits_K \theta_K(m) \otimes \prod \xi_i^{[k_i]}
$$
由上述讨论及 $\nabla$ 拓扑拟幂零（故该和有限！），此定义有意义。
计算表明
$$
c_{g, h} \circ c_{f, g} = c_{f, h}
$$
对任意第三个态射
$h:(U,T,\delta)\longrightarrow(X,T_e,\bar\gamma)$ 成立。
并且 $c_{f,f}=1$。因此这些映射都是同构，从而模 $\mathcal F_T$
与 $f$ 的选取无关。

\medskip\noindent
若 $a:(U',T',\delta')\to(U,T,\delta)$ 是 $\text{Cris}(X/S)$
中仿射对象之间的态射，则取 $f'=f\circ a$，显然存在典范同构
$a^*\mathcal F_T\to\mathcal F_{T'}$。略去此映射与 $f$ 的选取无关的验证。
以这些映射为限制映射，显然得到 $\text{Cris}(X/S)$ 中仿射对象所成
全子范畴上的一个拟凝聚模晶体。略去它延拓为整个
$\text{Cris}(X/S)$ 上晶体的证明。我们也略去以下证明：此过程为函子，
并且与 引理 \ref{crystalline-lemma-crystals-on-affine} 中构造的函子互为拟逆。
\end{proof}

\begin{lemma}
\label{crystalline-lemma-crystals-on-affine-smooth}
在 情形 \ref{crystalline-situation-affine} 中，设 $A\to P'\to C$ 是环映射，
$A\to P'$ 光滑，而 $P'\to C$ 满射且核为 $J'$。令 $D'$ 为
$D_{P',\gamma}(J')$ 的 $p$-进完备化。则存在如上的数据
$A\to P\to C$ 的一种选取，以及除幂 $A$-代数同态
$$
a : D \longrightarrow D',\quad b : D' \longrightarrow D
$$
它们与映射 $D\to C$ 及 $D'\to C$ 相容，并且
$a\circ b=\text{id}_{D'}$。这些映射诱导以下两类偶对所成范畴之间的等价：
在 $D$ 上满足
(\ref{crystalline-item-complete})、(\ref{crystalline-item-connection})、
(\ref{crystalline-item-integrable}) 与 (\ref{crystalline-item-topologically-quasi-nilpotent})
的偶对 $(M,\nabla)$，以及在 $D'$ 上满足
(\ref{crystalline-item-complete})、(\ref{crystalline-item-connection})、
(\ref{crystalline-item-integrable}) 与
(\ref{crystalline-item-topologically-quasi-nilpotent})\footnote{对 $D'$ 上的
$(M',\nabla')$ 表述此条件颇为棘手，见证明。} 的偶对 $(M',\nabla')$。
特别地，命题 \ref{crystalline-proposition-crystals-on-affine} 的范畴等价
对取值于 $D'$ 上偶对的相应函子也成立。
\end{lemma}

\begin{proof}
选取 $A$ 上多项式代数 $P=A[y_1,\ldots,y_m]$ 及满射 $P\to P'$，
并以复合 $P\to P'\to C$ 定义 $P\to C$。由除幂包络与完备化的函子性，
得到满射 $a:D\to D'$。取足够大的 $e$，使 $D_e$ 是 $C$ 在 $A$
上的除幂加厚。于是 $D_e\to C$ 是核局部幂零的满射；见
“除幂代数”引理 \ref{dpa-lemma-nil}。置
$D'_e=D'/p^eD'$，则 $D_e\to D'_e$ 的核局部幂零。因此由
“代数” 引理 \ref{algebra-lemma-smooth-strong-lift}，可找到映射
$P'\to D'_e$ 的提升 $\beta_e:P'\to D_e$。注意，对任意 $i\geq0$，
$D_{e+i+1}\to D_{e+i}\times_{D'_{e+i}}D'_{e+i+1}$
是核平方为零的满射，因为 $p^{e+i}D\to p^{e+i}D'$ 满射。
依次把通常的提升性质（“代数” 命题
\ref{algebra-proposition-smooth-formally-smooth}）应用于各图
$$
\xymatrix{
P' \ar[r] & D_{e + i} \times_{D'_{e + i}} D'_{e + i + 1} \\
A \ar[u] \ar[r] & D_{e + i + 1} \ar[u]
}
$$
可知存在 $A$-代数映射 $\beta:P'\to D$，它与 $a$ 的复合为给定映射
$P'\to D'$。由除幂包络的泛性质得到映射
$D_{P',\gamma}(J')\to D$。由于 $D$ 是 $p$-进完备的，得到
$b:D'\to D$，满足 $a\circ b=\text{id}_{D'}$。

\medskip\noindent
考虑基变换函子
$$
F : (M, \nabla) \longmapsto
(M \otimes^\wedge_{D, a} D', \nabla')
\quad\text{且}\quad
G : (M', \nabla') \longmapsto
(M' \otimes^\wedge_{D', b} D, \nabla)
$$
它们作用于满足 (\ref{crystalline-item-complete})、(\ref{crystalline-item-connection}) 与
(\ref{crystalline-item-integrable}) 的带联络模；见注
\ref{crystalline-remark-base-change-connection}。由于
$a\circ b=\text{id}_{D'}$，$F\circ G$ 是恒等函子。若
$G(M',\nabla')$ 满足性质
(\ref{crystalline-item-topologically-quasi-nilpotent})，我们便称
$(M',\nabla')$ 也满足该性质。形式论证表明，为完成证明，只需在
$(M,\nabla)$ 满足四个条件 (\ref{crystalline-item-complete})、
(\ref{crystalline-item-connection})、(\ref{crystalline-item-integrable}) 与
(\ref{crystalline-item-topologically-quasi-nilpotent}) 时，证明
$G(F(M,\nabla))$ 同构于 $(M,\nabla)$。为此使用 命题
\ref{crystalline-proposition-crystals-on-affine} 证明中的函子性同构
$$
c_{\text{id}_D, b \circ a} :
M \otimes_{D, \text{id}_D} D
\longrightarrow
M \otimes_{D, b \circ a} D
$$
（它要求 $\nabla$ 拓扑拟幂零，而这已在假设中）。
最后还需证明此映射是水平的，即与联络相容；此验证从略。

\medskip\noindent
证明中的最后一个断言随即得出。
\end{proof}

\begin{remark}
\label{crystalline-remark-equivalence-more-general}
命题 \ref{crystalline-proposition-crystals-on-affine} 的等价在下述情形仍成立：
从满射 $P\to C$ 出发，其中 $P/A$ 满足 “代数” 引理
\ref{algebra-lemma-smooth-strong-lift} 的强提升性质。其证明可仿照
引理 \ref{crystalline-lemma-crystals-on-affine-smooth} 的证明。
（若以后需要，将在此补充细节。）此结果大概也有直接证明；不过采用
多项式环的优点在于，各环 $D(n)$ 是除幂多项式环的 $p$-进完备化，
从而简化代数运算。
\end{remark}



\section{上同调的一般评注}
\label{crystalline-section-cohomology-lqc}

\noindent
本节作一些准备，把仿射概形晶体位点上模的上同调转化为代数问题。

\begin{lemma}
\label{crystalline-lemma-vanishing-lqc}
在 情形 \ref{crystalline-situation-global} 中，设 $\mathcal F$ 是
$\text{Cris}(X/S)$ 上局部拟凝聚的 $\mathcal O_{X/S}$-模。则
$$
H^p((U, T, \delta), \mathcal{F}) = 0
$$
对所有 $p>0$ 以及所有 $T$ 或 $U$ 仿射的 $(U,T,\delta)$ 成立。
\end{lemma}

\begin{proof}
由于 $U\to T$ 是加厚，$U$ 仿射当且仅当 $T$ 仿射；见“极限”
引理 \ref{limits-lemma-affine}。把“位点上同调”引理
\ref{sites-cohomology-lemma-cech-vanish-collection} 应用于仿射对象
$(U,T,\delta)$ 所成的族 $\mathcal B$，以及仿射开覆盖
$\mathcal U=\{(U_i,T_i,\delta_i)\to(U,T,\delta)\}$ 所成的族
$\text{Cov}$。这种覆盖的 {\v C}ech 复形
${\check C}^*(\mathcal U,\mathcal F)$，恰是拟凝聚
$\mathcal O_T$-模 $\mathcal F_T$ 关于仿射概形 $T$ 的仿射开覆盖
$\{T_i\to T\}$ 的 {\v C}ech 复形（这里使用了 $\mathcal F$
局部拟凝聚的假设）。因此，由“概形上同调”引理
\ref{coherent-lemma-cech-cohomology-quasi-coherent} 与
\ref{coherent-lemma-quasi-coherent-affine-cohomology-zero}，
其 {\v C}ech 上同调为零。于是“位点上同调”引理
\ref{sites-cohomology-lemma-cech-vanish-collection} 的假设均成立，结论得证。
\end{proof}

\begin{lemma}
\label{crystalline-lemma-compare}
在情形 \ref{crystalline-situation-global} 中，进一步假设 $X$ 与 $S$
都是仿射概形。考虑 $\text{Cris}(X/S)$ 中由除幂加厚
$(X,T,\delta)$ 所成的全子范畴 $\mathcal C$，并赋予混沌拓扑
（见 “位点论” 例 \ref{sites-example-indiscrete}）。对任意局部拟凝聚的
$\mathcal O_{X/S}$-模 $\mathcal F$，有
$$
R\Gamma(\mathcal{C}, \mathcal{F}|_\mathcal{C}) =
R\Gamma(\text{Cris}(X/S), \mathcal{F})
$$
\end{lemma}

\begin{proof}
以 $\text{AffineCris}(X/S)$ 表示 $\text{Cris}(X/S)$ 中由 $U$ 与 $T$
均仿射的对象 $(U,T,\delta)$ 所成的全子范畴。规定
$\text{AffineCris}(X/S)$ 的态射族
$\{(U_i,T_i,\delta_i)\to(U,T,\delta)\}_{i\in I}$ 是覆盖，当且仅当它在
$\text{Cris}(X/S)$ 中是覆盖，由此把该范畴变成位点。在此定义下，包含函子
$$
\text{AffineCris}(X/S) \longrightarrow \text{Cris}(X/S)
$$
是 “位点论” 定义 \ref{sites-definition-special-cocontinuous-functor}
意义下的特殊余连续函子。证明与 “拓扑” 引理
\ref{topologies-lemma-affine-big-site-Zariski} 的证明完全相同。
因此，沿上述包含函子限制后，$\text{Cris}(X/S)$ 上的层拓扑斯与
$\text{AffineCris}(X/S)$ 上的层拓扑斯相同。于是只需对包含
$\mathcal C\subset\text{AffineCris}(X/S)$ 证明相应断言。

\medskip\noindent
下文不再说明地使用 $\mathcal C$ 与 $\text{AffineCris}(X/S)$
具有积和纤维积这一事实（细节略；见引理
\ref{crystalline-lemma-divided-power-thickening-fibre-products}）。包含函子
$u:\mathcal C\to\text{AffineCris}(X/S)$ 全忠实、连续，并与积和纤维积交换。
我们断言它定义环化位点态射
$$
f :
(\text{AffineCris}(X/S), \mathcal{O}_{X/S})
\longrightarrow
(\Sh(\mathcal{C}), \mathcal{O}_{X/S}|_\mathcal{C})
$$
为此使用 “位点论” 引理 \ref{sites-lemma-directed-morphism}。由于
$\mathcal C$ 具有纤维积且 $u$ 与之交换，各范畴
$\mathcal I^u_{(U,T,\delta)}$ 是有向范畴的不交并（由 “位点论” 引理
\ref{sites-lemma-almost-directed} 与 “范畴” 引理
\ref{categories-lemma-split-into-directed}）。因此只需证明
$\mathcal I^u_{(U,T,\delta)}$ 连通。非空性由引理
\ref{crystalline-lemma-set-generators} 得出：由于 $U$ 与 $T$ 仿射，该引理说明
$\mathcal C$ 中至少有一个对象 $(X,T',\delta')$ 及一个除幂加厚态射
$(U,T,\delta)\to(X,T',\delta')$。连通性由 $\mathcal C$ 具有积且
$u$ 与积交换得出（与 “位点论” 引理 \ref{sites-lemma-directed}
的证明比较）。

\medskip\noindent
注意 $f_*\mathcal F=\mathcal F|_\mathcal C$。因此，若对 $p>0$ 有
$R^pf_*\mathcal F=0$，引理即得证；见“位点上同调”引理
\ref{sites-cohomology-lemma-apply-Leray}。由“位点上同调”引理
\ref{sites-cohomology-lemma-higher-direct-images}，只需对所有
$(X,T,\delta)$ 证明
$H^p(\text{AffineCris}(X/S)/(X,T,\delta),\mathcal F)=0$。
这由引理 \ref{crystalline-lemma-vanishing-lqc} 得出，因为位点
$\text{AffineCris}(X/S)/(X,T,\delta)$ 的拓扑斯等价于该引理中使用的
位点 $\text{Cris}(X/S)/(X,T,\delta)$ 的拓扑斯。
\end{proof}

\begin{lemma}
\label{crystalline-lemma-complete}
在 情形 \ref{crystalline-situation-affine} 中，置
$\mathcal C=(\text{Cris}(C/A))^{opp}$ 与
$\mathcal C^\wedge=(\text{Cris}^\wedge(C/A))^{opp}$，并赋予混沌拓扑；
记号见注 \ref{crystalline-remark-completed-affine-site}。存在拓扑斯态射
$$
g : \Sh(\mathcal{C}) \longrightarrow \Sh(\mathcal{C}^\wedge)
$$
使得若 $\mathcal F$ 是 $\mathcal C$ 上的交换群层，则
$$
R^pg_*\mathcal{F}(B \to C, \delta) =
\left\{
\begin{matrix}
\lim_e \mathcal{F}(B_e \to C, \delta) & \text{若 }p = 0 \\
R^1\lim_e \mathcal{F}(B_e \to C, \delta) & \text{若 }p = 1 \\
0 & \text{其他情形}
\end{matrix}
\right.
$$
其中当 $e\gg0$ 时，$B_e=B/p^eB$。
\end{lemma}

\begin{proof}
范畴之间的任意函子都在同方向定义混沌拓扑斯之间的态射；例如，可把这种
函子看作位点之间的余连续函子，见 “位点论” 节
\ref{sites-section-cocontinuous-morphism-topoi}。略去对 $g_*\mathcal F$
之描述的证明。注意，命题中只有当 $e$ 足够大时，
$(B_e\to C,\delta)$ 才是 $\text{Cris}(C/A)$ 的对象。
设 $\mathcal I$ 是 $\mathcal C$ 上的内射交换群层。则转移映射
$$
\mathcal{I}(B_e \to C, \delta) \leftarrow
\mathcal{I}(B_{e + 1} \to C, \delta)
$$
满射，因为态射
$$
(B_e \to C, \delta)
\longrightarrow
(B_{e + 1} \to C, \delta)
$$
在范畴 $\mathcal C$ 中是单态。因此，对内射交换群层，引理所示公式两边
相同。取 $\mathcal F$ 的一个内射消解即可得到结论（层是预层，故正合性
在各对象上截面群的层面检验）。
\end{proof}

\begin{lemma}
\label{crystalline-lemma-category-with-covering}
设 $\mathcal C$ 是赋予混沌拓扑的范畴。设 $X$ 是 $\mathcal C$ 的对象，
并且 $\mathcal C$ 的每个对象都有一个指向 $X$ 的态射。假设
$\mathcal C$ 具有二元积。则对 $\mathcal C$ 上每个交换群层
$\mathcal F$，总上同调 $R\Gamma(\mathcal C,\mathcal F)$ 由复形
$$
\mathcal{F}(X) \to \mathcal{F}(X \times X) \to
\mathcal{F}(X \times X \times X) \to \ldots
$$
表示；该复形与余单纯交换群 $[n]\mapsto\mathcal F(X^n)$ 相伴。
\end{lemma}

\begin{proof}
由于任意预层在 $\mathcal C$ 上都是层，对所有 $q>0$ 有
$H^q(X^p,\mathcal F)=0$。关于 $X$ 的假设是 $h_X\to*$ 满射。
利用 $H^q(X,\mathcal F)=H^q(h_X,\mathcal F)$ 与
$H^q(\mathcal C,\mathcal F)=H^q(*,\mathcal F)$，可见该断言是
“位点上同调”引理
\ref{sites-cohomology-lemma-cech-to-cohomology-sheaf-sets} 的特例。
\end{proof}









\section{余单纯准备}
\label{crystalline-section-cohomology}

\noindent
本节比较晶体上同调与德拉姆上同调。我们遵循 \cite{Bhatt}。

\begin{example}
\label{crystalline-example-cosimplicial-module}
设 $A_*$ 是任意余单纯环。考虑由下式定义的余单纯模 $M_*$：
$$
M_n = \bigoplus\nolimits_{i = 0, ..., n} A_n e_i
$$
对映射 $f:[n]\to[m]$，把 $M_*(f):M_n\to M_m$ 定义为将 $e_i$
映到 $e_{f(i)}$ 的唯一 $A_*(f)$-线性映射。我们断言 $M_*$
上的恒等映射同伦于 $0$。事实上，一个同伦由余单纯模映射
$$
h : M_* \longrightarrow \Hom(\Delta[1], M_*)
$$
给出；见第 \ref{crystalline-section-cosimplicial} 节。对
$j\in\{0,\ldots,n+1\}$，令 $\alpha^n_j:[n]\to[1]$ 为由
$\alpha^n_j(i)=0\Leftrightarrow i<j$ 定义的映射。于是
$\Delta[1]_n=\{\alpha^n_0,\ldots,\alpha^n_{n+1}\}$，相应地
$\Hom(\Delta[1],M_*)_n=\prod_{j=0,\ldots,n+1}M_n$；见“单纯方法”第
\ref{simplicial-section-homotopy} 节与第
\ref{simplicial-section-homotopy-cosimplicial} 节。我们不用这个积表示，
而把 $\Hom(\Delta[1],M_*)_n$ 的元素看作函数
$\Delta[1]_n\to M_n$。用此记号，把 $h$ 的次数 $n$ 分量定义为
$$
h_n(e_i)(\alpha^n_j) =
\left\{
\begin{matrix}
e_{i} & \text{若} & i < j \\
0 & \text{其他情形} 
\end{matrix}
\right.
$$
先验证 $h$ 是余单纯模态射。即对 $f:[n]\to[m]$，证明
\begin{equation}
\label{crystalline-equation-cosimplicial-morphism}
h_m \circ M_*(f) = \Hom(\Delta[1], M_*)(f) \circ h_n
\end{equation}
(\ref{crystalline-equation-cosimplicial-morphism}) 左端先在 $e_i$ 处取值，
再在 $\alpha^m_j$ 处取值，所得为
$$
h_m(e_{f(i)})(\alpha^m_j) =
\left\{
\begin{matrix}
e_{f(i)} & \text{若} & f(i) < j \\
0 & \text{其他情形}
\end{matrix}
\right.
$$
注意 $\alpha^m_j\circ f=\alpha^n_{j'}$，其中
$0\leq j'\leq n+1$ 是使 $f(i)<j$ 当且仅当 $i<j'$ 的唯一指标。
因此，(\ref{crystalline-equation-cosimplicial-morphism}) 右端先在 $e_i$ 处取值，
再在 $\alpha^m_j$ 处取值，所得为
$$
M_*(f)(h_n(e_i)(\alpha^m_j \circ f) =
M_*(f)(h_n(e_i)(\alpha^n_{j'})) =
\left\{
\begin{matrix}
e_{f(i)} & \text{若} & i < j' \\
0 & \text{其他情形} 
\end{matrix}
\right.
$$
由对 $j'$ 的描述可知两结果相等。因此 $h$ 是余单纯模映射。
令 $0:\Delta[0]\to\Delta[1]$ 与 $1:\Delta[0]\to\Delta[1]$
为显然映射，并以
$ev_0,ev_1:\Hom(\Delta[1],M_*)\to M_*$ 表示相应取值映射。
读者容易验证，复合
$$
ev_0 \circ h, ev_1 \circ h : M_* \longrightarrow M_*
$$
分别为 $1$ 与 $0$，故 $h$ 是所需的 $1$ 与 $0$ 之间的同伦。
\end{example}

\begin{lemma}
\label{crystalline-lemma-vanishing-omega-1}
采用 (\ref{crystalline-equation-omega-Dn}) 中的记号，复形
$$
\Omega_{D(0)} \to \Omega_{D(1)} \to \Omega_{D(2)} \to \ldots
$$
作为余单纯 $D(*)$-模同伦于零。
\end{lemma}

\begin{proof}
我们将使用 “单纯方法” 引理
\ref{simplicial-lemma-functorial-homotopy} 的原则，更具体地说，
使用 引理 \ref{crystalline-lemma-homotopy-tensor}：任意函子都把（余）单纯对象之间
的同伦映射变为同伦映射。引理中的复形等于余单纯模
$$
M_* = \left(
\Omega_{P/A} \to
\Omega_{P \otimes_A P/A} \to
\Omega_{P \otimes_A P \otimes_A P/A} \to \ldots
\right)
$$
沿余单纯环映射 $P\otimes_A\ldots\otimes_AP\to D(n)$ 基变换后再作
$p$-进完备化；这由引理
\ref{crystalline-lemma-module-differentials-divided-power-envelope} 得出，
见 (\ref{crystalline-equation-omega-D}) 后的评注。因此，只需证明余单纯模 $M_*$
同伦于零（这里使用基变换与 $p$-进完备化）。利用
$\mathbf Z\to A$ 的基变换，甚至可假设 $A=\mathbf Z$ 且
$P=\mathbf Z[\{x_i\}_{i\in I}]$。此时 $P^{\otimes n+1}$ 是以各元素
$$
x_i(e) = 1 \otimes \ldots \otimes x_i \otimes \ldots \otimes 1
$$
为变量的多项式代数，其中 $x_i$ 位于第 $e$ 个位置。复形的各模以
$\text d x_i(e)$ 为自由生成元。注意，若 $f:[n]\to[m]$，则
$$
M_*(f)(\text{d}x_i(e)) = \text{d}x_i(f(e))
$$
因此，$M_*$ 是按 $I$ 求和的、例
\ref{crystalline-example-cosimplicial-module} 中所研究模的副本之直和，结论得证。
\end{proof}

\begin{lemma}
\label{crystalline-lemma-vanishing}
采用 (\ref{crystalline-equation-Dn}) 与 (\ref{crystalline-equation-omega-Dn}) 中的记号，
给定任意余单纯 $D(*)$-模 $M_*$ 及 $i>0$，余单纯模
$$
M_0 \otimes^\wedge_{D(0)} \Omega^i_{D(0)} \to
M_1 \otimes^\wedge_{D(1)} \Omega^i_{D(1)} \to
M_2 \otimes^\wedge_{D(2)} \Omega^i_{D(2)} \to \ldots
$$
同伦于零，其中 $\Omega^i_{D(n)}$ 是 $\Omega_{D(n)}$
的第 $i$ 个外幂的 $p$-进完备化。
\end{lemma}

\begin{proof}
由引理 \ref{crystalline-lemma-vanishing-omega-1}，$\Omega_{D(*)}$ 的自同态
$0$ 与 $1$ 同伦。应用函子 $\wedge^i$，可知余单纯模
$\wedge^i\Omega_{D(*)}$ 也如此；见引理
\ref{crystalline-lemma-homotopy-tensor}。再次应用同一引理，得到其 $p$-进完备化
$\Omega^i_{D(*)}$ 同伦等价于零。与 $M_*$ 作张量积，得到
$M_*\otimes_{D(*)}\Omega^i_{D(*)}$ 同伦于零；再次见引理
\ref{crystalline-lemma-homotopy-tensor}。最后应用 $p$-进完备化函子即完成证明。
\end{proof}



\section{除幂 Poincar\'e 引理}
\label{crystalline-section-poincare}

\noindent
这里只给出最简单的版本。

\begin{lemma}
\label{crystalline-lemma-poincare}
设 $A$ 是环，$P=A\langle x_i\rangle$ 是 $A$ 上的除幂多项式环。
对任意 $A$-模 $M$，复形
$$
0 \to M \to
M \otimes_A P \to
M \otimes_A \Omega^1_{P/A, \delta} \to
M \otimes_A \Omega^2_{P/A, \delta} \to \ldots
$$
正合。令 $D$ 为 $P$ 的 $p$-进完备化，并令 $\Omega^i_D$ 为
$\Omega_{D/A,\delta}$ 第 $i$ 个外幂的 $p$-进完备化。对任意
$p$-进完备的 $A$-模 $M$，复形
$$
0 \to M \to
M \otimes^\wedge_A D \to
M \otimes^\wedge_A \Omega^1_D \to
M \otimes^\wedge_A \Omega^2_D \to \ldots
$$
正合。
\end{lemma}

\begin{proof}
只需证明 $A$-模复形
$$
E :
(0 \to A \to P \to \Omega^1_{P/A, \delta} \to
\Omega^2_{P/A, \delta} \to \ldots)
$$
同伦等价于零。对每个多重指标 $K=(k_i)$，考虑子复形 $E(K)$，
其次数 $j$ 分量为
$$
\bigoplus\nolimits_{I = \{i_1, \ldots, i_j\} \subset \text{Supp}(K)}
A
\prod\nolimits_{i \not \in I} x_i^{[k_i]}
\prod\nolimits_{i \in I} x_i^{[k_i - 1]}
\text{d}x_{i_1} \wedge \ldots \wedge \text{d}x_{i_j}
$$
由于 $E=\bigoplus E(K)$，只需证明每个复形 $E(K)$ 同伦于零。
若 $K=0$，则 $E(K):(A\to A)$ 同伦于零。若 $K$ 的支集
$S$ 非空且有限，则复形 $E(K)$ 同构于
$$
0 \to A \to
\bigoplus\nolimits_{s \in S} A \to
\wedge^2(\bigoplus\nolimits_{s \in S} A) \to
\ldots \to \wedge^{\# S}(\bigoplus\nolimits_{s \in S} A) \to 0
$$
后者同伦于零，例如由“代数进阶”引理
\ref{more-algebra-lemma-homotopy-koszul-abstract} 可知。
\end{proof}

\noindent
例 \ref{crystalline-example-integrate} 解释了下述引理的另一种（更直接的）方法。

\begin{lemma}
\label{crystalline-lemma-relative-poincare}
设 $A$ 是环，$(B,I,\delta)$ 是除幂环，其中 $B$ 是 $A$-代数。
令 $P=B\langle x_i\rangle$ 为 $B$ 上除幂多项式环，照例配备除幂理想
$J=IP+B\langle x_i\rangle_+$。设 $M$ 是配备可积联络
$\nabla:M\to M\otimes_B\Omega^1_{B/A,\delta}$ 的 $B$-模。
则德拉姆复形的映射
$$
M \otimes_B \Omega^*_{B/A, \delta}
\longrightarrow
M \otimes_P \Omega^*_{P/A, \delta}
$$
是拟同构。令 $D$、相应地 $D'$ 分别为 $B$、相应地 $P$ 的
$p$-进完备化，并令 $\Omega^i_D$、相应地 $\Omega^i_{D'}$ 分别为
$\Omega^i_{B/A,\delta}$、相应地 $\Omega^i_{P/A,\delta}$ 的
$p$-进完备化。设 $M$ 是配备整联络
$\nabla:M\to M\otimes^\wedge_D\Omega^1_D$ 的 $p$-进完备 $D$-模。
则德拉姆复形的映射
$$
M \otimes^\wedge_D \Omega^*_D
\longrightarrow
M \otimes^\wedge_D \Omega^*_{D'}
$$
是拟同构。
\end{lemma}

\begin{proof}
考虑 $\Omega^*_{B/A,\delta}$ 上由子复形
$F^i(\Omega^*_{B/A,\delta})=\sigma_{\geq i}\Omega^*_{B/A,\delta}$
给出的递减滤过 $F^*$；见 “同调” 节
\ref{homology-section-truncations}。由下式，它在
$\Omega^*_{P/A,\delta}$ 上诱导递减滤过 $F^*$：
$$
F^i(\Omega^*_{P/A, \delta}) =
F^i(\Omega^*_{B/A, \delta}) \wedge \Omega^*_{P/A, \delta}.
$$
我们有可分裂短正合列
$$
0 \to \Omega^1_{B/A, \delta} \otimes_B P \to
\Omega^1_{P/A, \delta} \to
\Omega^1_{P/B, \delta} \to 0
$$
其中最后一个模以 $\text d x_i$ 为自由基。由此可知
$F^i(\Omega^*_{P/A,\delta})\to\Omega^*_{P/A,\delta}$
逐项为可分裂单射，并且
$$
\text{gr}^i_F(\Omega^*_{P/A, \delta}) =
\Omega^i_{B/A, \delta} \otimes_B \Omega^*_{P/B, \delta}
$$
为复形的等式。因此可由下式在
$M\otimes_B\Omega^*_{B/A,\delta}$ 上定义滤过 $F^*$：
$$
F^i(M \otimes_B \Omega^*_{P/A, \delta}) =
M \otimes_B F^i(\Omega^*_{P/A, \delta})
$$
并且复形满足
$$
\text{gr}^i_F(M \otimes_B \Omega^*_{P/A, \delta}) =
M \otimes_B \Omega^i_{B/A, \delta} \otimes_B \Omega^*_{P/B, \delta}
$$
由引理 \ref{crystalline-lemma-poincare}，这些复形中的每一个都拟同构于置于次数
$0$ 的 $M\otimes_B\Omega^i_{B/A,\delta}$。因此，引理中第一个所示映射
是滤过复形的态射，并在分次片上诱导拟同构。这蕴含它自身是拟同构；
例如可使用与滤过复形相伴的谱序列，见 “同调” 节
\ref{homology-section-filtered-complex}。

\medskip\noindent
第二个拟同构的证明完全相同。
\end{proof}



\section{仿射情形的上同调}
\label{crystalline-section-cohomology-affine}

\noindent
回到第 \ref{crystalline-section-quasi-coherent-crystals} 节中研究的情形。
从 $(A,I,\gamma)$ 与 $A/I\to C$ 出发，置 $X=\Spec(C)$ 与
$S=\Spec(A)$。然后选取 $A$ 上多项式环 $P$ 及核为 $J$ 的满射
$P\to C$。由 (\ref{crystalline-equation-D}) 与 (\ref{crystalline-equation-Dn}) 得到
$D$ 和 $D(n)$。置 $T(n)_e=\Spec(D(n)/p^eD(n))$，使
$(X,T(n)_e,\delta(n))$ 是 $\text{Cris}(X/S)$ 的对象。
设 $\mathcal F$ 是 $\mathcal O_{X/S}$-模层，并对
$n=0,1,2,3,\ldots$ 置
$$
M(n) = \lim_e \Gamma((X, T(n)_e, \delta(n)), \mathcal{F})
$$
这构成余单纯环 $D(0),D(1),D(2),\ldots$ 上的余单纯模。

\begin{proposition}
\label{crystalline-proposition-compute-cohomology}
采用上述记号，假设：
\begin{enumerate}
\item $\mathcal F$ 局部拟凝聚；
\item 对 $\text{Cris}(X/S)$ 中任意态射
$(U,T,\delta)\to(U',T',\delta')$，若 $f:T\to T'$ 是闭浸入，
则映射 $c_f:f^*\mathcal F_{T'}\to\mathcal F_T$ 满射。
\end{enumerate}
则复形
$$
M(0) \to M(1) \to M(2) \to \ldots
$$
计算 $R\Gamma(\text{Cris}(X/S),\mathcal F)$。
\end{proposition}

\begin{proof}
利用假设 (1) 与 引理 \ref{crystalline-lemma-compare}，可知
$R\Gamma(\text{Cris}(X/S),\mathcal F)$ 同构于
$R\Gamma(\mathcal C,\mathcal F)$。注意，引理 \ref{crystalline-lemma-compare}
与 \ref{crystalline-lemma-complete} 中使用的范畴 $\mathcal C$ 相同。
设 $f:T\to T'$ 是 (2) 中的闭浸入。映射
$c_f:f^*\mathcal F_{T'}\to\mathcal F_T$ 满射，当且仅当
$\mathcal F_{T'}\to f_*\mathcal F_T$ 满射。因此，若 $\mathcal F$
满足 (1) 与 (2)，便得到 $T'$ 上拟凝聚 $\mathcal O_{T'}$-模的短正合列
$$
0 \to \mathcal{K} \to \mathcal{F}_{T'} \to f_*\mathcal{F}_T \to 0
$$
见“概形”第 \ref{schemes-section-quasi-coherent} 节，尤其是
引理 \ref{schemes-lemma-push-forward-quasi-coherent}。因此，若 $T'$ 仿射，
则由 $H^1(T',\mathcal K)$ 消失可知限制映射
$\mathcal F(U',T',\delta')\to\mathcal F(U,T,\delta)$ 满射；见
“概形上同调”引理
\ref{coherent-lemma-quasi-coherent-affine-cohomology-zero}。
于是 引理 \ref{crystalline-lemma-complete} 中各逆系统的转移映射满射。
由此得到对所有 $p\geq1$，
$R^pg_*(\mathcal F|_\mathcal C)=0$，其中 $g$ 如 引理
\ref{crystalline-lemma-complete}。范畴 $\mathcal C^\wedge$ 的对象 $D$ 满足 引理
\ref{crystalline-lemma-category-with-covering} 的假设。这由引理
\ref{crystalline-lemma-generator-completion} 及下式可知：
$$
D \times \ldots \times D = D(n)
$$
在 $\mathcal C$ 中成立，因为 $D(n)$ 是
$\text{Cris}^\wedge(C/A)$ 中 $D$ 的 $n+1$ 重余积；见
引理 \ref{crystalline-lemma-property-Dn}。结论得证。
\end{proof}

\begin{lemma}
\label{crystalline-lemma-cohomology-is-zero}
假设与记号同命题 \ref{crystalline-proposition-compute-cohomology}。则
$$
H^j(\text{Cris}(X/S), \mathcal{F} \otimes_{\mathcal{O}_{X/S}} \Omega^i_{X/S})
= 0
$$
对所有 $i>0$ 与所有 $j\geq0$ 成立。
\end{lemma}

\begin{proof}
由引理 \ref{crystalline-lemma-omega-locally-quasi-coherent}，
$\mathcal H=\mathcal F\otimes_{\mathcal O_{X/S}}\Omega^i_{X/S}$
也满足命题 \ref{crystalline-proposition-compute-cohomology} 的假设 (1) 与 (2)。记
$$
M(n)_e=\Gamma((X,T(n)_e,\delta(n)),\mathcal F)
$$
于是 $M(n)=\lim_eM(n)_e$。则
\begin{align*}
\lim_e \Gamma((X, T(n)_e, \delta(n)), \mathcal{H}) & =
\lim_e M(n)_e \otimes_{D(n)_e} \Omega_{D(n)}/p^e\Omega_{D(n)} \\
& = \lim_e M(n)_e \otimes_{D(n)} \Omega_{D(n)}
\end{align*}
由引理 \ref{crystalline-lemma-vanishing}，余单纯模
$$
M(0)_e \otimes_{D(0)} \Omega^i_{D(0)} \to
M(1)_e \otimes_{D(1)} \Omega^i_{D(1)} \to
M(2)_e \otimes_{D(2)} \Omega^i_{D(2)} \to \ldots
$$
同伦于零。由于转移映射 $M(n)_{e+1}\to M(n)_e$ 满射，
相伴复形的逆极限非周期\footnote{事实上，因为各同伦彼此拼合，
它们甚至同伦于零，不过这里不需要这一点。之所以采用这种迂回论证，
是因为在引理的假设下，我们不知道
$M(n)_e=M(n)_{e+1}/p^eM(n)_{e+1}$，所以极限
$\lim_eM(n)_e\otimes_{D(n)}\Omega^i_{D(n)}$ 并不是
$M(n)\otimes_{D(n)}\Omega^i_{D(n)}$ 的 $p$-进完备化。
若 $\mathcal F$ 是晶体，该等式则成立。}。
于是由命题 \ref{crystalline-proposition-compute-cohomology}，
$\mathcal H$ 的上同调消失。
\end{proof}

\begin{proposition}
\label{crystalline-proposition-compute-cohomology-crystal}
假设同命题 \ref{crystalline-proposition-compute-cohomology}，但现在假设
$\mathcal F$ 是拟凝聚模中的晶体。令 $(M,\nabla)$ 为相应的 $D$
上带联络模；见命题 \ref{crystalline-proposition-crystals-on-affine}。
则复形
$$
M \otimes^\wedge_D \Omega^*_D
$$
计算 $R\Gamma(\text{Cris}(X/S),\mathcal F)$。
\end{proposition}

\begin{proof}
我们使用与双复形 $K^{*,*}$ 相伴的两个谱序列来证明，其中
$$
K^{a, b} = M \otimes_D^\wedge \Omega^a_{D(b)}
$$
已有结果如下。引理 \ref{crystalline-lemma-vanishing} 表明每一列
$K^{a,*}$（$a>0$）非周期。命题
\ref{crystalline-proposition-compute-cohomology} 表明第一列 $K^{0,*}$ 拟同构于
$R\Gamma(\text{Cris}(X/S),\mathcal F)$。因此，与双复形相伴的第一个
谱序列表明，$R\Gamma(\text{Cris}(X/S),\mathcal F)$ 与
$\text{Tot}(K^{*,*})$ 之间存在典范拟同构。

\medskip\noindent
接着考虑各行 $K^{*,b}$。由引理 \ref{crystalline-lemma-structure-Dn}，
$b+1$ 个映射 $D\to D(b)$ 中每一个都把 $D(b)$ 表示为 $D$
上某个除幂多项式代数的 $p$-进完备化。因此引理
\ref{crystalline-lemma-relative-poincare} 表明映射
$$
M \otimes^\wedge_D\Omega^*_D
\longrightarrow
M \otimes^\wedge_{D(b)} \Omega^*_{D(b)} = K^{*, b}
$$
是拟同构。注意，这些映射在上同调上定义{\it 同一个}映射（甚至在导出范畴
中定义同一个映射），因为其逆由余对角映射 $D(b)\to D$ 给出，
它对应于乘法映射 $P\otimes_A\ldots\otimes_AP\to P$。
因此，第二个谱序列的 $E_1$ 页为
$$
E_1^{a, b} = H^a(M \otimes^\wedge_D\Omega^*_D)
$$
其微分为
$$
E_1^{a, 0} \xrightarrow{0}
E_1^{a, 1} \xrightarrow{1}
E_1^{a, 2} \xrightarrow{0}
E_1^{a, 3} \xrightarrow{1} \ldots
$$
因为每个微分都是给定等同
$H^a(M\otimes^\wedge_D\Omega^*_D)=E_1^{a,0}=E_1^{a,1}=\ldots$
的交错和。因此，$E_2$ 页在第一行等于
$H^a(M\otimes^\wedge_D\Omega^*_D)$，在其他各处为零。
由此，$M\otimes^\wedge_D\Omega^*_D$ 与第一行的等同诱导它与
$\text{Tot}(K^{*,*})$ 之间的拟同构。
\end{proof}

\begin{lemma}
\label{crystalline-lemma-compute-cohomology-crystal-smooth}
假设同命题 \ref{crystalline-proposition-compute-cohomology-crystal}。
设 $A \to P' \to C$ 为环映射，其中 $A \to P'$ 光滑，而 $P' \to C$
为核等于 $J'$ 的满射。令 $D'$ 为 $D_{P', \gamma}(J')$ 的 $p$-进完备化。
设 $(M', \nabla')$ 为与 $\mathcal{F}$ 对应的 $D'$ 上偶对，见
引理 \ref{crystalline-lemma-crystals-on-affine-smooth}。则复形
$$
M' \otimes^\wedge_{D'} \Omega^*_{D'}
$$
计算 $R\Gamma(\text{Cris}(X/S), \mathcal{F})$。
\end{lemma}

\begin{proof}
依引理 \ref{crystalline-lemma-crystals-on-affine-smooth} 选取
$a : D \to D'$ 与 $b : D' \to D$。
注意，基变换 $M = M' \otimes_{D', b} D$ 连同其联络 $\nabla$
对应于 $\mathcal{F}$。因此由
命题 \ref{crystalline-proposition-compute-cohomology-crystal}，
$M \otimes^\wedge_D \Omega_D^*$ 计算 $\mathcal{F}$ 的晶体上同调。
故只需证明由 $a$ 与 $b$ 诱导的基变换映射
$$
M' \otimes^\wedge_{D'} \Omega^*_{D'}
\longrightarrow
M \otimes^\wedge_D \Omega^*_D
\quad\text{且}\quad
M \otimes^\wedge_D \Omega^*_D
\longrightarrow
M' \otimes^\wedge_{D'} \Omega^*_{D'}
$$
均为拟同构。由于 $a \circ b = \text{id}_{D'}$，其中一个方向的复合
是复形 $M' \otimes^\wedge_{D'} \Omega^*_{D'}$ 上的恒等映射。
因此只需证明映射
$$
M \otimes^\wedge_D \Omega^*_D
\longrightarrow
M \otimes^\wedge_D \Omega^*_D
$$
由 $b \circ a : D \to D$ 诱导时为拟同构。（注意，两端确为同一复形，
因为 $M = M' \otimes^\wedge_{D', b} D$，从而
$M \otimes^\wedge_{D, b \circ a} D =
M' \otimes^\wedge_{D', b \circ a \circ b} D =
M' \otimes^\wedge_{D', b} D = M$。）事实上我们断言：对任意与到 $C$
的增广相容的除幂 $A$-代数同态 $\rho : D \to D$，其诱导映射
$M \otimes^\wedge_D \Omega^*_D \to M \otimes^\wedge_{D, \rho} \Omega^*_D$
均为拟同构。

\medskip\noindent
写成 $\rho(x_i) = x_i + z_i$。由于 $\rho$ 与到 $C$ 的增广相容，
各元素 $z_i$ 均属于 $D$ 的除幂理想。因此可把映射 $\rho$ 分解为复合
$$
D \xrightarrow{\sigma} D\langle \xi_i \rangle^\wedge \xrightarrow{\tau} D
$$
其中第一条映射由 $x_i \mapsto x_i + \xi_i$ 给出，第二条映射是把
$\xi_i$ 映到 $z_i$ 的除幂 $D$-代数映射。（这里用到了多项式代数、
除幂多项式代数、除幂包络以及 $p$-进完备化的泛性质。）注意，
$D\langle \xi_i \rangle^\wedge$ 存在一个{\it 自同构} $\alpha$，满足
$\alpha(x_i) = x_i - \xi_i$ 且 $\alpha(\xi_i) = \xi_i$。
把 引理 \ref{crystalline-lemma-relative-poincare} 应用于 $\alpha \circ \sigma$
（它把 $x_i$ 映到 $x_i$），并利用 $\alpha$ 为同构，可知 $\sigma$
诱导 $M \otimes^\wedge_D \Omega^*_D$ 与下述复形之间的拟同构：
$M \otimes^\wedge_{D, \sigma} \Omega^*_{D\langle x_i \rangle^\wedge}$.
另一方面，映射 $\tau$ 以 $D \to D\langle x_i \rangle^\wedge$、
$x_i \mapsto x_i$ 为左逆。再次使用 引理 \ref{crystalline-lemma-relative-poincare}，
可知 $\tau$ 诱导下述两复形之间的拟同构：
$M \otimes^\wedge_{D, \sigma} \Omega^*_{D\langle x_i \rangle^\wedge}$
与 $M \otimes^\wedge_{D, \tau \circ \sigma} \Omega^*_D$。复合这两个
拟同构，即得 $\rho$ 所诱导的拟同构
$M \otimes^\wedge_D \Omega^*_D \to M \otimes^\wedge_{D, \rho} \Omega^*_D$，
如所要求。
\end{proof}










\section{两个反例}
\label{crystalline-section-examples}

\noindent
在讨论晶体上同调的一些成功应用之前，先给出两个例子，以说明当所论
概形在基上不固有或为奇异概形时，晶体上同调为何表现不佳。第一个例子
见 \cite{BO}。

\begin{example}
\label{crystalline-example-torsion}
令 $A = \mathbf{Z}_p$，并在除幂理想 $(p)$ 上赋予其唯一的除幂结构
$\gamma$。令 $C = \mathbf{F}_p[x, y]/(x^2, xy, y^2)$。选取表示
$$
C = P/J = \mathbf{Z}_p[x, y]/(x^2, xy, y^2, p)
$$
令 $D = D_{P, \gamma}(J)^\wedge$，其除幂理想
$(\bar J, \bar \gamma)$ 如第 \ref{crystalline-section-quasi-coherent-crystals} 节。
仍以 $x,y$ 表示它们在 $D$ 中的像。考虑元素
$$
\tau = \bar\gamma_p(x^2)\bar\gamma_p(y^2) - \bar\gamma_p(xy)^2 \in D
$$
注意 $p\tau=0$，因为在 $D$ 中
$$
p! \bar\gamma_p(x^2) \bar\gamma_p(y^2) =
x^{2p} \bar\gamma_p(y^2) = \bar\gamma_p(x^2y^2) =
x^py^p \bar\gamma_p(xy) = p! \bar\gamma_p(xy)^2
$$
。又注意到 $\Omega_D$ 中 $\text{d}\tau=0$，因为
\begin{align*}
\text{d}(\bar\gamma_p(x^2) \bar\gamma_p(y^2))
& =
\bar\gamma_{p - 1}(x^2)\bar\gamma_p(y^2)\text{d}x^2 +
\bar\gamma_p(x^2)\bar\gamma_{p - 1}(y^2)\text{d}y^2 \\
& =
2 x \bar\gamma_{p - 1}(x^2)\bar\gamma_p(y^2)\text{d}x +
2 y \bar\gamma_p(x^2)\bar\gamma_{p - 1}(y^2)\text{d}y \\
& =
2/(p - 1)!( x^{2p - 1} \bar\gamma_p(y^2)\text{d}x +
y^{2p - 1} \bar\gamma_p(x^2)\text{d}y ) \\
& =
2/(p - 1)!
(x^{p - 1} \bar\gamma_p(xy^2)\text{d}x +
y^{p - 1} \bar\gamma_p(x^2y)\text{d}y) \\
& =
2/(p - 1)!
(x^{p - 1}y^p \bar\gamma_p(xy)\text{d}x +
x^py^{p - 1} \bar\gamma_p(xy)\text{d}y) \\
& =
2 \bar\gamma_{p - 1}(xy) \bar\gamma_p(xy)(y\text{d}x + x \text{d}y) \\
& = 
\text{d}(\bar\gamma_p(xy)^2)
\end{align*}
最后断言 $D$ 中 $\tau\neq0$。为此，只需构造 $\text{Cris}(C/S)$ 的一个
对象 $(B \to \mathbf{F}_p[x, y]/(x^2, xy, y^2), \delta)$，使得 $\tau$
在 $B$ 中的像非零。取
$$
B = \mathbf{F}_p[x, y, u, v]/(x^3, x^2y, xy^2, y^3, xu, yu, xv, yv, u^2, v^2)
$$
并取其到 $C$ 的显然满射。令 $K=\Ker(B\to C)$，考虑映射
$$
\delta_p : K \longrightarrow K,\quad
ax^2 + bxy + cy^2 + du + ev + fuv \longmapsto a^pu + c^pv
$$
可检验它满足“除幂代数”引理
\ref{dpa-lemma-need-only-gamma-p} 的假设 (1)、(2)、(3)，故定义一个除幂结构。
此外，$\tau$ 映到 $uv$，而后者在 $B$ 中非零。置
$X=\Spec(C)$ 与 $S=\Spec(A)$。由此得到以下结论：
\begin{enumerate}
\item $H^0(\text{Cris}(X/S), \mathcal{O}_{X/S})$ 含有 $p$-挠；
\item Frobenius 拉回
$F^* : H^0(\text{Cris}(X/S), \mathcal{O}_{X/S})
\to H^0(\text{Cris}(X/S), \mathcal{O}_{X/S})$ 不是单射。
\end{enumerate}
具体地，由命题 \ref{crystalline-proposition-compute-cohomology-crystal}，
$\tau$ 定义 $H^0(\text{Cris}(X/S), \mathcal{O}_{X/S})$ 的一个非零挠元。
同样，$F^*(\tau)=\sigma(\tau)$，其中 $\sigma:D\to D$ 由 $P$ 上任一
Frobenius 提升诱导。若取 $\sigma(x)=x^p$ 与 $\sigma(y)=y^p$，则直接
计算可得 $F^*(\tau)=0$。
\end{example}

\noindent
下一个例子表明，即使对仿射 $n$-空间，晶体上同调也不给出所期望的结果。

\begin{example}
\label{crystalline-example-affine-n-space}
令 $A=\mathbf Z_p$，并在除幂理想 $(p)$ 上赋予其唯一的除幂结构
$\gamma$。令 $C=\mathbf F_p[x_1,\ldots,x_r]$。选取表示
$$
C = P/J = P/pP\quad\text{其中}\quad P = \mathbf{Z}_p[x_1, \ldots, x_r]
$$
由 “除幂代数” 引理 \ref{dpa-lemma-gamma-extends}，
$pP$ 带有除幂。故置 $D=P^\wedge$ 并取除幂理想 $(p)$，便得到
第 \ref{crystalline-section-quasi-coherent-crystals} 节中的情形。于是
$R\Gamma(\text{Cris}(X/S), \mathcal{O}_{X/S})$ 由复形
$$
D \to \Omega^1_D \to \Omega^2_D \to \ldots \to \Omega^r_D
$$
表示，见命题 \ref{crystalline-proposition-compute-cohomology-crystal}。
假设 $r>0$，可得：
\begin{enumerate}
\item $X=\mathbf A^r_{\mathbf F_p}$ 在 $S=\Spec(\mathbf Z_p)$ 上的晶体
结构层之晶体上同调仅在次数 $0,\ldots,r$ 可能非零。
\item $H^0(\text{Cris}(X/S), \mathcal{O}_{X/S})=\mathbf Z_p$。
\item 上同调群 $H^r(\text{Cris}(X/S), \mathcal{O}_{X/S})$ 是无限群，
且不是挠阿贝尔群。
\item 上同调群 $H^r(\text{Cris}(X/S), \mathcal{O}_{X/S})$
对 $p$-进拓扑不分离。
\end{enumerate}
前两条尚属合理，第 (3)、(4) 条却令人不安！只要具体写出上面展示的
复形，便立即可见这些断言。这里仅考察 $r=1$。此时所考察的是下述
$p$-进完备模的两项复形：
$$
\text{d} :
D = \left(
\bigoplus\nolimits_{n \geq 0} \mathbf{Z}_p x^n
\right)^\wedge
\longrightarrow
\Omega^1_D = \left(
\bigoplus\nolimits_{n \geq 1} \mathbf{Z}_p x^{n - 1}\text{d}x
\right)^\wedge
$$
除右端缺少第一个直和项外，该映射由
$\text{diag}(0,1,2,3,4,\ldots)$ 给出。显然，
$\bigoplus_{n>0}\mathbf Z_p/n\mathbf Z_p$ 是余核的一个子群，故余核
为无限群。事实上，元素
$$
\omega = \sum\nolimits_{e > 0} p^e x^{p^{2e} - 1}\text{d}x
$$
显然不是余核中的挠元。然而情况还更糟。考虑元素
$$
\eta = \sum\nolimits_{e > 0} p^e x^{p^e - 1}\text{d}x
$$
对每个 $t>0$，元素 $\eta$ 模 $\text d$ 的像同余于
$\sum_{e>t}p^ex^{p^e-1}\text d x$，而后者可被 $p^t$ 整除。因此
$\eta$ 在 $H^1(\text{Cris}(X/S),\mathcal O_{X/S})$ 中的上同调类属于
$\bigcap p^tH^1(\text{Cris}(X/S),\mathcal O_{X/S})$。但 $\eta$ 不在
$\text d$ 的像中，因为否则它必须是某个
$a+\sum_{e>0}x^{p^e}$（其中 $a\in\mathbf Z_p$）的像，而这个元素不属于
左端。因此 $\eta$ 的上同调类非零，第 (4) 条得证。（事实上，$\eta$
的上同调类不是挠元。）好的一面是，作为一个 $p$-进完备模复形的
上同调，各群 $H^i(\text{Cris}(X/S),\mathcal O_{X/S})$ 都是导出完备的
$\mathbf Z_p$-模；见“代数进阶”第
\ref{more-algebra-section-derived-completion} 节。
\end{example}








\section{应用}
\label{crystalline-section-applications}

\noindent
本节汇集前述各节内容的一些应用。

\begin{proposition}
\label{crystalline-proposition-compare-with-de-Rham}
处于 情形 \ref{crystalline-situation-global}。设 $\mathcal F$ 为
$\text{Cris}(X/S)$ 上拟凝聚模中的晶体。复形的截断映射
$$
(\mathcal{F} \to
\mathcal{F} \otimes_{\mathcal{O}_{X/S}} \Omega^1_{X/S} \to
\mathcal{F} \otimes_{\mathcal{O}_{X/S}} \Omega^2_{X/S} \to \ldots)
\longrightarrow \mathcal{F}[0],
$$
虽非拟同构，但施加 $Ru_{X/S,*}$ 后成为拟同构。事实上，对任意 $i>0$，
有
$$
Ru_{X/S, *}(\mathcal{F} \otimes_{\mathcal{O}_{X/S}} \Omega^i_{X/S}) = 0.
$$
\end{proposition}

\begin{proof}
由引理 \ref{crystalline-lemma-automatic-connection} 得到命题中所示的德拉姆复形。
简记 $\mathcal H=\mathcal F\otimes\Omega^i_{X/S}$。设 $X'\subset X$
为仿射开子概形，且映入某个仿射开子概形 $S'\subset S$。则
$$
(Ru_{X/S, *}\mathcal{H})|_{X'_{Zar}} =
Ru_{X'/S', *}(\mathcal{H}|_{\text{Cris}(X'/S')}),
$$
见引理 \ref{crystalline-lemma-localize}。因此引理 \ref{crystalline-lemma-cohomology-is-zero}
表明，$Ru_{X/S,*}\mathcal H$ 是 $X_{Zar}$ 上的层复形，且其在任意仿射
开集上的上同调均为零。由于 $X$ 的拓扑有一个由仿射开集组成的基，
这说明 $Ru_{X/S,*}\mathcal H$ 拟同构于零。
\end{proof}

\begin{remark}
\label{crystalline-remark-vanishing}
命题 \ref{crystalline-proposition-compare-with-de-Rham} 的证明表明，结论
$$
Ru_{X/S, *}(\mathcal{F} \otimes_{\mathcal{O}_{X/S}} \Omega^i_{X/S}) = 0
$$
对任意满足 命题 \ref{crystalline-proposition-compute-cohomology} 中条件 (1)、(2)
的 $\mathcal O_{X/S}$-模 $\mathcal F$，在 $i>0$ 时均成立。这适用于以下
非晶体：所有 $\Omega^i_{X/S}$，以及任意形如
$\underline{\mathcal F}$ 的层，其中 $\mathcal F$ 是拟凝聚
$\mathcal O_X$-模。特别地，它适用于层
$\underline{\mathcal O_X}=\underline{\mathbf G_a}$。但要构造德拉姆复形，
需要类似 引理 \ref{crystalline-lemma-automatic-connection} 的结果，而这要求
$\mathcal F$ 为晶体。因此，（目前）满足 命题
\ref{crystalline-proposition-compare-with-de-Rham} 完整断言的模层，恰组成拟凝聚模中
晶体的范畴。
\end{remark}

\noindent
处于 情形 \ref{crystalline-situation-global}。设 $\mathcal F$ 为
$\text{Cris}(X/S)$ 上拟凝聚模中的晶体，并设 $(U,T,\delta)$ 为
$\text{Cris}(X/S)$ 的对象。命题
\ref{crystalline-proposition-compare-with-de-Rham} 使我们能够构造典范映射
\begin{equation}
\label{crystalline-equation-restriction}
R\Gamma(\text{Cris}(X/S), \mathcal{F})
\longrightarrow
R\Gamma(T, \mathcal{F}_T \otimes_{\mathcal{O}_T} \Omega^*_{T/S, \delta})
\end{equation}
确切地说，有 $R\Gamma(\text{Cris}(X/S),\mathcal F)=
R\Gamma(\text{Cris}(X/S),\mathcal F\otimes\Omega^*_{X/S})$；可把整体
上同调类限制到 $T$，且由引理 \ref{crystalline-lemma-module-of-differentials}，
$\Omega_{X/S}$ 限制为 $\Omega_{T/S,\delta}$。

















\section{若干进一步结果}
\label{crystalline-section-missing}

\noindent
本节列出一些尚缺证明的结果。我们先把它们表述为一系列备注；只有在补入
详细证明后，才会将其改写为正式的引理和命题。

\begin{remark}[高阶正像]
\label{crystalline-remark-compute-direct-image}
设 $p$ 为素数，并设
$(S,\mathcal I,\gamma)\to(S',\mathcal I',\gamma')$ 为
$\mathbf Z_{(p)}$ 上除幂概形的态射。设
$$
\xymatrix{
X \ar[r]_f \ar[d] & X' \ar[d] \\
S_0 \ar[r] & S'_0
}
$$
为概形态射的交换图，并假设 $p$ 在 $X$ 与 $X'$ 上局部幂零。设
$\mathcal F$ 为 $\text{Cris}(X/S)$ 上的 $\mathcal O_{X/S}$-模。则
$Rf_{\text{cris},*}\mathcal F$ 可如下计算。

\medskip\noindent
给定 $\text{Cris}(X'/S')$ 的对象 $(U',T',\delta')$，置
$U=X\times_{X'}U'=f^{-1}(U')$（它是 $X$ 的开子概形）。以
$(T_0,T,\delta)$ 表示 $S$ 上的除幂概形，使得
$$
\xymatrix{
T \ar[r] \ar[d] & T' \ar[d] \\
S \ar[r] & S'
}
$$
在除幂概形范畴中为笛卡尔图，见引理 \ref{crystalline-lemma-fibre-product}。
存在诱导态射 $U\to T_0$，从而得到态射
$(U/T)_{\text{cris}}\to(X/S)_{\text{cris}}$，见注
\ref{crystalline-remark-functoriality-cris}。令 $\mathcal F_U$ 为 $\mathcal F$ 的拉回，
并令 $\tau_{U/T}:(U/T)_{\text{cris}}\to T_{Zar}$ 为结构态射。则有
\begin{equation}
\label{crystalline-equation-identify-pushforward}
\left(Rf_{\text{cris}, *}\mathcal{F}\right)_{T'} =
R(T \to T')_*\left(R\tau_{U/T, *} \mathcal{F}_U \right)
\end{equation}
其中左端是限制（见第 \ref{crystalline-section-sheaves} 节）。

\medskip\noindent
提示：首先证明 $\text{Cris}(U/T)$ 是 $\text{Cris}(X/S)$ 在集合层
$f_{\text{cris}}^{-1}h_{(U',T',\delta')}$ 处的局部化（取 “位点论” 引理
\ref{sites-lemma-localize-topos-site} 的意义）。其次，利用内射
$\mathcal O_{X/S}$-模的拉回是 $\text{Cris}(U/T)$ 上的内射
$\mathcal O_{U/T}$-模，把断言化归到 $\mathcal F$ 为内射模以及模的正像
这一情形。最后检验结果对普通正像成立。
\end{remark}

\begin{remark}[Mayer-Vietoris]
\label{crystalline-remark-mayer-vietoris}
在 注 \ref{crystalline-remark-compute-direct-image} 的情形中，设有开覆盖
$X=X'\cup X''$，并记 $X'''=X'\cap X''$。令 $f'$、$f''$ 与 $f''$
分别为 $f$ 在 $X'$、$X''$ 与 $X'''$ 上的限制。再令
$\mathcal F'$、$\mathcal F''$ 与 $\mathcal F'''$ 分别为 $\mathcal F$
在 $X'$、$X''$ 与 $X'''$ 的晶体位点上的限制。则存在可区别三角
$$
Rf_{\text{cris}, *}\mathcal{F}
\longrightarrow
Rf'_{\text{cris}, *}\mathcal{F}' \oplus Rf''_{\text{cris}, *}\mathcal{F}''
\longrightarrow
Rf'''_{\text{cris}, *}\mathcal{F}'''
\longrightarrow
Rf_{\text{cris}, *}\mathcal{F}[1]
$$
于 $D(\mathcal O_{X'/S'})$ 中。

\medskip\noindent
提示：这是下述事实的形式推论：子范畴 $\text{Cris}(X'/S)$、
$\text{Cris}(X''/S)$、$\text{Cris}(X'''/S)$ 对应于
$\text{Cris}(X/S)$ 的终层之开子对象，且第三个是前两个的交。
\end{remark}

\begin{remark}[{\v C}ech 复形]
\label{crystalline-remark-cech-complex}
设 $p$ 为素数，并设 $(A,I,\gamma)$ 为除幂环，其中 $A$ 是
$\mathbf Z_{(p)}$-代数。置 $S=\Spec(A)$ 与 $S_0=\Spec(A/I)$。设 $X$
为 $S_0$ 上的分离概形\footnote{此假设并非严格必要；利用超覆盖可将
本备注的构造推广到一般情形。}，且 $p$ 在 $X$ 上局部幂零。设
$\mathcal F$ 为拟凝聚 $\mathcal O_{X/S}$-模中的晶体。

\medskip\noindent
选取 $X$ 的仿射开覆盖
$X=\bigcup_{\lambda\in\Lambda}U_\lambda$。写成
$U_\lambda=\Spec(C_\lambda)$。选取 $A$ 上的多项式代数 $P_\lambda$
及满射 $P_\lambda\to C_\lambda$。固定这些选择后，可构造一个计算
$R\Gamma(\text{Cris}(X/S),\mathcal F)$ 的 {\v C}ech 复形。

\medskip\noindent
给定 $n\geq0$ 及 $\lambda_0,\ldots,\lambda_n\in\Lambda$，记
$U_{\lambda_0\ldots\lambda_n}=U_{\lambda_0}\cap\ldots\cap U_{\lambda_n}$。
由假设它是仿射概形。写成
$U_{\lambda_0\ldots\lambda_n}=\Spec(C_{\lambda_0\ldots\lambda_n})$，并置
$$
P_{\lambda_0 \ldots \lambda_n} =
P_{\lambda_0} \otimes_A \ldots \otimes_A P_{\lambda_n}
$$
它带有到 $C_{\lambda_0\ldots\lambda_n}$ 的典范满射。以
$J_{\lambda_0\ldots\lambda_n}$ 表示其核，并令
$D_{\lambda_0\ldots\lambda_n}$ 为
$P_{\lambda_0\ldots\lambda_n}$ 中 $J_{\lambda_0\ldots\lambda_n}$
相对于 $\gamma$ 的除幂包络之 $p$-进完备化。令
$M_{\lambda_0\ldots\lambda_n}$ 为通过 命题
\ref{crystalline-proposition-crystals-on-affine} 与 $\mathcal F$ 在
$\text{Cris}(U_{\lambda_0\ldots\lambda_n}/S)$ 上的限制相对应的
$P_{\lambda_0\ldots\lambda_n}$-模。由构造得到余单纯除幂环 $D(*)$，
其次数 $n$ 部分为环
$$
D(n) =
\prod\nolimits_{\lambda_0 \ldots \lambda_n}
D_{\lambda_0 \ldots \lambda_n}
$$
（这里使用除幂包络的函子性，以及在以类似方式定义的环 $P(*)$ 上的
平凡余单纯结构）。由于 $M_{\lambda_0\ldots\lambda_n}$ 是 $\mathcal F$
在对象 $\Spec(D_{\lambda_0\ldots\lambda_n})$ 上的“值”，故由规则
$$
M(n) = \prod\nolimits_{\lambda_0 \ldots \lambda_n}
M_{\lambda_0 \ldots \lambda_n}
$$
定义的 $M(*)$ 构成余单纯 $D(*)$-模。现断言
$$
R\Gamma(\text{Cris}(X/S), \mathcal{F}) = s(M(*))
$$
其中 $s(-)$ 表示与余单纯模相伴的上链复形（见 “单纯方法” 节
\ref{simplicial-section-dold-kan-cosimplicial}）。

\medskip\noindent
提示：其证明类似于 命题 \ref{crystalline-proposition-compute-cohomology} 的
证明（特别地，该结果对满足该命题假设的任意模成立）。
\end{remark}

\begin{remark}[交错 {\v C}ech 复形]
\label{crystalline-remark-alternating-cech-complex}
设 $p$ 为素数，并设 $(A,I,\gamma)$ 为除幂环，其中 $A$ 是
$\mathbf Z_{(p)}$-代数。置 $S=\Spec(A)$ 与 $S_0=\Spec(A/I)$。设 $X$
为 $S_0$ 上的分离拟紧概形，且 $p$ 在 $X$ 上局部幂零。设
$\mathcal F$ 为拟凝聚 $\mathcal O_{X/S}$-模中的晶体。

\medskip\noindent
选取 $X$ 的有限仿射开覆盖
$X=\bigcup_{\lambda\in\Lambda}U_\lambda$，并在 $\Lambda$ 上选取一个全序。
写成 $U_\lambda=\Spec(C_\lambda)$。选取 $A$ 上的多项式代数
$P_\lambda$ 及满射 $P_\lambda\to C_\lambda$。固定这些选择后，可构造
一个计算 $R\Gamma(\text{Cris}(X/S),\mathcal F)$ 的交错 {\v C}ech 复形。

\medskip\noindent
采用 注 \ref{crystalline-remark-cech-complex} 中引入的记号。以
$\Omega_{\lambda_0\ldots\lambda_n}$ 表示
$D_{\lambda_0\ldots\lambda_n}$ 在 $A$ 上与除幂结构相容的微分模之
$p$-进完备化。令 $\nabla$ 为 命题
\ref{crystalline-proposition-crystals-on-affine} 所给出的
$M_{\lambda_0\ldots\lambda_n}$ 上可积联络。考虑双复形
$M^{\bullet,\bullet}$，其各项为
$$
M^{n, m} =
\bigoplus\nolimits_{\lambda_0 < \ldots < \lambda_n}
M_{\lambda_0 \ldots \lambda_n}
\otimes^\wedge_{D_{\lambda_0 \ldots \lambda_n}}
\Omega^m_{D_{\lambda_0 \ldots \lambda_n}}.
$$
微分 $d_1$（使 $n$ 增加）取通常的 {\v C}ech 微分，而微分 $d_2$
取联络，即德拉姆复形的微分。我们断言
$$
R\Gamma(\text{Cris}(X/S), \mathcal{F}) = \text{Tot}(M^{\bullet, \bullet})
$$
其中 $\text{Tot}(-)$ 表示与双复形相伴的全复形，见 “同调”
定义 \ref{homology-definition-associated-simple-complex}。

\medskip\noindent
提示：由命题 \ref{crystalline-proposition-compare-with-de-Rham} 有
$$
R\Gamma(\text{Cris}(X/S), \mathcal{F}) = R\Gamma(\text{Cris}(X/S),
\mathcal{F} \otimes_{\mathcal{O}_{X/S}} \Omega_{X/S}^\bullet)
$$
。公式右端正是覆盖 $X=\bigcup_{\lambda\in\Lambda}U_\lambda$
（它诱导 $\text{Cris}(X/S)$ 的终层之一个开覆盖）与复形
$\mathcal F\otimes_{\mathcal O_{X/S}}\Omega_{X/S}^\bullet$ 的交错
{\v C}ech 复形；见命题
\ref{crystalline-proposition-compute-cohomology-crystal}。结论现由位点上同调的一个
一般结果推出：若交错 {\v C}ech 复形在每个部分上给出正确答案，它就
计算上同调（以后在此插入引用）。
\end{remark}

\begin{remark}[拟凝聚性]
\label{crystalline-remark-quasi-coherent}
在 注 \ref{crystalline-remark-compute-direct-image} 的情形中，假设 $S\to S'$
与 $X\to S_0$ 都拟紧且拟分离。则对拟凝聚 $\mathcal O_{X/S}$-模中的
晶体 $\mathcal F$，各层 $R^if_{\text{cris},*}\mathcal F$ 均局部拟凝聚。

\medskip\noindent
提示：须证明其在 $T'$ 上的限制是拟凝聚 $\mathcal O_{T'}$-模，其中
$(U',T',\delta')$ 为 $\text{Cris}(X'/S')$ 的任意对象。只需在 $T'$
仿射时证明。利用公式 (\ref{crystalline-equation-identify-pushforward})、$T\to T'$
拟紧且拟分离这一事实（因为 $T$ 在 $T'$ 沿 $S\to S'$ 的基变换上
仿射），以及“概形上同调”引理
\ref{coherent-lemma-quasi-coherence-higher-direct-images}，可知只需证明
各层 $R^i\tau_{U/T,*}\mathcal F_U$ 拟凝聚。注意 $U\to T_0$ 也拟紧且
拟分离；见“概形”引理 \ref{schemes-lemma-quasi-compact-permanence}
与 \ref{schemes-lemma-quasi-compact-permanence}。

\medskip\noindent
这把问题化为：当 $p$ 在 $S$ 上局部幂零时，证明
$R^i\tau_{X/S,*}\mathcal F$ 在 $S$ 上拟凝聚。这里 $\tau_{X/S}$ 是结构
态射，见注 \ref{crystalline-remark-structure-morphism}。可以在 $S$ 上局部地
工作，故可假设 $S$ 仿射（见引理 \ref{crystalline-lemma-localize}）。对覆盖 $X$
的仿射开集数目作归纳，并使用 Mayer--Vietoris（注
\ref{crystalline-remark-mayer-vietoris}），可把问题化归到 $X$ 也仿射的情形
（如“概形上同调”引理
\ref{coherent-lemma-quasi-coherence-higher-direct-images} 的证明）。设
$X=\Spec(C)$、$S=\Spec(A)$，使得 $(A,I,\gamma)$ 与 $A\to C$ 如
情形 \ref{crystalline-situation-affine}。依第
\ref{crystalline-section-quasi-coherent-crystals} 节选取 $A$ 上多项式代数 $P$ 及满射
$P\to C$。令 $(M,\nabla)$ 为与 $\mathcal F$ 对应的模，见命题
\ref{crystalline-proposition-crystals-on-affine}。应用命题
\ref{crystalline-proposition-compute-cohomology-crystal}，可知
$R\Gamma(\text{Cris}(X/S),\mathcal F)$ 由 $M\otimes_D\Omega_D^*$ 表示。
注意，由于 $p$ 在 $A$ 中幂零，无需完备化！须证明这与在
$S=\Spec(A)$ 中取主开集相容。设 $g\in A$。类似地，
$R\Gamma(\text{Cris}(X_g/S_g),\mathcal F)$ 由
$M_g\otimes_{D_g}\Omega_{D_g}^*$ 计算（这里再次使用无需 $p$-进完备化）。
由于局部化是 $A$-模上的正合函子，结论随即成立。
\end{remark}

\begin{remark}[有界性]
\label{crystalline-remark-bounded-cohomology}
在 注 \ref{crystalline-remark-compute-direct-image} 的情形中，假设 $S\to S'$
拟紧且拟分离，而 $X\to S_0$ 为有限型且拟分离。则存在整数 $i_0$，
使得对拟凝聚 $\mathcal O_{X/S}$-模中的任意晶体 $\mathcal F$，当
$i>i_0$ 时均有 $R^if_{\text{cris},*}\mathcal F=0$。

\medskip\noindent
提示：仿照注 \ref{crystalline-remark-quasi-coherent} 的论证（使用“概形上同调”引理
\ref{coherent-lemma-quasi-coherence-higher-direct-images}），问题化为在
命题 \ref{crystalline-proposition-compute-cohomology-crystal} 的情形中，当 $C$
是 $A$ 上有限型代数时，证明 $i\gg0$ 有
$H^i(\text{Cris}(X/S),\mathcal F)=0$。这是显然的，因为可以选取有限
多项式代数，而 $i\gg0$ 时 $\Omega_D^i=0$。
\end{remark}

\begin{remark}[具体的有界性]
\label{crystalline-remark-bounded-cohomology-over-point}
在 情形 \ref{crystalline-situation-global} 中，设 $\mathcal F$ 为拟凝聚
$\mathcal O_{X/S}$-模中的晶体。假设 $S_0$ 仅有一个点，且
$X\to S_0$ 为有限表示。
\begin{enumerate}
\item 若 $\dim X=d$ 且 $X/S_0$ 的嵌入维数为 $e$，则 $i>d+e$ 时
$H^i(\text{Cris}(X/S),\mathcal F)=0$。
\item 若 $X$ 分离且可由 $q$ 个仿射开集覆盖，并且 $X/S_0$ 的嵌入维数
为 $e$，则 $i>q+e$ 时 $H^i(\text{Cris}(X/S),\mathcal F)=0$。
\end{enumerate}
提示：在情形 (1) 中可使用
$$
H^i(\text{Cris}(X/S), \mathcal{F}) = H^i(X_{Zar}, Ru_{X/S, *}\mathcal{F})
$$
以及如下事实：$Ru_{X/S,*}\mathcal F$ 局部地由一个德拉姆复形计算，
该复形利用 $X$ 到 $S$ 上维数为 $e$ 的光滑概形中的嵌入构造
（见引理 \ref{crystalline-lemma-compute-cohomology-crystal-smooth}）。这些德拉姆
复形在所有次数 $>e$ 处为零。因此 (1) 由 “上同调” 命题
\ref{cohomology-proposition-vanishing-Noetherian} 推出。在情形 (2) 中，
使用交错 {\v C}ech 复形（见注
\ref{crystalline-remark-alternating-cech-complex}）化归到 $X$ 仿射的情形。仿射
情形下，利用与 $X$ 到 $S$ 上维数为 $e$ 的光滑概形中某个嵌入相伴的
德拉姆复形证明结果（构造这样的嵌入需要一些工作）。
\end{remark}

\begin{remark}[基变换映射]
\label{crystalline-remark-base-change}
在 注 \ref{crystalline-remark-compute-direct-image} 的情形中，假设
$S=\Spec(A)$ 与 $S'=\Spec(A')$ 仿射。设 $\mathcal F'$ 为
$\mathcal O_{X'/S'}$-模，并设 $\mathcal F$ 为 $\mathcal F'$ 的拉回。
则存在典范基变换映射
$$
L(S' \to S)^*R\tau_{X'/S', *}\mathcal{F}'
\longrightarrow
R\tau_{X/S, *}\mathcal{F}
$$
其中 $\tau_{X/S}$ 与 $\tau_{X'/S'}$ 是结构态射，见注
\ref{crystalline-remark-structure-morphism}。在整体截面上，这给出基变换映射
\begin{equation}
\label{crystalline-equation-base-change-map}
R\Gamma(\text{Cris}(X'/S'), \mathcal{F}') \otimes^\mathbf{L}_{A'} A
\longrightarrow
R\Gamma(\text{Cris}(X/S), \mathcal{F})
\end{equation}
于 $D(A)$ 中。

\medskip\noindent
提示：把“位点上同调”注
\ref{sites-cohomology-remark-base-change} 的极一般基变换映射与典范映射
$Lf_{\text{cris}}^*\mathcal{F}' \to
f_{\text{cris}}^*\mathcal{F}' = \mathcal{F}$.
\end{remark}

\begin{remark}[基变换同构]
\label{crystalline-remark-base-change-isomorphism}
若以下条件全部满足，则映射 (\ref{crystalline-equation-base-change-map}) 为同构：
\begin{enumerate}
\item $p$ 在 $A'$ 中幂零；
\item $\mathcal F'$ 是拟凝聚 $\mathcal O_{X'/S'}$-模中的晶体；
\item $X'\to S'_0$ 为拟紧拟分离态射；
\item $X = X' \times_{S'_0} S_0$,
\item $\mathcal F'$ 是平坦 $\mathcal O_{X'/S'}$-模；
\item $X'\to S'_0$ 为局部完全交态射（见“态射进阶”
定义 \ref{more-morphisms-definition-lci}；例如，当 $X'\to S'_0$
为 syntomic 或光滑态射时成立）；
\item $X'$ 与 $S_0$ 在 $S'_0$ 上 Tor 独立（见“代数进阶”
定义 \ref{more-algebra-definition-tor-independent}；例如，当
$S_0\to S'_0$ 或 $X'\to S'_0$ 中任一个平坦时成立）。
\end{enumerate}
提示：条件 (1) 意味着在以下论证中，$p$-进完备化不起作用，因而可以忽略。
利用条件 (3) 与 Mayer--Vietoris（见注
\ref{crystalline-remark-mayer-vietoris}），问题化归到 $X'$ 仿射的情形。事实上，
由条件 (6)，进一步缩小后可假设 $X'=\Spec(C')$，并给定表示
$C'=A'/I'[x_1,\ldots,x_n]/(\bar f'_1,\ldots,\bar f'_c)$，其中
$\bar f'_1,\ldots,\bar f'_c$ 是 $A'/I'$ 中的 Koszul-正则列。（这意味着
光滑局部地，$\bar f'_1,\ldots,\bar f'_c$ 构成正则列；见“代数进阶”引理
\ref{more-algebra-lemma-Koszul-regular-flat-locally-regular}。）选取
$\bar f'_i$ 到元素 $f'_i\in A'[x_1,\ldots,x_n]$ 的提升。由 (4)，
$X=\Spec(C)$，其中
$C=A/I[x_1,\ldots,x_n]/(\bar f_1,\ldots,\bar f_c)$，而
$f_i\in A[x_1,\ldots,x_n]$ 是 $f'_i$ 的像。由性质 (7)，
$\bar f_1,\ldots,\bar f_c$ 是 $A/I[x_1,\ldots,x_n]$ 中的
Koszul-正则列。$A'[x_1,\ldots,x_n]$ 中理想
$I'A'[x_1,\ldots,x_n]+(f'_1,\ldots,f'_c)$ 相对于 $\gamma'$ 的除幂包络为
$$
D' = A'[x_1, \ldots, x_n]\langle \xi_1, \ldots, \xi_c \rangle/(\xi_i - f'_i)
$$
，见引理 \ref{crystalline-lemma-describe-divided-power-envelope}。然后可检验
$\xi_1-f'_1,\ldots,\xi_n-f'_n$ 是环
$A'[x_1,\ldots,x_n]\langle\xi_1,\ldots,\xi_c\rangle$ 中的
Koszul-正则列。类似地，$A[x_1,\ldots,x_n]$ 中理想
$IA[x_1,\ldots,x_n]+(f_1,\ldots,f_c)$ 相对于 $\gamma$ 的除幂包络为
$$
D = A[x_1, \ldots, x_n]\langle \xi_1, \ldots, \xi_c\rangle/(\xi_i - f_i)
$$
，且 $\xi_1-f_1,\ldots,\xi_n-f_n$ 是环
$A[x_1,\ldots,x_n]\langle\xi_1,\ldots,\xi_c\rangle$ 中的
Koszul-正则列。于是 $D'\otimes_{A'}^\mathbf L A=D$。条件 (2) 蕴含
$\mathcal F'$ 对应于一个由带联络 $D'$-模组成的偶对 $(M',\nabla)$，
见命题 \ref{crystalline-proposition-crystals-on-affine}。则
$M=M'\otimes_{D'}D$ 对应于拉回 $\mathcal F$。由假设 (5)，$M'$ 是平坦
$D'$-模，故
$$
M = M' \otimes_{D'} D = M' \otimes_{D'} D' \otimes_{A'}^\mathbf{L} A
= M' \otimes_{A'}^\mathbf{L} A
$$
由于微分模 $\Omega_{D'}$ 与 $\Omega_D$（定义见节
\ref{crystalline-section-quasi-coherent-crystals}）是由相同生成元生成的自由
$D'$-模，可得
$$
M \otimes_D \Omega^\bullet_D =
M' \otimes_{D'} \Omega^\bullet_{D'} \otimes_{D'} D =
M' \otimes_{D'} \Omega^\bullet_{D'} \otimes_{A'}^\mathbf{L} A
$$
，再由命题 \ref{crystalline-proposition-compute-cohomology-crystal} 即得所需。
\end{remark}

\begin{remark}[Rlim]
\label{crystalline-remark-rlim}
设 $p$ 为素数，并设 $(A,I,\gamma)$ 为除幂环，其中 $A$ 是
$\mathbf Z_{(p)}$-代数，且 $p$ 在 $A/I$ 中幂零。置 $S=\Spec(A)$ 与
$S_0=\Spec(A/I)$。设 $X$ 为 $S_0$ 上概形，且 $p$ 在 $X$ 上局部幂零。
设 $\mathcal F$ 为任意 $\mathcal O_{X/S}$-模。当 $e\gg0$ 时，
$(p^e)\subset I$ 被 $\gamma$ 保持，见“除幂代数”引理
\ref{dpa-lemma-extend-to-completion}。对 $e\gg0$ 置
$S_e=\Spec(A/p^eA)$。则 $\text{Cris}(X/S_e)$ 是
$\text{Cris}(X/S)$ 的全子范畴，并以 $\mathcal F_e$ 表示 $\mathcal F$
在 $\text{Cris}(X/S_e)$ 上的限制。于是
$$
R\Gamma(\text{Cris}(X/S), \mathcal{F}) =
R\lim_e R\Gamma(\text{Cris}(X/S_e), \mathcal{F}_e)
$$

\medskip\noindent
提示：只需对内射的 $\mathcal F$ 证明。此时各层 $\mathcal F_e$ 也都是
内射模，转移映射 $\Gamma(\mathcal F_{e+1})\to\Gamma(\mathcal F_e)$
满射，并且 $\Gamma(\mathcal F)=\lim_e\Gamma(\mathcal F_e)$；这是因为按
$\text{Cris}(X/S)$ 的定义，其任意对象局部地都是某个范畴
$\text{Cris}(X/S_e)$ 的对象。
\end{remark}

\begin{remark}[比较]
\label{crystalline-remark-comparison}
设 $p$ 为素数，并设 $(A,I,\gamma)$ 为除幂环，其中 $p$ 在 $A$ 中幂零。
置 $S=\Spec(A)$ 与 $S_0=\Spec(A/I)$。设 $Y$ 为 $S$ 上光滑概形，并置
$X=Y\times_SS_0$。设 $\mathcal F$ 为拟凝聚 $\mathcal O_{X/S}$-模中的
晶体。则
\begin{enumerate}
\item $\gamma$ 延拓为 $X$ 在 $Y$ 中的理想上的除幂结构，使
$(X,Y,\gamma)$ 成为 $\text{Cris}(X/S)$ 的对象；
\item 限制 $\mathcal F_Y$（见第 \ref{crystalline-section-sheaves} 节）带有典范
可积联络
$\nabla : \mathcal{F}_Y \to
\mathcal{F}_Y \otimes_{\mathcal{O}_Y} \Omega_{Y/S}$；
\item 有
$$
R\Gamma(\text{Cris}(X/S), \mathcal{F}) =
R\Gamma(Y, \mathcal{F}_Y \otimes_{\mathcal{O}_Y} \Omega^\bullet_{Y/S})
$$
于 $D(A)$ 中。
\end{enumerate}
提示：(1) 见“除幂代数”引理
\ref{dpa-lemma-gamma-extends}；(2) 见引理
\ref{crystalline-lemma-automatic-connection}。对 (3)，注意存在映射
(\ref{crystalline-equation-restriction})。当 $X$ 仿射时，此映射为同构，见引理
\ref{crystalline-lemma-compute-cohomology-crystal-smooth}。这说明在
$Y_{Zar}=X_{Zar}$ 上，复形 $Ru_{X/S,*}\mathcal F$ 与
$\mathcal F_Y\otimes\Omega^\bullet_{Y/S}$ 拟同构。由于
$R\Gamma(\text{Cris}(X/S),\mathcal F)=
R\Gamma(X_{Zar},Ru_{X/S,*}\mathcal F)$，结论成立。
\end{remark}

\begin{remark}[完美性]
\label{crystalline-remark-perfect}
设 $p$ 为素数，并设 $(A,I,\gamma)$ 为除幂环，其中 $p$ 在 $A$ 中幂零。
置 $S=\Spec(A)$ 与 $S_0=\Spec(A/I)$。设 $X$ 为 $S_0$ 上固有光滑概形。
设 $\mathcal F$ 为有限局部自由拟凝聚 $\mathcal O_{X/S}$-模中的晶体。
则 $R\Gamma(\text{Cris}(X/S),\mathcal F)$ 是 $D(A)$ 中的完美对象。

\medskip\noindent
提示：由注 \ref{crystalline-remark-base-change-isomorphism} 有
$$
R\Gamma(\text{Cris}(X/S), \mathcal{F}) \otimes_A^\mathbf{L} A/I
\cong
R\Gamma(\text{Cris}(X/S_0), \mathcal{F}|_{\text{Cris}(X/S_0)})
$$
由注 \ref{crystalline-remark-comparison} 有
$$
R\Gamma(\text{Cris}(X/S_0), \mathcal{F}|_{\text{Cris}(X/S_0)}) =
R\Gamma(X, \mathcal{F}_X \otimes \Omega^\bullet_{X/S_0})
$$
在德拉姆复形上使用朴素滤过，可知只要各复形
$$
R\Gamma(X, \mathcal{F}_X \otimes \Omega^q_{X/S_0})
$$
在 $D(A/I)$ 中完美，前一展示复形就在 $D(A/I)$ 中完美；见“代数进阶”
引理 \ref{more-algebra-lemma-two-out-of-three-perfect}。利用
$\mathcal F_X\otimes\Omega^q_{X/S_0}$ 是有限局部自由层以及
$X\to S_0$ 固有且平坦，凝聚上同调中的标准论证可知此条件成立
（以后在此插入引用）。应用“代数进阶”引理
\ref{more-algebra-lemma-perfect-modulo-nilpotent-ideal}，可知
$$
R\Gamma(\text{Cris}(X/S), \mathcal{F}) \otimes_A^\mathbf{L} A/I^n
$$
对所有 $n$ 都是 $D(A/I^n)$ 中的完美对象。除非 $A$ 为 Noether 环，
这还不够。确切地说，尽管由 $p$ 幂零的假设，$I$ 局部幂零，见
“除幂代数”引理 \ref{dpa-lemma-nil}，仍不能推出存在 $n$
使 $I^n=0$。反例为 $\mathbf F_p\langle x\rangle$。当
$\mathcal F=\mathcal O_{X/S}$ 时，一般情形可用
\url{https://math.columbia.edu/~dejong/wordpress/?p=2227}
中的论证证明。当系数 $\mathcal F$ 非平凡时，\cite{Faltings-very} 的
论证似乎如下。由“代数进阶”引理
\ref{more-algebra-lemma-perfect-modulo-nilpotent-ideal}，化归到 $pA=0$
的情形。此时 Frobenius 映射 $A\to A$、$a\mapsto a^p$ 分解为
$A\to A/I\xrightarrow{\varphi}A$（因为 $x\in I$ 时 $x^p=0$）。置
$X^{(1)}=X\otimes_{A/I,\varphi}A$。$X$ 的绝对 Frobenius 态射经过态射
$F_X:X\to X^{(1)}$ 分解（某种相对 Frobenius）。仿射局部地，若
$X=\Spec(C)$，则 $X^{(1)}=\Spec(C\otimes_{A/I,\varphi}A)$，而 $F_X$
对应于 $C\otimes_{A/I,\varphi}A\to C$、$c\otimes a\mapsto c^pa$。
这定义环化拓扑斯的态射
$$
(X/S)_{\text{cris}}
\xrightarrow{(F_X)_{\text{cris}}}
(X^{(1)}/S)_{\text{cris}}
\xrightarrow{u_{X^{(1)}/S}}
\Sh(X^{(1)}_{Zar})
$$
，其复合记为 $\text{Frob}_X$。然后通过局部计算证明
$R\text{Frob}_{X,*}\mathcal F$ 可由 $\mathcal O_{X^{(1)}}$-模的一个
完美复形表示（！）。
\end{remark}

\begin{remark}[完备情形的完美性]
\label{crystalline-remark-complete-perfect}
设 $p$ 为素数，并设 $(A,I,\gamma)$ 为除幂环，其中 $A$ 是 $p$-进完备环，
且 $p$ 在 $A/I$ 中幂零。置 $S=\Spec(A)$ 与 $S_0=\Spec(A/I)$。设 $X$
为 $S_0$ 上固有光滑概形，并设 $\mathcal F$ 为有限局部自由拟凝聚
$\mathcal O_{X/S}$-模中的晶体。则
$R\Gamma(\text{Cris}(X/S),\mathcal F)$ 是 $D(A)$ 中的完美对象。

\medskip\noindent
提示：已知 $K=R\Gamma(\text{Cris}(X/S),\mathcal F)$ 是 $A/p^eA$ 上
各上同调的导出极限 $K=R\lim K_e$，见注 \ref{crystalline-remark-rlim}。
由注 \ref{crystalline-remark-perfect}，每个 $K_e$ 都是 $D(A/p^eA)$ 中的完美
复形。由于 $A$ 是 $p$-进完备的，结论由“代数进阶”引理
\ref{more-algebra-lemma-Rlim-perfect-gives-complete} 推出。
\end{remark}

\begin{remark}[完备情形的比较]
\label{crystalline-remark-complete-comparison}
设 $p$ 为素数，并设 $(A,I,\gamma)$ 为除幂环，其中 $A$ 为 Noether
$p$-进完备环，且 $p$ 在 $A/I$ 中幂零。置 $S=\Spec(A)$ 与
$S_0=\Spec(A/I)$。设 $Y$ 为 $S$ 上固有光滑概形，并置
$X=Y\times_SS_0$。设 $\mathcal F$ 为拟凝聚 $\mathcal O_{X/S}$-模中的
有限型晶体。则
\begin{enumerate}
\item 存在带可积联络的凝聚 $\mathcal O_Y$-模 $\mathcal F_Y$
$$
\nabla :
\mathcal{F}_Y
\longrightarrow
\mathcal{F}_Y \otimes_{\mathcal{O}_Y} \Omega_{Y/S}
$$
，使得 $\mathcal F_Y/p^e\mathcal F_Y$ 是 注
\ref{crystalline-remark-comparison} 中所得的 $A/p^eA$ 上带联络模；
\item 有
$$
R\Gamma(\text{Cris}(X/S), \mathcal{F}) =
R\Gamma(Y, \mathcal{F}_Y \otimes_{\mathcal{O}_Y} \Omega^\bullet_{Y/S})
$$
于 $D(A)$ 中。
\end{enumerate}
提示：$\mathcal F_Y$ 的存在性即 Grothendieck 存在定理（以后在此插入
引用）。两端都由模 $p^e$ 版本的 $R\lim$ 计算，故得到上同调的同构
（左端见注 \ref{crystalline-remark-rlim}；右端使用形式函数定理，见“概形上同调”
定理 \ref{coherent-theorem-formal-functions}）。由注
\ref{crystalline-remark-comparison}，各模 $p^e$ 版本彼此同构。
\end{remark}




\section{沿纯不可分映射拉回}
\label{crystalline-section-pull-back-along-pth-root}

\noindent
所谓 $\alpha_p$-覆盖，是指形如
$$
X' = \Spec(C[z]/(z^p - c)) \longrightarrow \Spec(C) = X
$$
的态射，其中 $C$ 为 $\mathbf F_p$-代数且 $c\in C$。等价地，$X'$ 是
$X$ 上的 $\alpha_p$-挠子。所谓{\it 迭代 $\alpha_p$-覆盖}
\footnote{这不是标准记号。}，是指特征 $p$ 的概形态射，它在目标上
局部地是有限多个 $\alpha_p$-覆盖的复合。本节证明，沿这类态射的拉回
在逆转素数 $p$ 后诱导晶体上同调的拟同构。事实上，我们将证明该结果
的一个精确版本。先给出一个需要若干记号才能表述的预备引理。

\medskip\noindent
假设给定环映射 $B\to B'$ 及满足 注
\ref{crystalline-remark-base-change-connection} 假设的商
$\Omega_B\to\Omega$ 与 $\Omega_{B'}\to\Omega'$。于是
(\ref{crystalline-equation-base-change-map-complexes}) 给出复形的典范映射
$$
c_M^\bullet :
M \otimes_B \Omega^\bullet
\longrightarrow
M \otimes_B (\Omega')^\bullet
$$
，其中 $M$ 为任意带可积联络
$\nabla:M\to M\otimes_B\Omega_B$ 的 $B$-模。

\medskip\noindent
再假设给定 $a\in B$、$z\in B'$ 及映射 $\theta:B'\to B'$，满足：
\begin{enumerate}
\item
\label{crystalline-item-d-a-zero}
$\text{d}(a) = 0$,
\item
\label{crystalline-item-direct-sum}
$\Omega'=B'\otimes_B\Omega\oplus B'\text d z$；对所有 $f\in B'$ 写成
$\text{d}(f) = \text{d}_1(f) + \partial_z(f) \text{d}z$
，其中 $\text d_1(f)\in B'\otimes\Omega$ 且 $\partial_z(f)\in B'$；
\item
\label{crystalline-item-theta-linear}
$\theta:B'\to B'$ 为 $B$-线性；
\item
\label{crystalline-item-integrate}
$\partial_z \circ \theta = a$,
\item
\label{crystalline-item-injective}
$B\to B'$ 泛单射（因而 $\Omega\to\Omega'$ 单射）；
\item
\label{crystalline-item-factor}
对所有 $f\in B'$，$af-\theta(\partial_z(f))\in B$；
\item
\label{crystalline-item-horizontal}
$(\theta \otimes 1)(\text{d}_1(f)) - \text{d}_1(\theta(f)) \in \Omega$
对所有 $f\in B'$，其中
$\theta \otimes 1 : B' \otimes \Omega \to B' \otimes \Omega$
\end{enumerate}
这些条件在逻辑上并非彼此独立。例如，假设 (\ref{crystalline-item-integrate}) 蕴含
$\partial_z(af-\theta(\partial_z(f)))=0$。因此，若 $B\to B'$ 的像恰为
被 $\partial_z$ 消去的元素全体，则 (\ref{crystalline-item-factor}) 随即成立。
对条件 (\ref{crystalline-item-horizontal}) 可作类似论证。

\begin{lemma}
\label{crystalline-lemma-find-homotopy}
在上述情形中，存在复形映射
$$
e_M^\bullet :
M \otimes_B (\Omega')^\bullet
\longrightarrow
M \otimes_B \Omega^\bullet
$$
，使 $c_M^\bullet\circ e_M^\bullet$ 与
$e_M^\bullet\circ c_M^\bullet$ 都同伦于乘以 $a$ 的映射。
\end{lemma}

\begin{proof}
本证明中的一切张量积均取于 $B$ 上。假设 (\ref{crystalline-item-direct-sum}) 蕴含
$$
M \otimes (\Omega')^i =
(B' \otimes M \otimes \Omega^i)
\oplus
(B' \text{d}z \otimes M \otimes \Omega^{i - 1})
$$
对所有 $i\geq0$ 成立。$M\otimes(\Omega')^i$ 的一组加法生成元由形如
$f\omega$ 与 $f\text d z\wedge\eta$ 的元素组成，其中 $f\in B'$、
$\omega\in M\otimes\Omega^i$ 且 $\eta\in M\otimes\Omega^{i-1}$。

\medskip\noindent
对 $f\in B'$，记
$$
\epsilon(f) = af - \theta(\partial_z(f))
\quad\text{且}\quad
\epsilon'(f) = (\theta \otimes 1)(\text{d}_1(f)) - \text{d}_1(\theta(f))
$$
由假设 (\ref{crystalline-item-factor}) 与 (\ref{crystalline-item-horizontal})，有
$\epsilon(f)\in B$ 且 $\epsilon'(f)\in\Omega$。按规则
$e^i_M(f\omega)=\epsilon(f)\omega$ 与
$e^i_M(f\text d z\wedge\eta)=\epsilon'(f)\wedge\eta$ 定义
$e_M^\bullet$。下面将看到，映射族 $e_M^i$ 构成复形映射。

\medskip\noindent
定义
$$
h^i : M \otimes_B (\Omega')^i \longrightarrow M \otimes_B (\Omega')^{i - 1}
$$
，对上述元素规定 $h^i(f\omega)=0$ 以及
$h^i(f\text d z\wedge\eta)=\theta(f)\eta$。我们断言
$$
\text{d} \circ h + h \circ \text{d} = a - c_M^\bullet \circ e_M^\bullet
$$
由 (\ref{crystalline-item-d-a-zero})，乘以 $a$ 是复形映射。又由假设
(\ref{crystalline-item-injective})，$c_M^\bullet$ 是单射复形映射，故可知
$e_M^\bullet$ 是复形映射。为证明断言，计算
\begin{align*}
(\text{d} \circ h + h \circ \text{d})(f \omega)
& =
h\left(\text{d}(f) \wedge \omega + f \nabla(\omega)\right)
\\
& =
\theta(\partial_z(f)) \omega
\\
& =
a f\omega - \epsilon(f)\omega 
\\
& =
a f \omega - c^i_M(e^i_M(f\omega))
\end{align*}
第二个等号成立是因为 $\nabla(\omega)$ 中不出现 $\text d z$，第三个
等号则由假设 (6) 得到。类似地，
\begin{align*}
(\text{d} \circ h + h \circ \text{d})(f \text{d}z \wedge \eta)
& =
\text{d}(\theta(f) \eta) +
h\left(\text{d}(f) \wedge \text{d}z \wedge \eta -
f \text{d}z \wedge \nabla(\eta)\right)
\\
& =
\text{d}(\theta(f)) \wedge \eta + \theta(f) \nabla(\eta)
- (\theta \otimes 1)(\text{d}_1(f)) \wedge \eta
- \theta(f) \nabla(\eta)
\\
& =
\text{d}_1(\theta(f)) \wedge \eta +
\partial_z(\theta(f)) \text{d}z \wedge \eta -
(\theta \otimes 1)(\text{d}_1(f)) \wedge \eta
\\
& =
a f \text{d}z \wedge \eta - \epsilon'(f) \wedge \eta \\
& = a f \text{d}z \wedge \eta - c^i_M(e^i_M(f \text{d}z \wedge \eta))
\end{align*}
第二个等号成立是因为
$\text{d}(f) \wedge \text{d}z \wedge \eta =
- \text{d}z \wedge \text{d}_1(f) \wedge \eta$.
第四个等号由假设 (\ref{crystalline-item-integrate}) 得到。另一方面，由定义立即有
$e^i_M(c^i_M(\omega))=\epsilon(1)\omega=a\omega$。引理得证。
\end{proof}

\begin{example}
\label{crystalline-example-integrate}
上述情形的一个标准例子如下：$B'=B\langle z\rangle$ 是除幂环
$(B,J,\delta)$ 上的除幂多项式环，而
$J'=B'_++JB'\subset B'$ 带除幂 $\delta'$。具体地，取
$\Omega=\Omega_{B,\delta}$ 与 $\Omega'=\Omega_{B',\delta'}$。
此时可取 $a=1$ 以及
$$
\theta( \sum b_m z^{[m]} ) = \sum b_m z^{[m + 1]}
$$
注意
$$
f - \theta(\partial_z(f)) = f(0)
$$
等于常数项。因此在此情形下，引理 \ref{crystalline-lemma-find-homotopy}
恢复晶体 Poincar\'e 引理（引理 \ref{crystalline-lemma-relative-poincare}）。
\end{example}

\begin{lemma}
\label{crystalline-lemma-computation}
处于 情形 \ref{crystalline-situation-affine}。假设 $D$ 与 $\Omega_D$ 如
(\ref{crystalline-equation-D}) 和 (\ref{crystalline-equation-omega-D})。设 $\lambda\in D$，并令
$D'$ 为下列环的 $p$-进完备化：
$$
D[z]\langle \xi \rangle/(\xi - (z^p - \lambda))
$$
；并令 $\Omega_{D'}$ 为 $D'$ 在 $A$ 上的除幂微分模之 $p$-进完备化。
对任意满足 (\ref{crystalline-item-complete})、(\ref{crystalline-item-connection})、
(\ref{crystalline-item-integrable}) 与 (\ref{crystalline-item-topologically-quasi-nilpotent})
的 $D$ 上偶对 $(M,\nabla)$，复形的典范映射
(\ref{crystalline-equation-base-change-map-complexes})
$$
c_M^\bullet : M \otimes_D^\wedge \Omega^\bullet_D
\longrightarrow
M \otimes_D^\wedge \Omega^\bullet_{D'}
$$
具有如下性质：存在反向映射 $e_M^\bullet$，使
$c_M^\bullet\circ e_M^\bullet$ 与 $e_M^\bullet\circ c_M^\bullet$
都同伦于乘以 $p$ 的映射。
\end{lemma}

\begin{proof}
取 $a=p$，使用 引理 \ref{crystalline-lemma-find-homotopy} 来证明。故须找出
$\theta:D'\to D'$ 并证明
(\ref{crystalline-item-d-a-zero}), (\ref{crystalline-item-direct-sum}), (\ref{crystalline-item-theta-linear}),
(\ref{crystalline-item-integrate}), (\ref{crystalline-item-injective}), (\ref{crystalline-item-factor}),
(\ref{crystalline-item-horizontal})。先汇集关于环 $D,D'$ 与模
$\Omega_D,\Omega_{D'}$ 的若干信息。

\medskip\noindent
写成
$$
D[z]\langle \xi \rangle/(\xi - (z^p - \lambda))
=
D\langle \xi \rangle[z]/(z^p - \xi - \lambda)
$$
可见 $D'$ 是下列自由 $D$-模的 $p$-进完备化：
$$
\bigoplus\nolimits_{i = 0, \ldots, p - 1}
\bigoplus\nolimits_{n \geq 0}
z^i \xi^{[n]} D
$$
，其中 $\xi^{[0]}=1$。因此 $D\to D'$ 有连续的 $D$-线性截面；特别地，
$D\to D'$ 泛单射，即 (\ref{crystalline-item-injective}) 成立。把 $D'$ 看作 $A$ 上
具有除幂理想 $\overline J'=\overline JD'+(\xi)$ 的除幂代数。于是 $D'$
也是 $D$ 中由 $z^p-\lambda$ 生成的理想之除幂包络的 $p$-进完备化，
见引理 \ref{crystalline-lemma-describe-divided-power-envelope}。故
$$
\Omega_{D'} = \Omega_D \otimes_D^\wedge D' \oplus D'\text{d}z
$$
，由引理 \ref{crystalline-lemma-module-differentials-divided-power-envelope} 得到。
这证明了 (\ref{crystalline-item-direct-sum})。注意 (\ref{crystalline-item-d-a-zero}) 显然成立。

\medskip\noindent
现在构造 $\theta$。（我们编写了一个 PARI/gp 脚本 theta.gp 来检验本证明
中的若干公式；该脚本位于 Stacks project 的 scripts 子目录中。）在此
之前先计算元素 $z^i\xi^{[n]}$ 的导数。有
$\text d z^i=iz^{i-1}\text d z$。当 $n\geq1$ 时，
$$
\text{d}\xi^{[n]} =
\xi^{[n - 1]} \text{d}\xi =
- \xi^{[n - 1]}\text{d}\lambda + p z^{p - 1} \xi^{[n - 1]}\text{d}z
$$
，因为 $\xi=z^p-\lambda$。当 $0<i<p$ 且 $n\geq1$ 时，有
\begin{align*}
\text{d}(z^i\xi^{[n]})
& =
iz^{i - 1}\xi^{[n]}\text{d}z + z^i\xi^{[n - 1]}\text{d}\xi \\
& =
iz^{i - 1}\xi^{[n]}\text{d}z + z^i\xi^{[n - 1]}\text{d}(z^p - \lambda) \\
& =
- z^i\xi^{[n - 1]}\text{d}\lambda +
(iz^{i - 1}\xi^{[n]} + pz^{i + p - 1}\xi^{[n - 1]})\text{d}z \\
& =
- z^i\xi^{[n - 1]}\text{d}\lambda +
(iz^{i - 1}\xi^{[n]} + pz^{i - 1}(\xi + \lambda)\xi^{[n - 1]})\text{d}z \\
& =
- z^i\xi^{[n - 1]}\text{d}\lambda +
((i + pn)z^{i - 1}\xi^{[n]} + p\lambda z^{i - 1}\xi^{[n - 1]})\text{d}z
\end{align*}
最后一个等号来自 $\xi\xi^{[n-1]}=n\xi^{[n]}$。因此
\begin{align*}
\partial_z(z^i) & = i z^{i - 1} \\
\partial_z(\xi^{[n]}) & = p z^{p - 1} \xi^{[n - 1]} \\
\partial_z(z^i\xi^{[n]}) & =
(i + pn) z^{i - 1} \xi^{[n]} + p \lambda z^{i - 1}\xi^{[n - 1]}
\end{align*}
受这些公式启发，按以下规则定义 $\theta$：
$$
\begin{matrix}
\theta(z^j)
& = & p\frac{z^{j + 1}}{j + 1}
& j = 0, \ldots p - 1, \\
\theta(z^{p - 1}\xi^{[m]})
& = & \xi^{[m + 1]}
& m \geq 1, \\
\theta(z^j \xi^{[m]})
& = &
\frac{p z^{j + 1} \xi^{[m]} - \theta(p\lambda z^j \xi^{[m - 1]})}{(j + 1 + pm)}
& 0 \leq j < p - 1, m \geq 1
\end{matrix}
$$
在最后一行中，对 $m$ 作归纳来定义所选的 $\theta$。展开可得（当
$0\leq j<p-1$ 且 $1\leq m$ 时）
$$
\theta(z^j \xi^{[m]}) =
\textstyle{\frac{p z^{j + 1} \xi^{[m]}}{(j + 1 + pm)} -
\frac{p^2 \lambda z^{j + 1} \xi^{[m - 1]}}{(j + 1 + pm)(j + 1 + p(m - 1))} +
\ldots +
\frac{(-1)^m p^{m + 1} \lambda^m z^{j + 1}}
{(j + 1 + pm) \ldots (j + 1)}}
$$
，不过下文不会使用此表达式。显然，$\theta$ 唯一延拓为 $D'$ 上
$p$-进连续的 $D$-线性映射。由构造，(\ref{crystalline-item-theta-linear}) 与
(\ref{crystalline-item-integrate}) 成立。尚须证明 (\ref{crystalline-item-factor}) 与
(\ref{crystalline-item-horizontal})。

\medskip\noindent
证明 (\ref{crystalline-item-factor}) 与 (\ref{crystalline-item-horizontal})。由于 $\theta$ 是
$D$-线性且连续的，只需证明
$p - \theta \circ \partial_z$ 与
$(\theta \otimes 1) \circ \text{d}_1 - \text{d}_1 \circ \theta$
分别作用于元素 $z^i\xi^{[n]}$ 时，所得分别属于 $D$ 与
$\Omega_D$\footnote{可直接计算：$p-\theta\circ\partial_z$ 作用于
$z^i\xi^{[n]}$ 时，除 $1$ 映到 $p$、$\xi$ 映到 $-p\lambda$ 外均为零。
而 $(\theta\otimes1)\circ\text d_1-\text d_1\circ\theta$ 作用于
$z^i\xi^{[n]}$ 时，除 $z^{p-1}\xi$ 映到 $-\lambda$ 外均为零。}。
置 $D_0=\mathbf Z_{(p)}[\lambda]$ 与
$D_0'=\mathbf Z_{(p)}[z,\lambda]\langle\xi\rangle/(\xi-z^p+\lambda)$。
注意上述每个表达式都是 $D_0'$ 或 $\Omega_{D_0'}$ 的元素。因此只需在
$D_0\to D_0'$ 的情形证明。注意 $D_0,D_0'$ 都是无挠环，且
$D_0\otimes\mathbf Q=\mathbf Q[\lambda]$、
$D'_0\otimes\mathbf Q=\mathbf Q[z,\lambda]$。故 $D_0\subset D'_0$ 是由
$\partial_z$ 消去的元素组成的子环，而 (\ref{crystalline-item-factor}) 由
(\ref{crystalline-item-integrate}) 推出；见紧邻 引理
\ref{crystalline-lemma-find-homotopy} 之前的讨论。类似地，
$\text{d}_1(f) = \partial_\lambda(f)\text{d}\lambda$，故
$$
\left((\theta \otimes 1) \circ \text{d}_1 - \text{d}_1 \circ \theta\right)(f)
=
\left(\theta(\partial_\lambda(f)) - \partial_\lambda(\theta(f))\right)
\text{d}\lambda
$$
对系数施加 $\partial_z$，得到
\begin{align*}
\partial_z\left(
\theta(\partial_\lambda(f)) - \partial_\lambda(\theta(f))
\right)
& =
p \partial_\lambda(f) - \partial_z(\partial_\lambda(\theta(f))) \\
& =
p \partial_\lambda(f) - \partial_\lambda(\partial_z(\theta(f))) \\
& =
p \partial_\lambda(f) - \partial_\lambda(p f) = 0
\end{align*}
故系数如所要求不依赖于 $z$。引理证毕。
\end{proof}

\noindent
注意，迭代 $\alpha_p$-覆盖 $X'\to X$（定义见本节引言）是有限局部自由的。
因此若 $X$ 连通，则 $X'\to X$ 的次数恒定且为 $p$ 的幂。

\begin{lemma}
\label{crystalline-lemma-pullback-along-p-power-cover}
设 $p$ 为素数，并设 $(S,\mathcal I,\gamma)$ 为 $\mathbf Z_{(p)}$ 上满足
$p\in\mathcal I$ 的除幂概形。置 $S_0=V(\mathcal I)\subset S$。设
$f:X'\to X$ 为 $S_0$ 上概形之间次数恒为 $q$ 的迭代
$\alpha_p$-覆盖。设 $\mathcal F$ 为 $X$ 上拟凝聚层中的任意晶体，并置
$\mathcal F'=f_{\text{cris}}^*\mathcal F$。在可区别三角
$$
Ru_{X/S, *}\mathcal{F}
\longrightarrow
f_*Ru_{X'/S, *}\mathcal{F}'
\longrightarrow
E
\longrightarrow
Ru_{X/S, *}\mathcal{F}[1]
$$
中，对象 $E$ 的各上同调层均被 $q$ 消去。
\end{lemma}

\begin{proof}
注意 $X'\to X$ 是同胚，故可等同 $X$ 与 $X'$ 的底拓扑空间。问题显然
在 $X$ 上是局部的，因此可假设 $X,X',S$ 均仿射，且 $X'\to X$ 写成复合
$$
X' = X_n \to X_{n - 1} \to X_{n - 2} \to \ldots \to X_0 = X
$$
，其中每个态射 $X_{i+1}\to X_i$ 都是 $\alpha_p$-覆盖。以
$\mathcal F_i$ 表示 $\mathcal F$ 到 $X_i$ 的拉回。只需证明每个映射
$$
R\Gamma(\text{Cris}(X_i/S), \mathcal{F}_i)
\longrightarrow
R\Gamma(\text{Cris}(X_{i + 1}/S), \mathcal{F}_{i + 1})
$$
都嵌入一个三角，其第三项的各上同调群被 $p$ 消去。（这里使用三角范畴
$D(X)$ 的公理 TR4；略去细节。）

\medskip\noindent
因此可假设 $S=\Spec(A)$、$X=\Spec(C)$、$X'=\Spec(C')$，且对某个
$c\in C$ 有 $C'=C[z]/(z^p-c)$。选取 $A$ 上多项式代数 $P$ 及满射
$P\to C$。令 $D$ 为 $P$ 中 $\Ker(P\to C)$ 的除幂包络之 $p$-进完备化，
如 (\ref{crystalline-equation-D})。置 $P'=P[z]$，并取把 $z$ 映到 $C'$ 中 $z$ 的类
的满射 $P'\to C'$。选取 $c\in C$ 的提升 $\lambda\in D$。于是，$P'$ 中
$\Ker(P'\to C')$ 的除幂包络之 $p$-进完备化 $D'$ 同构于下列环的
$p$-进完备化：
$D[z]\langle \xi \rangle/(\xi - (z^p - \lambda))$, see
，见引理 \ref{crystalline-lemma-computation} 及其证明。因此，利用命题
\ref{crystalline-proposition-compute-cohomology-crystal} 中拟凝聚模中晶体的上同调
计算，本引理立即推出结论。
\end{proof}

\noindent
下述引理中的界大概并非最优。

\begin{lemma}
\label{crystalline-lemma-pullback-along-p-power-cover-cohomology}
采用 引理 \ref{crystalline-lemma-pullback-along-p-power-cover} 的记号与假设，则映射
$$
f^* :
H^i(\text{Cris}(X/S), \mathcal{F})
\longrightarrow
H^i(\text{Cris}(X'/S), \mathcal{F}')
$$
的核与余核均被 $q^{i+1}$ 消去。
\end{lemma}

\begin{proof}
对象 $E$ 仅在次数 $-1$ 及以上可能有非零上同调层，故谱序列
$H^a(\mathcal H^b(E))\Rightarrow H^{a+b}(E)$ 收敛。把这一事实与可区别
三角所伴的上同调长正合列结合，即得所述界。
\end{proof}

\noindent
在 情形 \ref{crystalline-situation-global} 中，假设 $p\in\mathcal I$。置
$$
X^{(1)} = X \times_{S_0, F_{S_0}} S_0.
$$
以 $F_{X/S_0}:X\to X^{(1)}$ 表示相对 Frobenius 态射。

\begin{lemma}
\label{crystalline-lemma-pullback-relative-frobenius}
在上述情形中，假设 $X\to S_0$ 光滑且相对维数为 $d$。则 $F_{X/S_0}$
是次数 $p^d$ 的迭代 $\alpha_p$-覆盖。因此引理
\ref{crystalline-lemma-pullback-along-p-power-cover} 与
\ref{crystalline-lemma-pullback-along-p-power-cover-cohomology} 适用于此情形。
特别地，对 $\text{Cris}(X^{(1)}/S)$ 上拟凝聚模中的任意晶体
$\mathcal G$，映射
$$
F_{X/S_0}^* : H^i(\text{Cris}(X^{(1)}/S), \mathcal{G})
\longrightarrow
H^i(\text{Cris}(X/S), F_{X/S_0, \text{cris}}^*\mathcal{G})
$$
的核与余核均被 $p^{d(i+1)}$ 消去。
\end{lemma}

\begin{proof}
只需证明第一条断言。为此可假设 $X$ 在 $\mathbf A^d_{S_0}$ 上为
\'etale，见 “态射” 引理
\ref{morphisms-lemma-smooth-etale-over-affine-space}。以
$\varphi:X\to\mathbf A^d_{S_0}$ 表示这一 \'etale 态射。此时
$X/S_0$ 的相对 Frobenius 嵌入交换图
$$
\xymatrix{
X \ar[d] \ar[r] & X^{(1)} \ar[d] \\
\mathbf{A}^d_{S_0} \ar[r] & \mathbf{A}^d_{S_0}
}
$$
，其中下方横箭头是 $\mathbf A^d_{S_0}$ 在 $S_0$ 上的相对 Frobenius
态射。它把所有坐标提升到第 $p$ 次幂，故为迭代 $\alpha_p$-覆盖。
最后注意该图为纤维方块；见 \'Etale “态射” 引理
\ref{etale-lemma-relative-frobenius-etale}。
\end{proof}













\section{Frobenius 在晶体上同调上的作用}
\label{crystalline-section-frobenius}

\noindent
本节证明，在逆转素数 $p$ 后，Frobenius 拉回在晶体上同调上诱导拟同构。
但即使只是表述此结论，也需要处于一个特殊情形。

\begin{situation}
\label{crystalline-situation-F-crystal}
在 情形 \ref{crystalline-situation-global} 中，假设：
\begin{enumerate}
\item 对某个满足 $p\in I$ 的除幂环 $(A,I,\gamma)$，有
$S=\Spec(A)$；
\item 给定除幂环同态 $\sigma:A\to A$，使得对所有 $x\in A$，
$\sigma(x)=x^p\bmod pA$。
\end{enumerate}
\end{situation}

\noindent
在 情形 \ref{crystalline-situation-F-crystal} 中，态射
$\Spec(\sigma):S\to S$ 是绝对 Frobenius $F_{S_0}:S_0\to S_0$ 的提升，
并且图
$$
\xymatrix{
X \ar[d] \ar[r]_{F_X} & X \ar[d] \\
S_0 \ar[r]^{F_{S_0}} & S_0
}
$$
交换，其中 $F_X:X\to X$ 是 $X$ 的绝对 Frobenius 态射。因此得到晶体
拓扑斯的态射
$$
(F_X)_{\text{cris}} :
(X/S)_{\text{cris}}
\longrightarrow
(X/S)_{\text{cris}}
$$
，见注 \ref{crystalline-remark-functoriality-cris}。下面依照 Saavedra 的记号
给出关于 $F$-晶体的术语，见 \cite{Saavedra}。

\begin{definition}
\label{crystalline-definition-F-crystal}
在 情形 \ref{crystalline-situation-F-crystal} 中，所谓{\it $X/S$ 上（相对于
$\sigma$）的 $F$-晶体}，是指偶对 $(\mathcal E,F_\mathcal E)$，它由有限
局部自由 $\mathcal O_{X/S}$-模中的晶体 $\mathcal E$ 与映射
$$
F_\mathcal{E} : (F_X)_{\text{cris}}^*\mathcal{E} \longrightarrow \mathcal{E}
$$
组成。若存在整数 $i\geq0$ 及映射
$V:\mathcal E\to(F_X)_{\text{cris}}^*\mathcal E$，使
$V\circ F_\mathcal E=p^i\text{id}$，则称该 $F$-晶体{\it 非退化}。
\end{definition}

\begin{remark}
\label{crystalline-remark-F-crystal-variants}
设 $(\mathcal E,F)$ 为 定义 \ref{crystalline-definition-F-crystal} 中的
$F$-晶体。文献中常把非退化条件纳入 $F$-晶体的定义，并且还常假设
$F\circ V=p^n\text{id}$。下述结果所需的是：存在整数 $j\geq0$，使
$\Ker(F)$ 与 $\Coker(F)$ 均被 $p^j$ 消去。若 $\mathcal E$ 的秩有界
（例如当 $X$ 拟紧时），则这两个条件都由定义中所表述的非退化条件
推出。确切地说，设 $R$ 为环，$r\geq1$ 为整数，并设
$K,L\in\text{Mat}(r\times r,R)$ 满足 $KL=p^i1_{r\times r}$。则
$\det(K)\det(L)=p^{ri}$。令 $L'$ 为 $L$ 的伴随矩阵，即
$L'L=LL'=\det(L)$。置 $K'=p^{ri}K$ 与 $j=ri+i$。由于 $KL=p^i$，有
$K'L=p^j1_{r\times r}$，并且
$$
L K' = L K \det(L) \det(M) = L K L L' \det(M) = L p^i L' \det(M) =
p^j 1_{r \times r}
$$
因此，若 $V$ 如 定义 \ref{crystalline-definition-F-crystal}，则取
$V'=p^NV$，其中 $N>i\cdot\text{rank}(\mathcal E)$，便有
$V'\circ F=p^{N+i}$ 及 $F\circ V'=p^{N+i}$。
\end{remark}

\begin{theorem}
\label{crystalline-theorem-cohomology-F-crystal}
在 情形 \ref{crystalline-situation-F-crystal} 中，设
$(\mathcal E,F_\mathcal E)$ 为非退化 $F$-晶体。假设 $A$ 为 $p$-进完备
Noether 环，且 $X\to S_0$ 固有光滑。则典范映射
$$
F_\mathcal{E} \circ (F_X)_{\text{cris}}^* :
R\Gamma(\text{Cris}(X/S), \mathcal{E}) \otimes^\mathbf{L}_{A, \sigma} A
\longrightarrow
R\Gamma(\text{Cris}(X/S), \mathcal{E})
$$
在逆转 $p$ 后成为同构。
\end{theorem}

\begin{proof}
先把该箭头写成三个箭头的复合。置
$$
X^{(1)} = X \times_{S_0, F_{S_0}} S_0
$$
并以 $F_{X/S_0}:X\to X^{(1)}$ 表示相对 Frobenius 态射。以
$\mathcal E^{(1)}$ 表示 $\mathcal E$ 沿 $\Spec(\sigma)$ 的基变换；换言之，
它是经由与交换图
$$
\xymatrix{
X^{(1)} \ar[r] \ar[d] & X \ar[d] \\
S \ar[r]^{\Spec(\sigma)} & S
}
$$
相伴的晶体拓扑斯态射，把 $\mathcal E$ 拉回到
$\text{Cris}(X^{(1)}/S)$ 所得的晶体。于是有基变换映射
\begin{equation}
\label{crystalline-equation-base-change-sigma}
R\Gamma(\text{Cris}(X/S), \mathcal{E}) \otimes^\mathbf{L}_{A, \sigma} A
\longrightarrow
R\Gamma(\text{Cris}(X^{(1)}/S), \mathcal{E}^{(1)})
\end{equation}
，见注 \ref{crystalline-remark-base-change}。注意
$F_{X/S_0}:X\to X^{(1)}$ 与投影 $X^{(1)}\to X$ 的复合是绝对
Frobenius 态射 $F_X$。故
$F_{X/S_0}^*\mathcal E^{(1)}=(F_X)_{\text{cris}}^*\mathcal E$。因此沿
$F_{X/S_0}$ 的拉回给出映射
\begin{equation}
\label{crystalline-equation-to-prove}
F_{X/S_0}^* :
R\Gamma(\text{Cris}(X^{(1)}/S), \mathcal{E}^{(1)})
\longrightarrow
R\Gamma(\text{Cris}(X/S), (F_X)^*_{\text{cris}}\mathcal{E})
\end{equation}
最后可用 $F_\mathcal E$ 得到映射
\begin{equation}
\label{crystalline-equation-F-E}
R\Gamma(\text{Cris}(X/S), (F_X)^*_{\text{cris}}\mathcal{E})
\longrightarrow
R\Gamma(\text{Cris}(X/S), \mathcal{E})
\end{equation}
定理中的映射是上述三个映射 (\ref{crystalline-equation-base-change-sigma})、
(\ref{crystalline-equation-to-prove}) 与 (\ref{crystalline-equation-F-E}) 的复合。由注
\ref{crystalline-remark-base-change-isomorphism}，第一个映射模 $p$ 的任意幂都是
拟同构。由于所涉复形在 $D(A)$ 中完美，它本身即为拟同构，见注
\ref{crystalline-remark-complete-perfect}。第三个映射在逆转 $p$ 后为拟同构，因为
$F_\mathcal E$ 有一个相差 $p$ 的某次幂的逆，见注
\ref{crystalline-remark-F-crystal-variants}。最后，由引理
\ref{crystalline-lemma-pullback-relative-frobenius}，第二个映射在逆转 $p$ 后为同构。
\end{proof}
